\section{Blow-up schemes; projecting cones; projective closure}
\label{section:II.8}

\subsection{Blow-up preschemes}
\label{subsection:II.8.1}

\begin{env}[8.1.1]
\label{II.8.1.1}
Let $Y$ be a prescheme, and for every integer $n\geq0$, let $\mathcal{I}_n$ be a
quasi-coherent ideal of $\mathcal{O}_Y$; suppose that the following conditions
are satisfied:
\begin{equation}
\label{II.8.1.1.1}
  \mathcal{I}_0=\mathcal{O}_Y,\qquad
  \mathcal{I}_n\subset\mathcal{I}_m\quad\text{for }m\leq n
\tag{8.1.1.1}
\end{equation}
\begin{equation}
\label{II.8.1.1.2}
  \mathcal{I}_m\mathcal{I}_n\subset\mathcal{I}_{m+n}
  \quad\text{for all }m,n.
\tag{8.1.1.2}
\end{equation}

\oldpage[II]{153}
Note that these hypotheses imply
\begin{equation}
\label{II.8.1.1.3}
  \mathcal{I}_1^n\subset\mathcal{I}_n.
\tag{8.1.1.3}
\end{equation}

Set
\begin{equation}
\label{II.8.1.1.4}
  \mathcal{S}=\bigoplus_{n\geq0}\mathcal{I}_n.
\tag{8.1.1.4}
\end{equation}

It follows from (\hyperref[II.8.1.1.1]{8.1.1.1}) and
(\hyperref[II.8.1.1.2]{8.1.1.2}) that $\mathcal{S}$ is a quasi-coherent
graded $\mathcal{O}_Y$-algebra, and therefore defines a $Y$-scheme
$X=\operatorname{Proj}(\mathcal{S})$. If $\mathcal{J}$ is an \emph{invertible}
ideal of $\mathcal{O}_Y$,
$\mathcal{I}_n\otimes_{\mathcal{O}_Y}\mathcal{J}^{\otimes n}$ is canonically
identified with $\mathcal{I}_n\mathcal{J}^n$. Therefore, if we replace the
$\mathcal{I}_n$ by the $\mathcal{I}_n\mathcal{J}^n$, which replaces
$\mathcal{S}$ by a quasi-coherent $\mathcal{O}_Y$-algebra
$\mathcal{S}_{(\mathcal{J})}$, then
$X_{(\mathcal{J})}=\operatorname{Proj}(\mathcal{S}_{(\mathcal{J})})$ is
canonically isomorphic to $X$ (\hyperref[II.3.1.8]{3.1.8}).
\end{env}

\begin{env}[8.1.2]
\label{II.8.1.2}
Suppose that $Y$ is \emph{locally integral}, so that the sheaf
$\mathcal{R}(Y)$ of rational functions is a quasi-coherent $\mathcal{O}_Y$-algebra
(\hyperref[I.7.3.7]{I, 7.3.7}). We will say that a
sub-$\mathcal{O}_Y$-module $\mathcal{I}$ of $\mathcal{R}(Y)$ is a
\emph{fractional ideal} of $\mathcal{R}(Y)$ if it is of \emph{finite type}
(\hyperref[0.5.2.1]{0, 5.2.1}). Suppose given, for every $n\geq0$, a
quasi-coherent fractional ideal $\mathcal{I}_n$ of $\mathcal{R}(Y)$, such
that $\mathcal{I}_0=\mathcal{O}_Y$ and such that condition
(\hyperref[II.8.1.1.2]{8.1.1.2}) is satisfied (but not necessarily the
second condition in (\hyperref[II.8.1.1.1]{8.1.1.1})); we can then still
define, by the formula (\hyperref[II.8.1.1.4]{8.1.1.4}), a quasi-coherent
graded $\mathcal{O}_Y$-algebra, and the corresponding $Y$-scheme
$X=\operatorname{Proj}(\mathcal{S})$; we will again have a canonical
isomorphism of $X$ onto $X_{(\mathcal{J})}$ for every invertible fractional
ideal $\mathcal{J}$ of $\mathcal{R}(Y)$.
\end{env}

\begin{definition}[8.1.3]
\label{II.8.1.3}
Let $Y$ be a prescheme (resp. a locally integral prescheme), and let
$\mathcal{I}$ be a quasi-coherent ideal of $\mathcal{O}_Y$ (resp. a quasi-coherent
fractional ideal of $\mathcal{R}(Y)$). We say that the $Y$-scheme
$X=\operatorname{Proj}(\bigoplus_{n\geq0}\mathcal{I}^n)$ is obtained by
\emph{blowing up the ideal $\mathcal{I}$}, or is the \emph{blow-up prescheme
of $Y$ relative to $\mathcal{I}$}. When $\mathcal{I}$ is a
quasi-coherent ideal of $\mathcal{O}_Y$ and $Y'$ is the closed sub-prescheme of $Y$
defined by $\mathcal{I}$, we also say that $X$ is the $Y$-scheme obtained by
blowing up $Y'$.
\end{definition}

By definition, $\mathcal{S}=\bigoplus_{n\geq0}\mathcal{I}^n$ is then
generated by $\mathcal{S}_1=\mathcal{I}$; if $\mathcal{I}$ is an
$\mathcal{O}_Y$-module of \emph{finite type}, $X$ is therefore \emph{projective} over
$Y$ (\hyperref[II.5.5.2]{5.5.2}). Without any hypothesis on $\mathcal{I}$, the
$\mathcal{O}_X$-module $\mathcal{O}_X(1)$ is \emph{invertible}
(\hyperref[II.3.2.5]{3.2.5}) and, by virtue of
(\hyperref[II.4.4.3]{4.4.3}), \emph{very ample} for the structural morphism $X\to Y$.

Note that if $f:X\to Y$ is the structural morphism, the restriction of
$f$ to $f^{-1}(Y-Y')$ is an \emph{isomorphism} onto $Y-Y'$ when
$\mathcal{I}$ is an ideal of $\mathcal{O}_Y$ and $Y'$ is the closed sub-prescheme
that it defines: indeed, the question being local on $Y$, it suffices to
suppose $\mathcal{I}=\mathcal{O}_Y$, and our assertion follows from
(\hyperref[II.3.1.7]{3.1.7}).

When $\mathcal{I}$ is replaced by $\mathcal{I}^d$ ($d>0$), the blow-up
$Y$-scheme $X$ is replaced by a canonically isomorphic $Y$-scheme $X'$
(\hyperref[II.8.1.1]{8.1.1}); similarly, for every invertible ideal (resp.
fractional ideal) $\mathcal{J}$, the blow-up prescheme
$X_{(\mathcal{J})}$ relative to the ideal $\mathcal{I}\mathcal{J}$ is
canonically isomorphic to $X$ (\hyperref[II.8.1.1]{8.1.1}).

In particular, when $\mathcal{I}$ is an \emph{invertible} ideal (resp.
fractional ideal), the $Y$-scheme obtained by blowing up
$\mathcal{I}$ is \emph{isomorphic} to $Y$ (\hyperref[II.3.1.7]{3.1.7}).

\begin{proposition}[8.1.4]
\label{II.8.1.4}
Let $Y$ be an integral prescheme.
\begin{enumerate}
\item[\rm(i)] For every sequence $(\mathcal{I}_n)$ of quasi-coherent fractional
ideals of $\mathcal{R}(Y)$ satisfying
(\hyperref[II.8.1.1.2]{8.1.1.2})

\oldpage[II]{154}
and such that $\mathcal{I}_0=\mathcal{O}_Y$, the $Y$-scheme
$X=\operatorname{Proj}(\bigoplus_{n\geq0}\mathcal{I}_n)$ is integral and the
structural morphism $f:X\to Y$ is dominant.
\item[\rm(ii)] Let $\mathcal{I}$ be a quasi-coherent fractional ideal of
$\mathcal{R}(Y)$, and let $X$ be the blow-up $Y$-scheme of $Y$ relative to
$\mathcal{I}$. If $\mathcal{I}\neq0$, the structural morphism $f:X\to Y$ is
then birational and surjective.
\end{enumerate}
\end{proposition}

\begin{enumerate}
\item[\rm(i)] follows from the fact that $\mathcal{S}=\bigoplus_{n\geq0}\mathcal{I}_n$
is an \emph{integral} $\mathcal{O}_Y$-algebra
(\hyperref[II.3.1.12]{3.1.12} and \hyperref[II.3.1.14]{3.1.14}), since
for every $y\in Y$, $\mathcal{O}_y$ is an integral ring
(\hyperref[I.5.1.4]{I, 5.1.4}).
\item[\rm(ii)] By (i), $X$ is integral; if moreover $x$ and $y$ are the
generic points of $X$ and $Y$, we have $f(x)=y$, and we must prove that $k(x)$
is of rank $1$ over $k(y)$. Now $x$ is also the generic point of the fibre
$f^{-1}(y)$; if $\psi$ is the canonical morphism $Z\to Y$, where
$Z=\operatorname{Spec}(k(y))$, the prescheme $f^{-1}(y)$ is identified with
$\operatorname{Proj}(\mathcal{S}')$, where
$\mathcal{S}'=\psi^*(\mathcal{S})$ (\hyperref[II.3.5.3]{3.5.3}). But it is
clear that $\mathcal{S}'=\bigoplus_{n\geq0}(\mathcal{I}_y)^n$, and since
$\mathcal{I}$ is a nonzero quasi-coherent fractional ideal of
$\mathcal{R}(Y)$, $\mathcal{I}_y\neq0$
(\hyperref[I.7.3.6]{I, 7.3.6}), whence $\mathcal{I}_y=k(y)$;
$\operatorname{Proj}(\mathcal{S}')$ is therefore identified with
$\operatorname{Spec}(k(y))$ (\hyperref[II.3.1.7]{3.1.7}), whence the
conclusion.
\end{enumerate}

We will prove a \emph{converse} of
(\hyperref[II.8.1.4]{8.1.4}) in
(\hyperref[III.2.3.8]{III, 2.3.8}).

\begin{env}[8.1.5]
\label{II.8.1.5}
Let us return to the situation and notation of
(\hyperref[II.8.1.1]{8.1.1}). By definition, the injection homomorphisms
$\mathcal{I}_{n+1}\to\mathcal{I}_n$
(\hyperref[II.8.1.1.1]{8.1.1.1}) define, for each $k\in\mathbf{Z}$, an
injective homomorphism of degree $0$ of graded $\mathcal{S}$-modules
\begin{equation}
\label{II.8.1.5.1}
  u_k:\mathcal{S}_+(k+1)\to\mathcal{S}(k)~;
\tag{8.1.5.1}
\end{equation}
since $\mathcal{S}_+(k+1)$ and $\mathcal{S}(k+1)$ are canonically
(TN)-isomorphic, there corresponds canonically to $u_k$ an injective homomorphism
of $\mathcal{O}_X$-modules (\hyperref[II.3.4.2]{3.4.2}):
\begin{equation}
\label{II.8.1.5.2}
  \widetilde{u}_k:\mathcal{O}_X(k+1)\to\mathcal{O}_X(k).
\tag{8.1.5.2}
\end{equation}

Recall on the other hand (\hyperref[II.3.2.6]{3.2.6}) that we have defined
canonical homomorphisms
\begin{equation}
\label{II.8.1.5.3}
  \lambda:\mathcal{O}_X(h)\otimes_{\mathcal{O}_X}\mathcal{O}_X(k)
  \to\mathcal{O}_X(h+k)
\tag{8.1.5.3}
\end{equation}
and since the diagram
\[
\xymatrix{
  \mathcal{S}(h)\otimes_{\mathcal{S}}\mathcal{S}(k)
    \otimes_{\mathcal{S}}\mathcal{S}(l)\ar[r]\ar[d]
  & \mathcal{S}(h+k)\otimes_{\mathcal{S}}\mathcal{S}(l)\ar[d]\\
  \mathcal{S}(h)\otimes_{\mathcal{S}}\mathcal{S}(k+l)\ar[r]
  & \mathcal{S}(h+k+l)
}
\]
is commutative, it follows from the functoriality of the $\lambda$
(\hyperref[II.3.2.6]{3.2.6}) that the homomorphisms
(\hyperref[II.8.1.5.3]{8.1.5.3}) define on
\begin{equation}
\label{II.8.1.5.4}
  \mathcal{S}_X=\bigoplus_{n\in\mathbf{Z}}\mathcal{O}_X(n)
\tag{8.1.5.4}
\end{equation}

\oldpage[II]{155}
a structure of quasi-coherent graded $\mathcal{O}_X$-algebra. Moreover, the
diagram
\[
\xymatrix{
  \mathcal{S}(h)\otimes_{\mathcal{S}}\mathcal{S}_+(k+1)\ar[r]\ar[d]_{1\otimes u_k}
  & \mathcal{S}_+(h+k+1)\ar[d]^{u_{k+h}}\\
  \mathcal{S}(h)\otimes_{\mathcal{S}}\mathcal{S}(k)\ar[r]
  & \mathcal{S}(h+k)
}
\]
is commutative; the functoriality of the $\lambda$ then shows that we have a
commutative diagram
\begin{equation}
\label{II.8.1.5.5}
\xymatrix{
  \mathcal{O}_X(h)\otimes_{\mathcal{O}_X}\mathcal{O}_X(k+1)
    \ar[r]^{\lambda}\ar[d]_{1\otimes\widetilde{u}_k}
  & \mathcal{O}_X(h+k+1)\ar[d]^{\widetilde{u}_{k+h}}\\
  \mathcal{O}_X(h)\otimes_{\mathcal{O}_X}\mathcal{O}_X(k)
    \ar[r]^{\lambda}
  & \mathcal{O}_X(h+k)
}
\tag{8.1.5.5}
\end{equation}
the horizontal arrows being the canonical homomorphisms. We can therefore
say that the $\widetilde{u}_k$ define an \emph{injective homomorphism} (of
degree $0$) \emph{of graded $\mathcal{S}_X$-modules}
\begin{equation}
\label{II.8.1.5.6}
  \widetilde{u}:\mathcal{S}_X(1)\to\mathcal{S}_X.
\tag{8.1.5.6}
\end{equation}
\end{env}

\begin{env}[8.1.6]
\label{II.8.1.6}
Keeping the notation of (\hyperref[II.8.1.5]{8.1.5}), let us now remark
that, for $n\geq0$, the composite homomorphism
$\widetilde{v}_n=\widetilde{u}_{n-1}\widetilde{u}_{n-2}\cdots\widetilde{u}_0$
is an \emph{injective} homomorphism $\mathcal{O}_X(n)\to\mathcal{O}_X$;
we denote by $\mathcal{I}_{n,X}$ its image, which is therefore a
quasi-coherent ideal of $\mathcal{O}_X$, \emph{isomorphic} to $\mathcal{O}_X(n)$.
Moreover, the diagram
\[
\xymatrix{
  \mathcal{O}_X(m)\otimes_{\mathcal{O}_X}\mathcal{O}_X(n)
    \ar[r]^{\lambda}\ar[d]_{\widetilde{v}_m\otimes\widetilde{v}_n}
  & \mathcal{O}_X(m+n)\ar[d]^{\widetilde{v}_{m+n}}\\
  \mathcal{O}_X\ar[r]_{\mathrm{id}}
  & \mathcal{O}_X
}
\]
is commutative for $m\geq0$, $n\geq0$. From this we conclude the following
inclusions
\begin{equation}
\label{II.8.1.6.1}
  \mathcal{I}_{0,X}=\mathcal{O}_X,\qquad
  \mathcal{I}_{n,X}\subset\mathcal{I}_{m,X}
  \quad\text{for }0\leq m\leq n
\tag{8.1.6.1}
\end{equation}
\begin{equation}
\label{II.8.1.6.2}
  \mathcal{I}_{m,X}\mathcal{I}_{n,X}\subset\mathcal{I}_{m+n,X}
  \quad\text{for }m\geq0,\ n\geq0.
\tag{8.1.6.2}
\end{equation}
\end{env}

\oldpage[II]{156}

\begin{proposition}[8.1.7]
\label{II.8.1.7}
Let $Y$ be a prescheme, $\mathcal{I}$ a quasi-coherent ideal of
$\mathcal{O}_Y$, and $X=\operatorname{Proj}(\bigoplus_{n\geq0}\mathcal{I}^n)$ the
$Y$-scheme obtained by blowing up $\mathcal{I}$. Then, for every
$n>0$, we have a canonical isomorphism
\begin{equation}
\label{II.8.1.7.1}
  \mathcal{O}_X(n)\xrightarrow{\sim}\mathcal{I}^n\mathcal{O}_X
  =\mathcal{I}_{n,X}
\tag{8.1.7.1}
\end{equation}
(cf. (\hyperref[0.4.3.5]{0, 4.3.5})), and consequently
$\mathcal{I}^n\mathcal{O}_X$ is a very ample invertible $\mathcal{O}_X$-module
if $n>0$.
\end{proposition}

The last assertion is immediate, since $\mathcal{O}_X(1)$ is invertible
(\hyperref[II.3.2.5]{3.2.5}) and very ample for $Y$ by definition
(\hyperref[II.4.4.3]{4.4.3} and \hyperref[II.4.4.9]{4.4.9}). On the
other hand, by definition, the image of $v_n$ is none other than
$\mathcal{I}^n\mathcal{S}$, and (\hyperref[II.8.1.7.1]{8.1.7.1}) therefore follows
from the exactness of the functor $\widetilde{\mathcal{M}}$
(\hyperref[II.3.2.4]{3.2.4}) and from the formula
(\hyperref[II.3.2.4.1]{3.2.4.1}).

\begin{corollary}[8.1.8]
\label{II.8.1.8}
Under the hypotheses of (\hyperref[II.8.1.7]{8.1.7}), if $f:X\to Y$ is the
structural morphism and $Y'$ the closed sub-prescheme of $Y$ defined by
$\mathcal{I}$, the closed sub-prescheme $X'=f^{-1}(Y')$ of $X$ is defined by
$\mathcal{I}\mathcal{O}_X$ (canonically isomorphic to $\mathcal{O}_X(1)$),
whence a canonical exact sequence
\begin{equation}
\label{II.8.1.8.1}
  0\to\mathcal{O}_X(1)\to\mathcal{O}_X\to\mathcal{O}_{X'}\to0.
\tag{8.1.8.1}
\end{equation}
\end{corollary}

This follows from (\hyperref[II.8.1.7.1]{8.1.7.1}) and
(\hyperref[I.4.4.5]{I, 4.4.5}).

\begin{env}[8.1.9]
\label{II.8.1.9}
Under the hypotheses of (\hyperref[II.8.1.7]{8.1.7}), we can make the
structure of the $\mathcal{I}_{n,X}$ more precise. Indeed, let us remark that the homomorphism
\[
  \widetilde{u}_{-1}:\mathcal{O}_X\to\mathcal{O}_X(-1)
\]
corresponds canonically to a section $s$ of $\mathcal{O}_X(-1)$ over
$X$, which we will call the \emph{canonical section} (relative to
$\mathcal{I}$) (\hyperref[0.5.1.1]{0, 5.1.1}). On the other hand, in the
diagram (\hyperref[II.8.1.5.5]{8.1.5.5}), the horizontal arrows are
isomorphisms (\hyperref[II.3.2.7]{3.2.7}); replacing $h$ by $k$ and $k$ by
$-1$ in this diagram, we obtain
$\widetilde{u}_k=1_k\otimes\widetilde{u}_{-1}$ ($1_h$ denoting the identity of
$\mathcal{O}_X(h)$), in other words the homomorphism $\widetilde{u}_k$ is
none other than \emph{tensor multiplication by the canonical section $s$}
(for every $k\in\mathbf{Z}$). The homomorphism $\widetilde{u}$
(\hyperref[II.8.1.5.6]{8.1.5.6}) can therefore also be interpreted in the same way.

Consequently, for every $n\geq0$, the homomorphism
$\widetilde{v}_n:\mathcal{O}_X(n)\to\mathcal{O}_X$ is none other than
tensor multiplication by $s^{\otimes n}$; we deduce from this:
\end{env}

\begin{corollary}[8.1.10]
\label{II.8.1.10}
With the notation of (\hyperref[II.8.1.8]{8.1.8}), the underlying space of
$X'$ is the set of $x\in X$ such that $s(x)=0$, $s$ denoting the canonical
section of $\mathcal{O}_X(-1)$.
\end{corollary}

Indeed, if $c_x$ is a generator of the stalk $(\mathcal{O}_X(1))_x$ at a
point $x$, $s_x\otimes c_x$ is canonically identified with a generator of the
stalk of $\mathcal{I}_{1,X}$ at the point $x$, and is therefore invertible if and
only if $s_x\notin\mathfrak{m}_x(\mathcal{O}_X(-1))_x$, in other words if
$s(x)\neq0$.

\begin{proposition}[8.1.11]
\label{II.8.1.11}
Let $Y$ be an integral prescheme, $\mathcal{I}$ a quasi-coherent fractional
ideal of $\mathcal{R}(Y)$, and $X$ the $Y$-scheme obtained by blowing up
$\mathcal{I}$. Then $\mathcal{I}\mathcal{O}_X$ is a very ample invertible
$\mathcal{O}_X$-module for $Y$.
\end{proposition}

The question being local on $Y$ (\hyperref[II.4.4.5]{4.4.5}), we may
restrict to the case where $Y=\operatorname{Spec}(A)$, where $A$ is an integral
ring with field of fractions $K$ and $\mathcal{I}=\widetilde{\mathfrak{I}}$,
$\mathfrak{I}$ being a fractional ideal of $K$; consequently there exists an
element $a\neq0$ of $A$ such that $a\mathfrak{I}\subset A$. Set
$S=\bigoplus_{n\geq0}\mathfrak{I}^n$; the map $x\mapsto ax$ is an
$A$-isomorphism of $\mathfrak{I}^{n+1}=(S(1))_n$ onto
$a\mathfrak{I}^{n+1}=a\mathfrak{I}S_n\subset\mathfrak{I}^n=S_n$,

\oldpage[II]{157}
and therefore defines a (TN)-isomorphism of degree $0$ of graded $S$-modules
$S_+(1)\to a\mathfrak{I}S$. Furthermore, $x\mapsto a^{-1}x$ is an
isomorphism of degree $0$ of graded $S$-modules
$a\mathfrak{I}S\xrightarrow{\sim}\mathfrak{I}S$. By composition
(\hyperref[II.3.2.4]{3.2.4}) we therefore obtain an isomorphism of
$\mathcal{O}_X$-modules
$\mathcal{O}_X(1)\xrightarrow{\sim}\mathcal{I}\mathcal{O}_X$, and since $S$
is generated by $S_1=\mathfrak{I}$, $\mathcal{O}_X(1)$ is invertible
(\hyperref[II.3.2.5]{3.2.5}) and very ample
(\hyperref[II.4.4.3]{4.4.3} and \hyperref[II.4.4.9]{4.4.9}), whence our
assertion.

\subsection{Preliminary results on localization in graded rings}
\label{subsection:II.8.2}

\begin{env}[8.2.1]
\label{II.8.2.1}
Let $S$ be a graded ring, where for the moment we do not assume the degrees
to be positive. We set
\begin{equation}
\label{II.8.2.1.1}
  S^{\geq}=\bigoplus_{n\geq0}S_n,
  \qquad
  S^{\leq}=\bigoplus_{n\leq0}S_n,
\tag{8.2.1.1}
\end{equation}
which are graded subrings of $S$, with degrees respectively all
positive and all negative. If $f$ is a homogeneous element of degree $d$
(positive or negative) of $S$, the ring of fractions $S_f=S'$ is again endowed
with a graded ring structure, by taking for $S'_n$ ($n\in\mathbf{Z}$)
the set of $x/f^k$, where $x\in S_{n+kd}$ ($k\geq0$); we again set
$S_{(f)}=S'_0$, and we write $S_f^{\geq}$ and $S_f^{\leq}$ for $S'^{\geq}$ and
$S'^{\leq}$ respectively. If $d>0$, we have
\begin{equation}
\label{II.8.2.1.2}
  (S^{\geq})_f=S_f
\tag{8.2.1.2}
\end{equation}
since, if $x\in S_{n+kd}$, with $n+kd<0$, we can write
$x/f^k=xf^h/f^{h+k}$ and $n+(h+k)d>0$ for $h$ sufficiently large and $>0$. We
conclude by definition that
\begin{equation}
\label{II.8.2.1.3}
  (S^{\geq})_{(f)}=(S_f^{\geq})_0=S_{(f)}.
\tag{8.2.1.3}
\end{equation}

If $M$ is a graded $S$-module, we similarly set
\begin{equation}
\label{II.8.2.1.4}
  M^{\geq}=\bigoplus_{n\geq0}M_n,
  \qquad
  M^{\leq}=\bigoplus_{n\leq0}M_n
\tag{8.2.1.4}
\end{equation}
which are respectively a graded $S^{\geq}$-module and a graded
$S^{\leq}$-module, and whose intersection is the $S_0$-module $M_0$. If $f\in S_d$, we
again define $M_f$ as a graded $S_f$-module by taking as elements of
degree $n$ the $z/f^k$, where $z\in M_{n+kd}$ ($k\geq0$); we denote by
$M_{(f)}$ the set of elements of degree $0$ of $M_f$, which is an
$S_{(f)}$-module, and we write $M_f^{\geq}$ and $M_f^{\leq}$ instead of
$(M_f)^{\geq}$ and $(M_f)^{\leq}$ respectively. If $d>0$, we see as
above that
\begin{equation}
\label{II.8.2.1.5}
  (M^{\geq})_f=M_f
\tag{8.2.1.5}
\end{equation}
and
\begin{equation}
\label{II.8.2.1.6}
  (M^{\geq})_{(f)}=(M_f^{\geq})_0=M_{(f)}.
\tag{8.2.1.6}
\end{equation}
\end{env}

\begin{env}[8.2.2]
\label{II.8.2.2}
Let $z$ be an indeterminate, which we will call the \emph{homogenization
variable}. If $S$ is a graded ring (graded in positive or
negative degrees), the polynomial algebra\footnote{There can be no
confusion here with the use of the notation $\widehat{S}$ to denote the separated
completion of a ring.}
\begin{equation}
\label{II.8.2.2.1}
  \widehat{S}=S[z]
\tag{8.2.2.1}
\end{equation}

\oldpage[II]{158}
is a graded $S$-algebra, when we take as the degree of $fz^n$ ($n\geq0$), where
$f$ is homogeneous,
\begin{equation}
\label{II.8.2.2.2}
  \deg(fz^n)=n+\deg f.
\tag{8.2.2.2}
\end{equation}
\end{env}

\begin{lemma}[8.2.3]
\label{II.8.2.3}
\begin{enumerate}
\item[\rm(i)] We have canonical isomorphisms of (ungraded) rings
\begin{equation}
\label{II.8.2.3.1}
  \widehat{S}_{(z)}\xrightarrow{\sim}
  \widehat{S}/(z-1)\widehat{S}\xrightarrow{\sim}S.
\tag{8.2.3.1}
\end{equation}
\item[\rm(ii)] We have a canonical isomorphism of (ungraded) rings
\begin{equation}
\label{II.8.2.3.2}
  \widehat{S}_{(f)}\xrightarrow{\sim}S_f^{\leq}
\tag{8.2.3.2}
\end{equation}
for every $f\in S_d$, with $d>0$.
\end{enumerate}
\end{lemma}

The first of the isomorphisms of (\hyperref[II.8.2.3.1]{8.2.3.1}) was
defined in (\hyperref[II.2.2.5]{2.2.5}) and the second is trivial;
the isomorphism $\widehat{S}_{(z)}\xrightarrow{\sim}S$ thus defined therefore
assigns to $xz^n/z^{n+k}$, where $\deg(x)=k$ ($k\geq-n$), the element $x$.
The homomorphism (\hyperref[II.8.2.3.2]{8.2.3.2}) assigns to
$xz^n/f^k$, where $\deg(x)=kd-n$, the element $x/f^k$, of degree $-n$ in
$S_f^{\leq}$, and it is again evident that we indeed have an isomorphism here.

\begin{env}[8.2.4]
\label{II.8.2.4}
Let $M$ be a graded $S$-module. It is clear that the $S$-module
\begin{equation}
\label{II.8.2.4.1}
  \widehat{M}=M\otimes_S\widehat{S}=M\otimes_S S[z]
\tag{8.2.4.1}
\end{equation}
is the direct sum of the $S$-modules $M\otimes_S Sz^n$, hence of the abelian groups
$M_k\otimes_S Sz^n$ ($k\in\mathbf{Z}$, $n\geq0$); we define on
$\widehat{M}$ a structure of graded $\widehat{S}$-module by taking
\begin{equation}
\label{II.8.2.4.2}
  \deg(x\otimes z^n)=n+\deg x
\tag{8.2.4.2}
\end{equation}
for every homogeneous $x$ in $M$. We leave it to the reader to prove
the analogue of (\hyperref[II.8.2.3]{8.2.3}):
\end{env}

\begin{lemma}[8.2.5]
\label{II.8.2.5}
\begin{enumerate}
\item[\rm(i)] We have a canonical di-isomorphism of (ungraded) modules
\begin{equation}
\label{II.8.2.5.1}
  \widehat{M}_{(z)}\xrightarrow{\sim}M.
\tag{8.2.5.1}
\end{equation}
\item[\rm(ii)] For every $f\in S_d$ ($d>0$), we have a di-isomorphism of
(ungraded) modules
\begin{equation}
\label{II.8.2.5.2}
  \widehat{M}_{(f)}\xrightarrow{\sim}M_f^{\leq}.
\tag{8.2.5.2}
\end{equation}
\end{enumerate}
\end{lemma}

\begin{env}[8.2.6]
\label{II.8.2.6}
Let $S$ be a ring graded in \emph{positive} degrees, and consider in $S$
the decreasing sequence of graded ideals
\begin{equation}
\label{II.8.2.6.1}
  S_{[n]}=\bigoplus_{m\geq n}S_m\qquad(n\geq0)
\tag{8.2.6.1}
\end{equation}
(in particular we have $S_{[0]}=S$, $S_{[1]}=S_+$). Since it is clear that
$S_{[m]}S_{[n]}\subset S_{[m+n]}$, we can define a \emph{graded ring}
$S^\natural$ by taking
\begin{equation}
\label{II.8.2.6.2}
  S^\natural=\bigoplus_{n\geq0}S_n^\natural
  \quad\text{with}\quad S_n^\natural=S_{[n]}.
\tag{8.2.6.2}
\end{equation}
$S_0^\natural$ is therefore the ring $S$ considered as an \emph{ungraded} ring,
and $S^\natural$ is consequently an $S_0^\natural$-algebra. For every
homogeneous element $f\in S_d$ ($d>0$), we denote by $f^\natural$ the element
$f$ considered as belonging to $S_{[d]}=S_d^\natural$. With this
notation:
\end{env}

\oldpage[II]{159}

\begin{lemma}[8.2.7]
\label{II.8.2.7}
Let $S$ be a ring graded in positive degrees, and $f$ a homogeneous element of
$S_d$ ($d>0$). We have canonical isomorphisms of rings
\begin{equation}
\label{II.8.2.7.1}
  S_f\xrightarrow{\sim}\bigoplus_{n\in\mathbf{Z}}S(n)_{(f)}
\tag{8.2.7.1}
\end{equation}
\begin{equation}
\label{II.8.2.7.2}
  (S_f^{\geq})_{f/1}\xrightarrow{\sim}S_f
\tag{8.2.7.2}
\end{equation}
\begin{equation}
\label{II.8.2.7.3}
  S^\natural_{(f^\natural)}\xrightarrow{\sim}S_f^{\geq}
\tag{8.2.7.3}
\end{equation}
of which the first two are isomorphisms of graded rings.
\end{lemma}

It is immediate by definition that
$(S_f)_n=(S(n)_f)_0$, whence the isomorphism
(\hyperref[II.8.2.7.1]{8.2.7.1}), which is none other than the identity. On the other
hand, since $f/1$ is invertible in $S_f$, there is a canonical isomorphism
$S_f\xrightarrow{\sim}(S_f^{\geq})_{f/1}=(S_f)_{f/1}$ by virtue of
(\hyperref[II.8.2.1.2]{8.2.1.2}) applied to $S_f$; the inverse
isomorphism is by definition the isomorphism
(\hyperref[II.8.2.7.2]{8.2.7.2}). Finally, if
$x=\sum_{m\geq n}y_m$ is an element of $S_{[n]}$, with $n=kd$, we assign
to the element $x/(f^\natural)^k$ the element
$\sum_m y_m/f^k$ of $S_f^{\geq}$, and we verify immediately that this defines an
isomorphism (\hyperref[II.8.2.7.3]{8.2.7.3}).

\begin{env}[8.2.8]
\label{II.8.2.8}
If $M$ is a graded $S$-module, we similarly set, for every $n\in\mathbf{Z}$,
\begin{equation}
\label{II.8.2.8.1}
  M_{[n]}=\bigoplus_{m\geq n}M_m
\tag{8.2.8.1}
\end{equation}
and since $S_{[m]}M_{[n]}\subset M_{[m+n]}$ ($m\geq0$), we can define a
graded $S^\natural$-module $M^\natural$ by taking
\begin{equation}
\label{II.8.2.8.2}
  M^\natural=\bigoplus_{n\in\mathbf{Z}}M_n^\natural,
  \quad\text{with}\quad M_n^\natural=M_{[n]}.
\tag{8.2.8.2}
\end{equation}
\end{env}

We again leave to the reader the proof of the following:

\begin{lemma}[8.2.9]
\label{II.8.2.9}
With the notation of (\hyperref[II.8.2.7]{8.2.7}) and
(\hyperref[II.8.2.8]{8.2.8}), we have the canonical di-isomorphisms of
modules
\begin{equation}
\label{II.8.2.9.1}
  M_f\xrightarrow{\sim}\bigoplus_{n\in\mathbf{Z}}M(n)_{(f)}
\tag{8.2.9.1}
\end{equation}
\begin{equation}
\label{II.8.2.9.2}
  (M_f^{\geq})_{f/1}\xrightarrow{\sim}M_f
\tag{8.2.9.2}
\end{equation}
\begin{equation}
\label{II.8.2.9.3}
  M^\natural_{(f^\natural)}\xrightarrow{\sim}M_f^{\geq}
\tag{8.2.9.3}
\end{equation}
of which the first two are di-isomorphisms of graded modules.
\end{lemma}

\begin{lemma}[8.2.10]
\label{II.8.2.10}
Let $S$ be a ring graded in positive degrees.
\begin{enumerate}
\item[\rm(i)] For $S^\natural$ to be an $S_0^\natural$-algebra of finite
type (resp. Noetherian), it is necessary and sufficient that $S$ be an $S_0$-algebra
of finite type (resp. Noetherian).
\item[\rm(ii)] For $S_{n+1}^\natural=S_1^\natural S_n^\natural$ to hold for
$n\geq n_0$, it is necessary and sufficient that $S_{n+1}=S_1S_n$ for $n\geq n_0$.
\item[\rm(iii)] For $S_n^\natural=S_1^{\natural n}$ to hold for $n\geq n_0$, it
is necessary and sufficient that $S_n=S_1^n$ for $n\geq n_0$.
\item[\rm(iv)] If $(f_\alpha)$ is a set of homogeneous elements of $S_+$
such that $S_+$ is contained in the radical in $S$ of the ideal of $S$ generated by
the $f_\alpha$, then $S_+^\natural$ is contained in the radical in $S^\natural$ of
the ideal of $S^\natural$ generated by the $f_\alpha^\natural$.
\end{enumerate}
\end{lemma}

\begin{enumerate}
\item[\rm(i)] If $S^\natural$ is an $S_0^\natural$-algebra of finite type,
$S_+=S_1^\natural$ is a module of finite type over $S=S_0^\natural$ by
(\hyperref[II.2.1.6]{2.1.6}, (i)), hence $S$ is an $S_0$-algebra of finite
type (\hyperref[II.2.1.4]{2.1.4}); if $S^\natural$ is a Noetherian
ring, so is $S_0^\natural=S$
(\hyperref[II.2.1.5]{2.1.5}). Conversely, if $S$ is an $S_0$-algebra

\oldpage[II]{160}
of finite type, we know (\hyperref[II.2.1.6]{2.1.6}, (ii)) that there exist
$h>0$ and $m_0>0$ such that $S_{n+h}=S_hS_n$ for $n\geq m_0$; we may
evidently suppose $m_0\geq h$. Moreover, the $S_m$ are $S_0$-modules
of finite type (\hyperref[II.2.1.6]{2.1.6}, (i)). This being so, if
$n\geq m_0+h$, we have $S_n^\natural=S_hS_{n-h}^\natural
=S_h^\natural S_{n-h}^\natural$; and if $m<m_0+h$, we have, setting
$E=S_{m_0}+\cdots+S_{m_0+h-1}$,
\[
  S_m^\natural=S_m+\cdots+S_{m_0+h-1}+S_hE+S_h^2E+\cdots.
\]
For $1\leq m\leq m_0$, let $G_m$ be the union of finite systems of
generators of the $S_0$-modules $S_i$, for $m\leq i\leq m_0+h-1$,
considered as a subset of $S_{[m]}$. For $m_0+1\leq m\leq m_0+h-1$, let
similarly $G_m$ be the union of finite systems of generators of the $S_0$-modules
$S_i$, for $m\leq i\leq m_0+h-1$, and of $S_hE$, considered as a subset of
$S_{[m]}$. It is clear that $S_m^\natural=S_0^\natural G_m$ for
$1\leq m\leq m_0+h-1$, and consequently the union $G$ of the $G_m$ for
$1\leq m\leq m_0+h-1$ is a system of generators of the
$S_0^\natural$-algebra $S^\natural$. We conclude that if $S=S_0^\natural$
is a Noetherian ring, so is $S^\natural$.

\item[\rm(ii)] It is clear that if $S_{n+1}=S_1S_n$ for $n\geq n_0$, we have
$S_{n+1}^\natural=S_1S_n^\natural$, and \emph{a fortiori}
$S_{n+1}^\natural=S_1^\natural S_n^\natural$ for $n\geq n_0$.
Conversely, this last relation can be written
\[
  S_{n+1}+S_{n+2}+\cdots=(S_1+S_2+\cdots)(S_n+S_{n+1}+\cdots)
\]
and comparing the terms of degree $n+1$ (in $S$) on the two sides gives
$S_{n+1}=S_1S_n$.

\item[\rm(iii)] If $S_n=S_1^n$ for $n\geq n_0$, we have
$S_n^\natural=S_1^n+S_1^{n+1}+\cdots$; since $S_1^\natural$ contains
$S_1+S_1^2+\cdots$, we have $S_n^\natural\subset S_1^{\natural n}$, and consequently
$S_n^\natural=S_1^{\natural n}$ for $n\geq n_0$. Conversely, the only
terms of $S_1^{\natural n}=(S_1+S_2+\cdots)^n$ that are of degree $n$ in
$S$ are those of $S_1^n$; the relation $S_n^\natural=S_1^{\natural n}$
therefore implies $S_n=S_1^n$.

\item[\rm(iv)] It suffices to prove that if an element $g\in S_{k+h}$ is
considered as an element of $S_k^\natural$ ($k>0$, $h\geq0$), then there
exists an integer $n>0$ such that, in $S_{kn}^\natural$, $g^n$ is a linear
combination of the $f_\alpha^\natural$ with coefficients in $S^\natural$. By
hypothesis, there is an integer $m_0$ such that for $m\geq m_0$ we have, \emph{in
$S$}, $g^m=\sum_\alpha c_{\alpha m}f_\alpha$, the indices $\alpha$ appearing
in this formula being \emph{independent of $m$}; moreover, we may
evidently suppose the $c_{\alpha m}$ homogeneous, with
\[
  \deg(c_{\alpha m})=m(k+h)-\deg f_\alpha
\]
in $S$. Then take $m_0$ sufficiently large that
$km_0>\deg f_\alpha$ for all the $f_\alpha$ that appear in $g^{m_0}$;
for every $\alpha$, let $c'_{\alpha m}$ be the element $c_{\alpha m}$ considered
as being of degree $km-\deg(f_\alpha)$ in $S^\natural$; we then have, in
$S^\natural$, $g^m=\sum_\alpha c'_{\alpha m}f_\alpha^\natural$, which
completes the proof.
\end{enumerate}

\begin{env}[8.2.11]
\label{II.8.2.11}
Consider the graded $S_0$-algebra
\begin{equation}
\label{II.8.2.11.1}
  S^\natural\otimes_S S_0=S^\natural/S_+S^\natural
  =\bigoplus_{n\geq0}S_{[n]}/S_+S_{[n]}.
\tag{8.2.11.1}
\end{equation}
Since $S_n$ is a quotient $S_0$-module of $S_{[n]}/S_+S_{[n]}$, we have a
canonical homomorphism of graded $S_0$-algebras
\begin{equation}
\label{II.8.2.11.2}
  S^\natural\otimes_S S_0\to S
\tag{8.2.11.2}
\end{equation}
which is evidently \emph{surjective}, and consequently corresponds
(\hyperref[II.2.9.2]{2.9.2}) to a canonical \emph{closed immersion}
\begin{equation}
\label{II.8.2.11.3}
  \operatorname{Proj}(S)\to
  \operatorname{Proj}(S^\natural\otimes_S S_0).
\tag{8.2.11.3}
\end{equation}
\end{env}

\oldpage[II]{161}

\begin{proposition}[8.2.12]
\label{II.8.2.12}
The canonical morphism (\hyperref[II.8.2.11.3]{8.2.11.3}) is bijective.
For the homomorphism (\hyperref[II.8.2.11.2]{8.2.11.2}) to be
(TN)-bijective, it is necessary and sufficient that there exist $n_0$ such that
$S_{n+1}=S_1S_n$ for $n\geq n_0$. If this last condition is satisfied,
(\hyperref[II.8.2.11.3]{8.2.11.3}) is an isomorphism; the converse is
true when $S$ is Noetherian.
\end{proposition}

To prove the first assertion, it suffices
(\hyperref[II.2.8.3]{2.8.3}) to prove that the kernel $\mathfrak{I}$ of
the homomorphism (\hyperref[II.8.2.11.2]{8.2.11.2}) consists of
\emph{nilpotent} elements. Now if $f\in S_{[n]}$ is an element whose class modulo
$S_+S_{[n]}$ belongs to this kernel, this means that $f\in S_{[n+1]}$;
the element $f^{n+1}$, considered as belonging to $S_{[n(n+1)]}$, then
belongs to $S_+S_{[n(n+1)]}$, since it can be written $f\cdot f^n$; hence the class of
$f^{n+1}$ modulo $S_+S_{[n(n+1)]}$ is zero, which proves our assertion.
Since the hypothesis $S_{n+1}=S_1S_n$ for $n\geq n_0$ is equivalent to
$S_{n+1}^{\natural}=S_1^{\natural}S_n^{\natural}$ for $n\geq n_0$
(\hyperref[II.8.2.10]{8.2.10}, (ii)), this hypothesis is equivalent by definition
to the fact that (\hyperref[II.8.2.11.2]{8.2.11.2}) is (TN)-injective, hence
(TN)-bijective, and then (\hyperref[II.8.2.11.3]{8.2.11.3}) is an
isomorphism by virtue of (\hyperref[II.2.9.1]{2.9.1}). Conversely, if
(\hyperref[II.8.2.11.3]{8.2.11.3}) is an isomorphism, the sheaf
$\widetilde{\mathfrak{I}}$ on
$\operatorname{Proj}(S^{\natural}\otimes_S S_0)$ is zero
(\hyperref[II.2.9.2]{2.9.2}, (i)); since
$S^{\natural}\otimes_S S_0$ is Noetherian, being a quotient of $S^{\natural}$
(\hyperref[II.8.2.10]{8.2.10}, (i)), we conclude from
(\hyperref[II.2.7.3]{2.7.3}) that $\mathfrak{I}$ satisfies condition (TN),
hence $S_{n+1}^{\natural}=S_1^{\natural}S_n^{\natural}$ for $n\geq n_0$, and
this completes the proof by virtue of
(\hyperref[II.8.2.10]{8.2.10}, (ii)).

\begin{env}[8.2.13]
\label{II.8.2.13}
Now consider the canonical injections $(S_+)^n\to S_{[n]}$, which
define an injective homomorphism of degree $0$ of graded rings
\begin{equation}
\label{II.8.2.13.1}
  \bigoplus_{n\geq0}(S_+)^n\to S^{\natural}.
\tag{8.2.13.1}
\end{equation}
\end{env}

\begin{proposition}[8.2.14]
\label{II.8.2.14}
For the homomorphism (\hyperref[II.8.2.13.1]{8.2.13.1}) to be a
(TN)-isomorphism, it is necessary and sufficient that there exist $n_0$ such that
$S_n=S_1^n$ for every $n\geq n_0$. When this is so, the morphism
corresponding to (\hyperref[II.8.2.13.1]{8.2.13.1}) is everywhere defined and
is an isomorphism
\[
  \operatorname{Proj}(S^{\natural})\xrightarrow{\sim}
  \operatorname{Proj}\left(\bigoplus_{n\geq0}(S_+)^n\right)~;
\]
the converse is true when $S$ is Noetherian.
\end{proposition}

The first two assertions are evident, in view of
(\hyperref[II.8.2.10]{8.2.10}, (iii)) and
(\hyperref[II.2.9.1]{2.9.1}). The third will follow from
(\hyperref[II.8.2.10]{8.2.10}, (i) and (iii)) and from the following lemma:

\begin{lemma}[8.2.14.1]
\label{II.8.2.14.1}
Let $T$ be a ring graded in positive degrees which is a $T_0$-algebra of finite
type. If the morphism corresponding to the injective homomorphism
$\bigoplus_{n\geq0}T_1^n\to T$ is everywhere defined and is an isomorphism
\[
  \operatorname{Proj}(T)\to
  \operatorname{Proj}\left(\bigoplus_{n\geq0}T_1^n\right),
\]
there exists $n_0$ such that $T_n=T_1^n$ for $n\geq n_0$.
\end{lemma}

Indeed, let $g_i$ ($1\leq i\leq r$) be generators of the $T_0$-module
$T_1$. The hypothesis first implies that the $D_+(g_i)$ cover
$\operatorname{Proj}(T)$ (\hyperref[II.2.8.1]{2.8.1}). Let
$(h_j)_{1\leq j\leq s}$ be a system of homogeneous elements of $T_+$, with
$\deg(h_j)=n_j$, which together with the $g_i$ form a system of generators of
the ideal $T_+$, or (what amounts to the same
(\hyperref[II.2.1.3]{2.1.3})) a system of generators of $T$ as a
$T_0$-algebra; if we set $T'=\bigoplus_{n\geq0}T_1^n$, the element
$h_j/g_i^{n_j}$ of the ring $T_{(g_i)}$ must by hypothesis belong to the
subring $T'_{(g_i)}$, so there exists an integer $k$ such that
$T_1^k h_j\subset T_1^{k+n_j}$ for every $j$. We conclude by induction
on $r$ that $T_1^k h_j^r\subset T'$ for every $r\geq1$, and by definition of the
$h_j$, we therefore have $T_1^kT\subset T'$. On the other hand, for every $j$
there is an integer $m_j$ such that $h_j^{m_j}$ belongs to the ideal of $T$ generated by
the $g_i$ (\hyperref[II.2.3.14]{2.3.14}), hence $h_j^{m_j}\in T_1T$, and

\oldpage[II]{162}
$h_j^{m_jk}\in T_1^kT\subset T'$. Consequently there is an integer $m_0\geq k$ such
that $h_j^m\in T_1^{mn_j}$ for $m\geq m_0$. This being so, if $q$ is the
largest of the integers $n_j$, the number $n_0=qsm_0+k$ answers the question.
Indeed, an element of $T_n$, for $n\geq n_0$, is a sum of monomials belonging
to $T_1^\alpha u$, where $u$ is a product of powers of the $h_j$; if
$\alpha\geq k$, it follows from the above that $T_1^\alpha u\subset T_1^n$;
otherwise, one of the exponents of the $h_j$ is $\geq m_0$, hence
$u\in T_1^\beta v$, where $\beta\geq k$ and $v$ is again a product of
powers of the $h_j$; we are then reduced to the preceding case, and we
finally see that $T_1^\alpha u\subset T_1^n$ in all cases.

\begin{remark}[8.2.15]
\label{II.8.2.15}
The condition $S_n=S_1^n$ for $n\geq n_0$ evidently implies
$S_{n+1}=S_1S_n$ for $n\geq n_0$, but the converse is false, even
when $S$ is assumed Noetherian. For example, let $K$ be a field,
$A=K[x]$, $B=K[y]/y^2K[y]$, where $x$ and $y$ are two indeterminates, $x$ being
taken of degree $1$ and $y$ of degree $2$, and let $S=A\otimes_KB$, so that
$S$ is a graded algebra over $K$ having a basis consisting of the elements $1$,
$x^n$ ($n\geq1$) and $x^ny$ ($n\geq0$). It is immediate that
$S_{n+1}=S_1S_n$ for $n\geq2$, but $S_1^n=Kx^n$ whereas
$S_n=Kx^n+Kx^{n-2}y$ for $n\geq2$.
\end{remark}

\subsection{Projecting cones}
\label{subsection:II.8.3}

\begin{env}[8.3.1]
\label{II.8.3.1}
Let $Y$ be a prescheme; throughout this subsection, we will be concerned only
with \emph{$Y$-preschemes} and \emph{$Y$-morphisms}. Let $\mathcal{S}$ be a
quasi-coherent $\mathcal{O}_Y$-algebra graded in \emph{positive degrees};
\emph{we will moreover suppose that $\mathcal{S}_0=\mathcal{O}_Y$}.
In accordance with the notation introduced in
(\hyperref[II.8.2.2]{8.2.2}), we set
\begin{equation}
\label{II.8.3.1.1}
  \widehat{\mathcal{S}}=\mathcal{S}[z]
  =\mathcal{S}\otimes_{\mathcal{O}_Y}\mathcal{O}_Y[z]
\tag{8.3.1.1}
\end{equation}
which we consider as an $\mathcal{O}_Y$-algebra graded in positive degrees by
defining the degrees by the formula
(\hyperref[II.8.2.2.2]{8.2.2.2}), so that for every affine open subset $U$
of $Y$, we have
\[
  \Gamma(U,\widehat{\mathcal{S}})=(\Gamma(U,\mathcal{S}))[z].
\]
In what follows, we set
\begin{equation}
\label{II.8.3.1.2}
  X=\operatorname{Proj}(\mathcal{S}),\qquad
  C=\operatorname{Spec}(\mathcal{S}),\qquad
  \widehat{C}=\operatorname{Proj}(\widehat{\mathcal{S}})
\tag{8.3.1.2}
\end{equation}
(where, in the definition of $C$, $\mathcal{S}$ is considered as an ungraded
$\mathcal{O}_Y$-algebra), and we will say that $C$ (resp.
$\widehat{C}$) is the \emph{affine cone} (resp. \emph{projective cone}) defined
by $\mathcal{S}$; we will also say ``cone'' instead of ``affine cone''. By
abuse of language, we will also say that $C$ (resp. $\widehat{C}$) is the
\emph{affine projecting cone of $X$} (resp. the \emph{projective projecting
cone of $X$}), it being understood that the prescheme $X$ is given in the form
$\operatorname{Proj}(\mathcal{S})$; finally, we will say that $\widehat{C}$ is
the \emph{projective closure} of $C$ (it being understood that the datum of
$\mathcal{S}$ is part of the structure of $C$).
\end{env}

\begin{proposition}[8.3.2]
\label{II.8.3.2}
There exist canonical $Y$-morphisms
\begin{equation}
\label{II.8.3.2.1}
  Y\xrightarrow{\varepsilon}C\xrightarrow{i}\widehat{C}
\tag{8.3.2.1}
\end{equation}
\begin{equation}
\label{II.8.3.2.2}
  X\xrightarrow{j}\widehat{C}
\tag{8.3.2.2}
\end{equation}

\oldpage[II]{163}
such that $\varepsilon$ and $j$ are closed immersions, and $i$ is an affine
morphism which is a dominant open immersion, for which
\begin{equation}
\label{II.8.3.2.3}
  i(C)=\widehat{C}-j(X)~;
\tag{8.3.2.3}
\end{equation}
moreover, $\widehat{C}$ is the smallest closed sub-prescheme of $\widehat{C}$
dominating $i(C)$.
\end{proposition}

To define $i$, we consider the open subset of $\widehat{C}$
\begin{equation}
\label{II.8.3.2.4}
  \widehat{C}_z=
  \operatorname{Spec}(\widehat{\mathcal{S}}/(z-1)\widehat{\mathcal{S}})
\tag{8.3.2.4}
\end{equation}
(\hyperref[II.3.1.4]{3.1.4}), $z$ being canonically identified with a
section of $\widehat{\mathcal{S}}$ over $Y$. The isomorphism
$i:C\xrightarrow{\sim}\widehat{C}_z$ then corresponds to the canonical
isomorphism (\hyperref[II.8.2.3.1]{8.2.3.1})
\[
  \widehat{\mathcal{S}}/(z-1)\widehat{\mathcal{S}}
  \xrightarrow{\sim}\mathcal{S}.
\]
The morphism $\varepsilon$ corresponds to the augmentation homomorphism
$\mathcal{S}\to\mathcal{S}_0=\mathcal{O}_Y$, with kernel
$\mathcal{S}_+$ (\hyperref[II.1.2.7]{1.2.7}), and since the latter is
surjective, $\varepsilon$ is a closed immersion
(\hyperref[II.1.4.10]{1.4.10}). Finally, $j$ similarly corresponds
(\hyperref[II.3.5.1]{3.5.1}) to the surjective homomorphism of degree $0$,
$\widehat{\mathcal{S}}\to\mathcal{S}$, which reduces to the identity on
$\mathcal{S}$ and is zero on $z\widehat{\mathcal{S}}$, which is its kernel;
$j$ is everywhere defined and is a closed immersion by virtue of
(\hyperref[II.3.6.2]{3.6.2}).

To prove the other assertions of
(\hyperref[II.8.3.2]{8.3.2}), we may evidently restrict to the case where
$Y=\operatorname{Spec}(A)$ is affine and $\mathcal{S}=\widetilde{S}$, where $S$
is a graded $A$-algebra, whence
$\widehat{\mathcal{S}}=(\widehat{S})^{\sim}$; the homogeneous elements $f$ of
$S_+$ are then identified with sections of $\widehat{\mathcal{S}}$ over
$Y$, and the open subset of $\widehat{C}$ denoted $D_+(f)$ in
(\hyperref[II.2.3.3]{2.3.3}) is also written $\widehat{C}_f$
(\hyperref[II.3.1.4]{3.1.4}); similarly, the open subset of $C$ denoted
$D(f)$ in (\hyperref[I.1.1.1]{I, 1.1.1}) is also written $C_f$
(\hyperref[0.5.5.2]{0, 5.5.2}). This being so, it follows from
(\hyperref[II.2.3.14]{2.3.14}) and from the definition of $\widehat{S}$ that
in the present case the open subsets $\widehat{C}_z=i(C)$ and $\widehat{C}_f$
($f$ homogeneous in $S_+$) form a \emph{covering} of
$\widehat{C}$. Moreover, with this notation, we have
\begin{equation}
\label{II.8.3.2.5}
  i^{-1}(\widehat{C}_f)=C_f~;
\tag{8.3.2.5}
\end{equation}
indeed,
$\widehat{C}_f\cap i(C)=\widehat{C}_f\cap\widehat{C}_z
=\widehat{C}_{fz}=\operatorname{Spec}(\widehat{S}_{(fz)})$. Now, if
$d=\deg(f)$, $\widehat{S}_{(fz)}$ is canonically isomorphic to
$(\widehat{S}_{(z)})_{f/z^d}$ (\hyperref[II.2.2.2]{2.2.2}), and it follows
from the definition of the isomorphism
(\hyperref[II.8.2.3.1]{8.2.3.1}) that the image of
$(\widehat{S}_{(z)})_{f/z^d}$ under the corresponding isomorphism of rings of
fractions is exactly $S_f$. Since $C_f=\operatorname{Spec}(S_f)$, this
proves (\hyperref[II.8.3.2.5]{8.3.2.5}) and shows at the same time that the
morphism $i$ is affine; moreover, the restriction of $i$ to $C_f$,
considered as a morphism into $\widehat{C}_f$, corresponds
(\hyperref[I.1.7.3]{I, 1.7.3}) to the canonical homomorphism
$\widehat{S}_{(f)}\to\widehat{S}_{(fz)}$, and by virtue of the above and of
(\hyperref[II.8.2.3.2]{8.2.3.2}), we can state the following result:

\begin{env}[8.3.2.6]
\label{II.8.3.2.6}
If $Y=\operatorname{Spec}(A)$ is affine and $\mathcal{S}=\widetilde{S}$,
then, for every homogeneous $f$ in $S_+$, $\widehat{C}_f$ is canonically
identified with $\operatorname{Spec}(S_f^{\leq})$, and the morphism
$C_f\to\widehat{C}_f$, the restriction of $i$, then corresponds to the canonical
injection $S_f^{\leq}\to S_f$.
\end{env}

Note now that (for $Y$ affine) the complement of
$\widehat{C}_z$ in $\widehat{C}=\operatorname{Proj}(\widehat{S})$

\oldpage[II]{164}
is, by definition, the set of graded prime ideals of $\widehat{S}$
containing $z$, that is, $j(X)$ by definition of $j$, which proves
(\hyperref[II.8.3.2.3]{8.3.2.3}).

Finally, to prove the last assertion of
(\hyperref[II.8.3.2]{8.3.2}), we may always suppose $Y$ affine. With
the preceding notation, note that in the ring $\widehat{S}$, $z$
is not a divisor of $0$; since $\overline{i(C)}=\widehat{C}$, it will suffice
to prove the following

\begin{lemma}[8.3.2.7]
\label{II.8.3.2.7}
Let $T$ be a ring graded in positive degrees,
$Z=\operatorname{Proj}(T)$, and $g$ a homogeneous element of $T$ of degree $d>0$.
If $g$ is not a divisor of $0$ in $T$, $Z$ is the smallest closed
sub-prescheme of $Z$ that dominates $Z_g=D_+(g)$.
\end{lemma}

By virtue of (\hyperref[I.4.1.9]{I, 4.1.9}) the question is local on $Z$;
for every homogeneous element $h\in T_e$ ($e>0$), it therefore suffices to prove that
$Z_h$ is the smallest closed sub-prescheme of $Z_h$ that dominates $Z_{gh}$;
it follows from the definitions and from (\hyperref[I.4.3.2]{I, 4.3.2}) that this
condition is equivalent to the fact that the canonical homomorphism
$T_{(h)}\to T_{(gh)}$ is \emph{injective}. Now this homomorphism is identified with
the canonical homomorphism
$T_{(h)}\to(T_{(h)})_{g^e/h^d}$ (\hyperref[II.2.2.3]{2.2.3}). But since
$g^e$ is not a divisor of $0$ in $T$, $g^e/h^d$ is not one in $T_h$
(nor, \emph{a fortiori}, in $T_{(h)}$), because the relation
$(g^e/h^d)(t/h^m)=0$ with $t\in T$ and $m>0$ implies the existence of an $n>0$
such that $h^ng^et=0$, whence $h^nt=0$, and consequently $t/h^m=0$ in $T_h$. This
therefore completes the proof (\hyperref[0.1.2.2]{0, 1.2.2}).

\begin{env}[8.3.3]
\label{II.8.3.3}
We will often identify the affine cone $C$ with the sub-prescheme induced by the
projective cone $\widehat{C}$ on the open subset $i(C)$, by means of the open
immersion $i$. The closed sub-prescheme of $C$ associated to the closed
immersion $\varepsilon$ is called the \emph{apex prescheme of $C$}; we also
say that $\varepsilon$, which is a $Y$-section of $C$, is the \emph{apex
section} or the \emph{zero section of $C$}; we can identify $Y$ with the
apex prescheme of $C$ by means of $\varepsilon$. Moreover,
$i\circ\varepsilon$ is a $Y$-section of $\widehat{C}$, hence also a
closed immersion (\hyperref[I.5.4.6]{I, 5.4.6}), corresponding to
the canonical surjective homomorphism of degree $0$,
$\widehat{\mathcal{S}}=\mathcal{S}[z]\to\mathcal{O}_Y[z]$
(cf. (\hyperref[II.3.1.7]{3.1.7})), with kernel
$\mathcal{S}_+[z]=\mathcal{S}_+\widehat{\mathcal{S}}$; the sub-prescheme of
$\widehat{C}$ associated to this closed immersion is also called the
\emph{apex prescheme of $\widehat{C}$}, and $i\circ\varepsilon$ the
\emph{apex section of $\widehat{C}$}; it can be identified with $Y$ by
$i\circ\varepsilon$. Finally, the closed sub-prescheme of $\widehat{C}$
associated to $j$ is called the \emph{locus at infinity of $\widehat{C}$}, and can
be identified with $X$ by means of $j$.
\end{env}

\begin{env}[8.3.4]
\label{II.8.3.4}
The sub-preschemes of $C$ (resp. $\widehat{C}$) induced respectively on
the \emph{open subsets}
\begin{equation}
\label{II.8.3.4.1}
  E=C-\varepsilon(Y),\qquad
  \widehat{E}=\widehat{C}-i(\varepsilon(Y))
\tag{8.3.4.1}
\end{equation}
are called respectively (by abuse of language) the \emph{punctured affine
cone} and the \emph{punctured projective cone} defined by $\mathcal{S}$; note
that in spite of this terminology, $E$ \emph{is not necessarily
affine} over $Y$, nor $\widehat{E}$ projective over $Y$
(cf. (\hyperref[II.8.4.3]{8.4.3})). When we identify $C$ with $i(C)$, we
therefore have, for the underlying spaces,
\begin{equation}
\label{II.8.3.4.2}
  C\cup\widehat{E}=\widehat{C},\qquad C\cap\widehat{E}=E
\tag{8.3.4.2}
\end{equation}
so that $\widehat{C}$ can be considered as obtained by
\emph{gluing} the open sub-preschemes $C$ and $\widehat{E}$; moreover,
by virtue of (\hyperref[II.8.3.2.3]{8.3.2.3}),
\begin{equation}
\label{II.8.3.4.3}
  E=\widehat{E}-j(X).
\tag{8.3.4.3}
\end{equation}

\oldpage[II]{165}
When $Y=\operatorname{Spec}(A)$ is affine, we have, with the notation of
(\hyperref[II.8.3.2]{8.3.2}),
\begin{equation}
\label{II.8.3.4.4}
  E=\bigcup C_f,\qquad
  \widehat{E}=\bigcup\widehat{C}_f,\qquad
  C_f=C\cap\widehat{C}_f
\tag{8.3.4.4}
\end{equation}
$f$ running through the set of homogeneous elements of $S_+$ (or only through a
subset $M$ of this set generating an ideal of $S_+$ whose radical in
$S_+$ is $S_+$ itself, in other words, such that the $X_f$ for $f\in M$
cover $X$ (\hyperref[II.2.3.14]{2.3.14})). The gluing of $C$ and
$\widehat{C}_f$ along $C_f$ is then determined by injection morphisms
$C_f\to C$, $C_f\to\widehat{C}_f$, which, as we have seen
(\hyperref[II.8.3.2.6]{8.3.2.6}), correspond respectively to the
canonical homomorphisms $S\to S_f$, $S_f^{\leq}\to S_f$.
\end{env}

\begin{proposition}[8.3.5]
\label{II.8.3.5}
With the notation of (\hyperref[II.8.3.1]{8.3.1}) and
(\hyperref[II.8.3.4]{8.3.4}), the morphism associated
(\hyperref[II.3.5.1]{3.5.1}) to the canonical injection
$\varphi:\mathcal{S}\to\widehat{\mathcal{S}}=\mathcal{S}[z]$ is a surjective
affine morphism (called the \emph{canonical retraction})
\begin{equation}
\label{II.8.3.5.1}
  p:\widehat{E}\to X
\tag{8.3.5.1}
\end{equation}
such that
\begin{equation}
\label{II.8.3.5.2}
  p\circ j=1_X.
\tag{8.3.5.2}
\end{equation}
\end{proposition}

To prove the proposition, we may restrict to the case where $Y$ is affine.
Taking into account the expression (\hyperref[II.8.3.4.4]{8.3.4.4}) of
$\widehat{E}$, the fact that the domain of definition $G(\varphi)$ of $p$ is
equal to $\widehat{E}$ will follow from the first of the following assertions:

\begin{env}[8.3.5.3]
\label{II.8.3.5.3}
If $Y=\operatorname{Spec}(A)$ is affine and $\mathcal{S}=\widetilde{S}$, we
have, for every homogeneous $f\in S_+$,
\begin{equation}
\label{II.8.3.5.4}
  p^{-1}(X_f)=\widehat{C}_f
\tag{8.3.5.4}
\end{equation}
and the restriction of $p$ to
$\widehat{C}_f=\operatorname{Spec}(S_f^{\leq})$, considered as a morphism from
$\widehat{C}_f$ to $X_f$, corresponds to the canonical injection
$S_{(f)}\to S_f^{\leq}$. If moreover $f\in S_1$, $\widehat{C}_f$ is isomorphic
to $X_f\otimes_{\mathbf{Z}}\mathbf{Z}[T]$ ($T$ an indeterminate).
\end{env}

Indeed, the formula (\hyperref[II.8.3.5.4]{8.3.5.4}) is only a
particular case of (\hyperref[II.2.8.1.1]{2.8.1.1}), and the second assertion
is none other than the definition of $\operatorname{Proj}(\varphi)$ when $Y$
is affine (\hyperref[II.2.8.1]{2.8.1}). Then the formula
(\hyperref[II.8.3.5.2]{8.3.5.2}) and the fact that $p$ is surjective follow
from the fact that the composite $\mathcal{S}\to\widehat{\mathcal{S}}\to\mathcal{S}$ of the
canonical homomorphisms is the identity on $\mathcal{S}$. Finally, the
last assertion of (\hyperref[II.8.3.5.3]{8.3.5.3}) comes from the fact that
$S_f^{\leq}$ is isomorphic to $S_{(f)}[T]$ when $f\in S_1$
(\hyperref[II.2.2.1]{2.2.1}).

\begin{corollary}[8.3.6]
\label{II.8.3.6}
The restriction
\begin{equation}
\label{II.8.3.6.1}
  \pi:E\to X
\tag{8.3.6.1}
\end{equation}
of $p$ to $E$ is a surjective affine morphism. If $Y$ is affine and $f$ is
homogeneous in $S_+$, we have
\begin{equation}
\label{II.8.3.6.2}
  \pi^{-1}(X_f)=C_f
\tag{8.3.6.2}
\end{equation}
and the restriction of $\pi$ to $C_f$ corresponds to the canonical injection
$S_{(f)}\to S_f$. If moreover $f\in S_1$, $C_f$ is isomorphic to
$X_f\otimes_{\mathbf{Z}}\mathbf{Z}[T,T^{-1}]$ ($T$ an indeterminate).
\end{corollary}

The formula (\hyperref[II.8.3.6.2]{8.3.6.2}) follows immediately from
(\hyperref[II.8.3.5.3]{8.3.5.3}) and
(\hyperref[II.8.3.2.5]{8.3.2.5}), and proves the surjectivity of $\pi$; we
have also seen that the immersion $i$, restricted to $C_f$, corresponds

\oldpage[II]{166}
to the injection $S_f^{\leq}\to S_f$ (\hyperref[II.8.3.2]{8.3.2}). Finally, the
last assertion is a consequence of the fact that for $f\in S_1$, $S_f$ is
isomorphic to $S_{(f)}[T,T^{-1}]$ (\hyperref[II.2.2.1]{2.2.1}).

\begin{remark}[8.3.7]
\label{II.8.3.7}
When $Y$ is affine, the elements of the underlying space of $E$ are the
prime ideals $\mathfrak{p}$ (not necessarily graded) of $S$ not
containing $S_+$, by virtue of the definition of the immersion $\varepsilon$
(\hyperref[II.8.3.2]{8.3.2}). For such an ideal $\mathfrak{p}$, the
$\mathfrak{p}\cap S_n$ evidently satisfy the conditions of
(\hyperref[II.2.1.9]{2.1.9}), so there exists a unique \emph{graded} prime
ideal $\mathfrak{q}$ of $S$ such that
$\mathfrak{q}\cap S_n=\mathfrak{p}\cap S_n$ for every $n$; the map
$\pi:E\to X$ of underlying spaces is then interpreted by means of the relation
\begin{equation}
\label{II.8.3.7.1}
  \pi(\mathfrak{p})=\mathfrak{q}.
\tag{8.3.7.1}
\end{equation}
Indeed, to verify this relation, it suffices to consider an $f$
homogeneous in $S_+$ such that $\mathfrak{p}\in D(f)$, and to observe that
$\mathfrak{q}_{(f)}$ is the inverse image of $\mathfrak{p}_f$ under
the injection $S_{(f)}\to S_f$.
\end{remark}

\begin{corollary}[8.3.8]
\label{II.8.3.8}
If $\mathcal{S}$ is generated by $\mathcal{S}_1$, the morphisms $p$ and
$\pi$ are of finite type; for every $x\in X$, the fibre $p^{-1}(x)$ is
isomorphic to $\operatorname{Spec}(k(x)[T])$ and the fibre $\pi^{-1}(x)$ is
isomorphic to $\operatorname{Spec}(k(x)[T,T^{-1}])$.
\end{corollary}

This follows immediately from (\hyperref[II.8.3.5]{8.3.5}) and
(\hyperref[II.8.3.6]{8.3.6}), observing that, when $Y$ is affine and
$S$ is generated by $S_1$, the $X_f$, for $f\in S_1$, form a covering
of $X$ (\hyperref[II.2.3.14]{2.3.14}).

\begin{remark}[8.3.9]
\label{II.8.3.9}
The punctured affine cone corresponding to the graded $\mathcal{O}_Y$-algebra
$\mathcal{O}_Y[T]$ ($T$ an indeterminate) is identified with
$G_m=\operatorname{Spec}(\mathcal{O}_Y[T,T^{-1}])$, since it is none other than
$C_T$, as we saw in (\hyperref[II.8.3.2]{8.3.2}) (see
(\hyperref[II.8.4.4]{8.4.4}) for a more general result). This prescheme
is canonically endowed with a structure of ``\emph{commutative $Y$-group
scheme}''. This is a notion that will be developed in detail later,
and that can be summarized briefly here as follows. A $Y$-group
scheme is a $Y$-scheme $G$ endowed with two $Y$-morphisms
$p:G\times_YG\to G$ and $s:G\to G$ that satisfy the conditions formally
analogous to the axioms of the composition law and of the symmetry map
in a group: the diagram
\[
  \xymatrix{
    G\times G\times G \ar[r]^{p\times1}\ar[d]_{1\times p}
      & G\times G \ar[d]^{p} \\
    G\times G \ar[r]_{p} & G
  }
\]
must be commutative (``associativity''), and we must have a condition which
corresponds, for groups, to the fact that the maps
\[
  (x,y)\mapsto(x,x^{-1},y)\mapsto(x,x^{-1}y)\mapsto x(x^{-1}y)
\]
and
\[
  (x,y)\mapsto(x,x^{-1},y)\mapsto(x,yx^{-1})\mapsto(yx^{-1})x
\]
must both reduce to $(x,y)\mapsto y$; the sequence of morphisms
corresponding, for example, to the first composite map is
\[
  G\times G\xrightarrow{(1,s)\times1}G\times G\times G
  \xrightarrow{1\times p}G\times G\xrightarrow{p}G
\]
and the reader will write the second sequence similarly.

\oldpage[II]{167}
It is immediate (\hyperref[I.3.4.3]{I, 3.4.3}) that giving a structure
of $Y$-group scheme on a $Y$-scheme $G$ is equivalent to giving, for
\emph{every} $Y$-prescheme $Z$, a \emph{group} structure on
the set $\operatorname{Hom}_Y(Z,G)$, these structures having to be such that
for every $Y$-morphism $Z\to Z'$, the corresponding map
$\operatorname{Hom}_Y(Z',G)\to\operatorname{Hom}_Y(Z,G)$ is a
group homomorphism. In the particular case of $G_m$ that we are
considering here, $\operatorname{Hom}_Y(Z,G)$ is identified with the set of
$Z$-sections of $Z\times_YG_m$ (\hyperref[I.3.3.14]{I, 3.3.14}), hence with
the set of $Z$-sections of
$\operatorname{Spec}(\mathcal{O}_Z[T,T^{-1}])$; finally, the same
reasoning as in (\hyperref[I.3.3.15]{I, 3.3.15}) shows that this
set is canonically identified with the set of \emph{invertible} elements
of the ring $\Gamma(Z,\mathcal{O}_Z)$, and the group structure on this
set is the structure coming from multiplication in the ring
$\Gamma(Z,\mathcal{O}_Z)$. The reader may verify that the morphisms $p$
and $s$ considered above are obtained in the following way: by
(\hyperref[II.1.2.7]{1.2.7}) and
(\hyperref[II.1.4.6]{1.4.6}) they correspond to homomorphisms of $\mathcal{O}_Y$-algebras
\begin{align*}
  \pi &: \mathcal{O}_Y[T,T^{-1}]
  \to\mathcal{O}_Y[T,T^{-1},T',T'^{-1}],\\
  \sigma &: \mathcal{O}_Y[T,T^{-1}]\to\mathcal{O}_Y[T,T^{-1}]
\end{align*}
and are entirely determined by giving $\pi(T)=TT'$ and
$\sigma(T)=T^{-1}$.

This being so, $G_m$ can be considered as a ``\emph{universal domain of
operators}'' for every \emph{affine cone}
$C=\operatorname{Spec}(\mathcal{S})$, where $\mathcal{S}$ is a
quasi-coherent $\mathcal{O}_Y$-algebra graded in positive degrees. This
means that we can canonically define a $Y$-morphism
$G_m\times_YC\to C$ which has the formal properties of an external law of a
set endowed with a group of operators; or again, as above for
group schemes, we give ourselves, for every $Y$-prescheme $Z$, an external law
on $\operatorname{Hom}_Y(Z,C)$ having the group
$\operatorname{Hom}_Y(Z,G_m)$ as set of operators, with the usual
axioms of sets endowed with a group of operators, and a compatibility
condition with the $Y$-morphisms $Z\to Z'$. In the present case, the
morphism $G_m\times_YC\to C$ is defined by giving the homomorphism of
$\mathcal{O}_Y$-algebras
\[
  \mathcal{S}\to
  \mathcal{S}\otimes_{\mathcal{O}_Y}\mathcal{O}_Y[T,T^{-1}]
  =\mathcal{S}[T,T^{-1}]
\]
which assigns to every section $s_n\in\Gamma(U,\mathcal{S}_n)$ ($U$ an open
subset of $Y$) the section
$s_nT^n\in\Gamma(U,\mathcal{S}\otimes_{\mathcal{O}_Y}
\mathcal{O}_Y[T,T^{-1}])$.

Conversely, suppose given a quasi-coherent $\mathcal{O}_Y$-algebra
$\mathcal{S}$, \emph{not graded a priori}, and, on
$C=\operatorname{Spec}(\mathcal{S})$, a structure of ``$Y$-scheme in
sets \emph{endowed with groups of operators}'' having as domain
of operators the $Y$-group scheme $G_m$; then we canonically deduce from it
a \emph{grading} of $\mathcal{O}_Y$-algebra on
$\mathcal{S}$. Indeed, giving a $Y$-morphism $G_m\times_YC\to C$
is equivalent to giving a homomorphism of $\mathcal{O}_Y$-algebras
$\psi:\mathcal{S}\to\mathcal{S}[T,T^{-1}]$, which can therefore be written
$\psi=\sum_{n\in\mathbf{Z}}\psi_nT^n$, where the
$\psi_n:\mathcal{S}\to\mathcal{S}$ are homomorphisms of
$\mathcal{O}_Y$-modules (for every section $s\in\Gamma(U,\mathcal{S})$, $U$ open in $Y$,
and every point $x$ of $U$, there exists an open neighbourhood $V$ of $x$ in $U$
such that $\psi_n(s)|_V=0$ except for finitely many indices $n$). One can
then verify that the axioms of sets endowed with a group of operators
imply that the $\psi_n(\mathcal{S})=\mathcal{S}_n$ define a
grading (in positive or negative degrees) of $\mathcal{O}_Y$-algebra on
$\mathcal{S}$, the $\psi_n$ being the corresponding projectors. We thus have
a notion of ``\emph{affine cone}'' structure on every affine $Y$-scheme,
defined in a ``geometric'' way, without recourse to a previously given grading.

\oldpage[II]{168}
We will not develop this point of view further here, and we leave it to the
reader to formulate precisely the definitions and results that
correspond to the preceding indications.
\end{remark}

\subsection{Projective closure of a vector bundle}
\label{subsection:II.8.4}

\begin{env}[8.4.1]
\label{II.8.4.1}
Let $Y$ be a prescheme and $\mathcal{E}$ a quasi-coherent $\mathcal{O}_Y$-module.
If we take for $\mathcal{S}$ the graded $\mathcal{O}_Y$-algebra
$\mathbf{S}_{\mathcal{O}_Y}(\mathcal{E})$, the definition
(\hyperref[II.8.3.1.1]{8.3.1.1}) shows that
$\widehat{\mathcal{S}}$ is identified with
$\mathbf{S}_{\mathcal{O}_Y}(\mathcal{E}\oplus\mathcal{O}_Y)$. Since the affine cone
$\operatorname{Spec}(\mathcal{S})$ defined by $\mathcal{S}$ is then by
definition $\mathbf{V}(\mathcal{E})$, and $\operatorname{Proj}(\mathcal{S})$
is by definition $\mathbf{P}(\mathcal{E})$, we therefore see that:
\end{env}

\begin{proposition}[8.4.2]
\label{II.8.4.2}
The projective closure of a vector bundle $\mathbf{V}(\mathcal{E})$ over
$Y$ is canonically isomorphic to
$\mathbf{P}(\mathcal{E}\oplus\mathcal{O}_Y)$, and the locus at infinity of the
latter is canonically isomorphic to $\mathbf{P}(\mathcal{E})$.
\end{proposition}

\begin{remark}[8.4.3]
\label{II.8.4.3}
Take for example $\mathcal{E}=\mathcal{O}_Y^r$ with $r\geq2$; then the
punctured cones $E$, $\widehat{E}$ defined by $\mathcal{S}$ are neither affine
nor projective over $Y$ if $Y\neq\varnothing$. The second assertion is
immediate, because $\widehat{C}=\mathbf{P}(\mathcal{O}_Y^{r+1})$ is projective over
$Y$ and the underlying spaces of $E$ and $\widehat{E}$ are non-closed open
subsets of $\widehat{C}$, so the canonical immersions
$E\to\widehat{E}$ and $\widehat{E}\to\widehat{C}$ are not projective
(\hyperref[II.5.5.3]{5.5.3}), and we conclude by
(\hyperref[II.5.5.5]{5.5.5, (v)}). On the other hand, supposing
$Y=\operatorname{Spec}(A)$ affine and, for example, $r=2$, we have
$C=\operatorname{Spec}(A[T_1,T_2])$, and $E$ is then the prescheme induced by
$C$ on the open subset $D(T_1)\cup D(T_2)$; now we have seen that the latter is
not affine (\hyperref[I.5.5.11]{I, 5.5.11});
\emph{a fortiori} $\widehat{E}$ cannot be affine, since $E$ is the open subset
where the section $z$ over $\widehat{E}$ does not vanish
(\hyperref[II.8.3.2]{8.3.2}).

However:
\end{remark}

\begin{proposition}[8.4.4]
\label{II.8.4.4}
If $\mathcal{L}$ is an invertible $\mathcal{O}_Y$-module, we have
canonical isomorphisms for the punctured cones $E$ and $\widehat{E}$
corresponding to $C=\mathbf{V}(\mathcal{L})$,
\begin{equation}
\label{II.8.4.4.1}
  \operatorname{Spec}\left(
    \bigoplus_{n\in\mathbf{Z}}\mathcal{L}^{\otimes n}
  \right)\xrightarrow{\sim}E
\tag{8.4.4.1}
\end{equation}
\begin{equation}
\label{II.8.4.4.2}
  \mathbf{V}(\mathcal{L}^{-1})\xrightarrow{\sim}\widehat{E}.
\tag{8.4.4.2}
\end{equation}

Moreover, there exists a canonical isomorphism of the projective closure of
$\mathbf{V}(\mathcal{L})$ onto that of $\mathbf{V}(\mathcal{L}^{-1})$,
transforming the zero section (resp. the locus at infinity) of the first into the
locus at infinity (resp. the zero section) of the second.
\end{proposition}

\begin{proof}
Here we have
$\mathcal{S}=\bigoplus_{n\geq0}\mathcal{L}^{\otimes n}$; the canonical
injection
\[
  \mathcal{S}\to\bigoplus_{n\in\mathbf{Z}}\mathcal{L}^{\otimes n}
\]
defines a canonical dominant morphism
\begin{equation}
\label{II.8.4.4.3}
  \operatorname{Spec}\left(
    \bigoplus_{n\in\mathbf{Z}}\mathcal{L}^{\otimes n}
  \right)
  \to\mathbf{V}(\mathcal{L})
  =\operatorname{Spec}\left(
    \bigoplus_{n\geq0}\mathcal{L}^{\otimes n}
  \right)
\tag{8.4.4.3}
\end{equation}
and it suffices to prove that this morphism is an isomorphism of the scheme
$\operatorname{Spec}(\bigoplus_{n\in\mathbf{Z}}\mathcal{L}^{\otimes n})$
onto $E$. The question being local on $Y$, we may suppose that
$Y=\operatorname{Spec}(A)$ is affine

\oldpage[II]{169}
and $\mathcal{L}=\mathcal{O}_Y$, hence $\mathcal{S}=(A[T])^{\sim}$ and
$\bigoplus_{n\in\mathbf{Z}}\mathcal{L}^{\otimes n}
=(A[T,T^{-1}])^{\sim}$. Now $A[T,T^{-1}]$ is the ring of fractions
$A[T]_T$ of $A[T]$, so (\hyperref[II.8.4.4.3]{8.4.4.3}) identifies the
left-hand side with the prescheme induced by $C=\mathbf{V}(\mathcal{L})$ on
the open subset $D(T)$; the complement $V(T)$ of this open subset in $C$ is
the underlying space of the closed sub-prescheme of $C$ defined by the ideal
$TA[T]$, that is, the zero section of $C$, hence $E=D(T)$.

On the other hand, the isomorphism (\hyperref[II.8.4.4.2]{8.4.4.2}) will be a
consequence of the last assertion, $\mathbf{V}(\mathcal{L}^{-1})$ being the
complement of the locus at infinity of its projective closure and $\widehat{E}$
the complement of the zero section of the projective closure of
$C=\mathbf{V}(\mathcal{L})$. Now these projective closures are
respectively $\mathbf{P}(\mathcal{L}^{-1}\oplus\mathcal{O}_Y)$ and
$\mathbf{P}(\mathcal{L}\oplus\mathcal{O}_Y)$; but we can write
$\mathcal{L}\oplus\mathcal{O}_Y
=\mathcal{L}\otimes(\mathcal{L}^{-1}\oplus\mathcal{O}_Y)$. The existence of
the desired canonical isomorphism then follows from
(\hyperref[II.4.1.4]{4.1.4}), and everything comes down to seeing that this isomorphism
exchanges zero sections and loci at infinity. For this we may restrict to the
case where $Y=\operatorname{Spec}(A)$ is affine,
$L=Ac$, $L^{-1}=Ac'$, the canonical isomorphism
$L\otimes L^{-1}\to A$ assigning to $c\otimes c'$ the element $1$
of $A$. Then $\mathbf{S}(L\oplus A)$ is the tensor product of $A[z]$ and
$\bigoplus_{n\geq0}Ac^{\otimes n}$, $\mathbf{S}(L^{-1}\oplus A)$ the tensor
product of $A[z]$ and $\bigoplus_{n\geq0}Ac'^{\otimes n}$, and
the isomorphism defined in (\hyperref[II.4.1.4]{4.1.4}) assigns to
$z^h\otimes c'^{\otimes(n-h)}$ the element
$z^{n-h}\otimes c^{\otimes h}$. Now, in
$\mathbf{P}(\mathcal{L}^{-1}\oplus\mathcal{O}_Y)$, the locus at infinity is
the set of points where the section $z$ vanishes, and the zero section is
the set of points where the section $c'$ vanishes; since we have analogous
definitions for $\mathbf{P}(\mathcal{L}\oplus\mathcal{O}_Y)$, our conclusion
follows immediately from the explicit description above.
\end{proof}

\subsection{Functorial behaviour}
\label{subsection:II.8.5}

\begin{env}[8.5.1]
\label{II.8.5.1}
Let $Y$, $Y'$ be two preschemes, $q:Y'\to Y$ a morphism, and
$\mathcal{S}$ (resp. $\mathcal{S}'$) a quasi-coherent $\mathcal{O}_Y$-algebra
(resp. $\mathcal{O}_{Y'}$-algebra) graded in
\emph{positive} degrees. Consider a $q$-morphism of graded algebras
\begin{equation}
\label{II.8.5.1.1}
  \varphi:\mathcal{S}\to\mathcal{S}'.
\tag{8.5.1.1}
\end{equation}
We know (\hyperref[II.1.5.6]{1.5.6}) that there corresponds to it canonically a
morphism
\[
  \Phi=\operatorname{Spec}(\varphi):
  \operatorname{Spec}(\mathcal{S}')\to\operatorname{Spec}(\mathcal{S})
\]
such that the diagram
\begin{equation}
\label{II.8.5.1.2}
  \xymatrix{
    C'\ar[r]^{\Phi}\ar[d] & C\ar[d]\\
    Y'\ar[r]_{q} & Y
  }
\tag{8.5.1.2}
\end{equation}
where we have set $C=\operatorname{Spec}(\mathcal{S})$,
$C'=\operatorname{Spec}(\mathcal{S}')$, is commutative.
\emph{Suppose moreover that
$\mathcal{S}_0=\mathcal{O}_Y$, $\mathcal{S}'_0=\mathcal{O}_{Y'}$}; let
$\varepsilon:Y\to C$ and $\varepsilon':Y'\to C'$ be the canonical immersions
(\hyperref[II.8.3.2]{8.3.2}); we then have a commutative diagram
\begin{equation}
\label{II.8.5.1.3}
  \xymatrix{
    Y'\ar[r]^{q}\ar[d]_{\varepsilon'} & Y\ar[d]^{\varepsilon}\\
    C'\ar[r]_{\Phi} & C
  }
\tag{8.5.1.3}
\end{equation}

\oldpage[II]{170}
which corresponds to the diagram
\[
  \xymatrix{
    \mathcal{S}\ar[r]^{\varphi}\ar[d] & \mathcal{S}'\ar[d]\\
    \mathcal{O}_Y\ar[r] & \mathcal{O}_{Y'}
  }
\]
where the vertical arrows are the augmentation homomorphisms, and whose
commutativity follows from the hypothesis that $\varphi$ is assumed to be a
homomorphism of \emph{graded} algebras.
\end{env}

\begin{proposition}[8.5.2]
\label{II.8.5.2}
If $E$ (resp. $E'$) is the punctured affine cone defined by $\mathcal{S}$
(resp. $\mathcal{S}'$), we have $\Phi^{-1}(E)\subset E'$; if moreover
$\operatorname{Proj}(\varphi):G(\varphi)\to\operatorname{Proj}(\mathcal{S})$
is everywhere defined (in other words if
$G(\varphi)=\operatorname{Proj}(\mathcal{S}')$), then we have
$\Phi^{-1}(E)=E'$, and conversely.
\end{proposition}

\begin{proof}
The first assertion follows from the commutativity of
(\hyperref[II.8.5.1.3]{8.5.1.3}). To prove the second, we may
restrict to the case where $Y=\operatorname{Spec}(A)$ and
$Y'=\operatorname{Spec}(A')$ are affine,
$\mathcal{S}=\widetilde{S}$, $\mathcal{S}'=\widetilde{S'}$. For every
homogeneous $f$ in $S_+$, if we set $f'=\varphi(f)$, we have
$\Phi^{-1}(C_f)=C'_{f'}$ (\hyperref[I.2.2.4.1]{I, 2.2.4.1}); to say that
$G(\varphi)=\operatorname{Proj}(S')$ means that in $S'_+$ the radical of
the ideal generated by the $f'=\varphi(f)$ is $S'_+$ itself
(\hyperref[II.2.8.1]{2.8.1}) and (\hyperref[II.2.3.14]{2.3.14}), and this
is again equivalent to saying that the $C'_{f'}$ cover $E'$
(\hyperref[II.8.3.4.4]{8.3.4.4}).
\end{proof}

\begin{env}[8.5.3]
\label{II.8.5.3}
The $q$-morphism $\varphi$ extends canonically to a $q$-morphism
of graded algebras
\begin{equation}
\label{II.8.5.3.1}
  \widehat{\varphi}:\widehat{\mathcal{S}}
  \to\widehat{\mathcal{S}'}
\tag{8.5.3.1}
\end{equation}
by setting $\widehat{\varphi}(z)=z$. Consequently we deduce from it a morphism
\[
  \widehat{\Phi}=\operatorname{Proj}(\widehat{\varphi}):
  G(\widehat{\varphi})\to\widehat{C}
  =\operatorname{Proj}(\widehat{\mathcal{S}})
\]
such that the diagram
\[
  \xymatrix{
    G(\widehat{\varphi})\ar[r]^{\widehat{\Phi}}\ar[d] &
    \widehat{C}\ar[d]\\
    Y'\ar[r]_{q} & Y
  }
\]
is commutative (\hyperref[II.3.5.6]{3.5.6}). It follows immediately from the
definitions that if we denote by $i:C\to\widehat{C}$ and
$i':C'\to\widehat{C'}$ the canonical open immersions
(\hyperref[II.8.3.2]{8.3.2}), we have
$i'(C')\subset G(\widehat{\varphi})$ and the diagram
\begin{equation}
\label{II.8.5.3.2}
  \xymatrix{
    C'\ar[r]^{\Phi}\ar[d]_{i'} & C\ar[d]^{i}\\
    G(\widehat{\varphi})\ar[r]_{\widehat{\Phi}} & \widehat{C}
  }
\tag{8.5.3.2}
\end{equation}
is commutative. Finally, if we set
$X=\operatorname{Proj}(\mathcal{S})$, $X'=\operatorname{Proj}(\mathcal{S}')$,
and if $j:X\to\widehat{C}$, $j':X'\to\widehat{C'}$ are the canonical closed
immersions (\hyperref[II.8.3.2]{8.3.2}), it again follows from the
definitions of these immersions that
$j'(G(\varphi))\subset G(\widehat{\varphi})$ and that the diagram

\oldpage[II]{171}
\begin{equation}
\label{II.8.5.3.3}
  \xymatrix{
    G(\varphi)\ar[r]^{\operatorname{Proj}(\varphi)}\ar[d]_{j'} &
    X\ar[d]^{j}\\
    G(\widehat{\varphi})\ar[r]_{\widehat{\Phi}} & \widehat{C}
  }
\tag{8.5.3.3}
\end{equation}
is commutative.
\end{env}

\begin{proposition}[8.5.4]
\label{II.8.5.4}
If $\widehat{E}$ (resp. $\widehat{E'}$) is the punctured projective cone defined
by $\mathcal{S}$ (resp. $\mathcal{S}'$), we have
$\widehat{\Phi}^{-1}(\widehat{E})\subset\widehat{E'}$; moreover, if
$p:\widehat{E}\to X$ and $p':\widehat{E'}\to X'$ are the canonical
retractions, we have
$p'(\widehat{\Phi}^{-1}(\widehat{E}))\subset G(\varphi)$, and the diagram
\begin{equation}
\label{II.8.5.4.1}
  \xymatrix{
    \widehat{\Phi}^{-1}(\widehat{E})
      \ar[r]^{\widehat{\Phi}}\ar[d]_{p'} &
    \widehat{E}\ar[d]^{p}\\
    G(\varphi)\ar[r]_{\operatorname{Proj}(\varphi)} & X
  }
\tag{8.5.4.1}
\end{equation}
is commutative. If $\operatorname{Proj}(\varphi)$ is everywhere defined, so is
$\widehat{\Phi}$, and we have
$\widehat{\Phi}^{-1}(\widehat{E})=\widehat{E'}$.
\end{proposition}

\begin{proof}
The first assertion follows from the commutativity of the diagrams
(\hyperref[II.8.5.1.3]{8.5.1.3}) and
(\hyperref[II.8.5.3.2]{8.5.3.2}), and the next two from the definition
of the canonical retractions (\hyperref[II.8.3.5]{8.3.5}) and from the definition
of $\widehat{\varphi}$. On the other hand, to see that $\widehat{\Phi}$ is
everywhere defined when $\operatorname{Proj}(\varphi)$ is, we
may restrict to considering the case where $Y=\operatorname{Spec}(A)$ and
$Y'=\operatorname{Spec}(A')$ are affine,
$\mathcal{S}=\widetilde{S}$, $\mathcal{S}'=\widetilde{S'}$; the hypothesis is
that, when $f$ runs through the set of homogeneous elements of $S_+$, the ideal
generated in $S'_+$ by the $\varphi(f)$ has as radical in $S'_+$ the ideal
$S'_+$ itself; we conclude immediately that the radical in $(S'[z])_+$ of
the ideal generated by $z$ and the $\varphi(f)$ is $(S'[z])_+$ itself, whence
our assertion; this proves at the same time that $\widehat{E'}$ is the union of the
$\widehat{C'}_{\varphi(f)}$, hence equal to
$\widehat{\Phi}^{-1}(\widehat{E})$.
\end{proof}

\begin{corollary}[8.5.5]
\label{II.8.5.5}
When $\operatorname{Proj}(\varphi)$ is everywhere defined, the inverse image
under $\widehat{\Phi}$ of the underlying space of the locus at infinity (resp. of the
apex prescheme) of $\widehat{C}$ is the underlying space of the locus at
infinity (resp. of the apex prescheme) of $\widehat{C'}$.
\end{corollary}

\begin{proof}
This follows immediately from (\hyperref[II.8.5.4]{8.5.4}) and
(\hyperref[II.8.5.2]{8.5.2}), taking into account the relations
(\hyperref[II.8.3.4.1]{8.3.4.1}) and
(\hyperref[II.8.3.4.2]{8.3.4.2}).
\end{proof}

\subsection{A canonical isomorphism for punctured cones}
\label{subsection:II.8.6}

\begin{env}[8.6.1]
\label{II.8.6.1}
Let $Y$ be a prescheme, $\mathcal{S}$ a quasi-coherent $\mathcal{O}_Y$-algebra
graded in positive degrees \emph{such that
$\mathcal{S}_0=\mathcal{O}_Y$}, and $X$ the $Y$-scheme
$\operatorname{Proj}(\mathcal{S})$. We are going to apply the results of
(\hyperref[subsection:II.8.5]{8.5}) to the case where we take $Y'=X$,
$q:X\to Y$ being the structural morphism; let
\begin{equation}
\label{II.8.6.1.1}
  \mathcal{S}_X=\bigoplus_{n\in\mathbf{Z}}\mathcal{O}_X(n)
\tag{8.6.1.1}
\end{equation}

\oldpage[II]{172}
which is a quasi-coherent graded $\mathcal{O}_X$-algebra, the multiplication
on it being defined by means of the canonical homomorphisms
(\hyperref[II.3.2.6.1]{3.2.6.1})
\[
  \mathcal{O}_X(m)\otimes_{\mathcal{O}_X}\mathcal{O}_X(n)
  \to\mathcal{O}_X(m+n)
\]
whose associativity is guaranteed by the commutative diagram
(\hyperref[II.2.5.11.4]{2.5.11.4}). We take for $\mathcal{S}'$ the
quasi-coherent graded sub-$\mathcal{O}_X$-algebra
$\mathcal{S}_X^{\geq}=\bigoplus_{n\geq0}\mathcal{O}_X(n)$ of
$\mathcal{S}_X$, with positive degrees.

Finally, we will consider the canonical $q$-morphism
\begin{equation}
\label{II.8.6.1.2}
  \alpha:\mathcal{S}\to\mathcal{S}_X^{\geq}
\tag{8.6.1.2}
\end{equation}
defined in (\hyperref[II.3.3.2.3]{3.3.2.3}) as a homomorphism
$\mathcal{S}\to q_*(\mathcal{S}_X)$, but which evidently maps
$\mathcal{S}$ into $q_*(\mathcal{S}_X^{\geq})$. We set
\begin{equation}
\label{II.8.6.1.3}
  C_X=\operatorname{Spec}(\mathcal{S}_X^{\geq}),\qquad
  \widehat{C}_X=\operatorname{Proj}(\mathcal{S}_X^{\geq}[z]),\qquad
  X'=\operatorname{Proj}(\mathcal{S}_X^{\geq})
\tag{8.6.1.3}
\end{equation}
and we denote by $E_X$ and $\widehat{E}_X$ the corresponding punctured affine
cone and punctured projective cone; we denote by
$\varepsilon_X:X\to C_X$, $i_X:C_X\to\widehat{C}_X$,
$j_X:X'\to\widehat{C}_X$, $p_X:\widehat{E}_X\to X'$,
$\pi_X:E_X\to X'$ the canonical morphisms defined in
(\hyperref[subsection:II.8.3]{8.3}).
\end{env}

\begin{proposition}[8.6.2]
\label{II.8.6.2}
The structural morphism $u:X'\to X$ is an isomorphism, and the morphism
$\operatorname{Proj}(\alpha)$ is everywhere defined and identical to $u$. The
morphism
$\operatorname{Proj}(\widehat{\alpha}):\widehat{C}_X\to\widehat{C}$ is
everywhere defined, and its restrictions to $\widehat{E}_X$ and $E_X$ are
isomorphisms onto $\widehat{E}$ and $E$ respectively. Finally, when $X'$ and
$X$ are identified by means of $u$, the morphisms $p_X$ and $\pi_X$
are identified with the structural morphisms of the $X$-preschemes $\widehat{E}_X$ and
$E_X$.
\end{proposition}

\begin{proof}
We may evidently restrict to the case where $Y=\operatorname{Spec}(A)$ is affine
and $\mathcal{S}=\widetilde{S}$; then $X$ is the union of the affine open subsets
$X_f$, where $f$ runs through the set of homogeneous elements of $S_+$, the ring
of $X_f$ being $S_{(f)}$. Moreover, it follows from
(\hyperref[II.8.2.7.1]{8.2.7.1}) that
\begin{equation}
\label{II.8.6.2.1}
  \Gamma(X_f,\mathcal{S}_X^{\geq})=S_f^{\geq}.
\tag{8.6.2.1}
\end{equation}

We therefore have $u^{-1}(X_f)=\operatorname{Proj}(S_f^{\geq})$. But if
$f\in S_d$ ($d>0$), $\operatorname{Proj}(S_f^{\geq})$ is canonically
isomorphic to $\operatorname{Proj}((S_f^{\geq})^{(d)})$
(\hyperref[II.2.4.7]{2.4.7}), and on the other hand
$(S_f^{\geq})^{(d)}=(S^{(d)})_f^{\geq}$ is identified with $S_{(f)}[T]$
(\hyperref[II.2.2.1]{2.2.1}) by the map $T\to f/1$; we conclude
(\hyperref[II.3.1.7]{3.1.7}) that the structural morphism
$u^{-1}(X_f)\to X_f$ is an isomorphism, whence the first assertion. To
prove the second, note that the restriction
$u^{-1}(X_f)\cap G(\alpha)\to X=\operatorname{Proj}(S)$ of
$\operatorname{Proj}(\alpha)$ corresponds to the canonical map
$x\to x/1$ from $S$ to $S_f^{\geq}$
(\hyperref[II.2.6.2]{2.6.2}); we first deduce from this that
$G(\alpha)=X'$; then, taking into account that
$u^{-1}(X_f)=(u^{-1}(X_f))_{f/1}$, it follows from
(\hyperref[II.2.8.1.1]{2.8.1.1}) that the image of $u^{-1}(X_f)$ under
$\operatorname{Proj}(\alpha)$ is contained in $X_f$, and the restriction of
$\operatorname{Proj}(\alpha)$ to $u^{-1}(X_f)$, considered as a morphism
into $X_f=\operatorname{Spec}(S_{(f)})$, is indeed identical to that of $u$.
Finally, the formula (\hyperref[II.8.3.5.4]{8.3.5.4}), applied to $p_X$
instead of $p$, shows that
$p_X^{-1}(u^{-1}(X_f))=
\operatorname{Spec}((S_f^{\geq})_{f/1}^{\leq})$, and this open subset is, by
(\hyperref[II.8.5.4.1]{8.5.4.1}), the inverse image under
$\operatorname{Proj}(\widehat{\alpha})$ of
$p^{-1}(X_f)=\operatorname{Spec}(S_f^{\leq})$
(\hyperref[II.8.3.5.3]{8.3.5.3}). Taking into account
(\hyperref[II.8.2.3.2]{8.2.3.2}), the restriction of
$\operatorname{Proj}(\widehat{\alpha})$ to
$p_X^{-1}(u^{-1}(X_f))$ corresponds to the inverse isomorphism of
(\hyperref[II.8.2.7.2]{8.2.7.2}), restricted to $S_f^{\leq}$, whence the
third point; the last assertion is evident by definition.

\oldpage[II]{173}
Note also that it follows from the commutative diagram
(\hyperref[II.8.5.3.2]{8.5.3.2}) that \emph{the restriction to $C_X$ of
$\operatorname{Proj}(\widehat{\alpha})$ is none other than the morphism
$\operatorname{Spec}(\alpha)$}.
\end{proof}

\begin{corollary}[8.6.3]
\label{II.8.6.3}
Considered as $X$-schemes, $\widehat{E}_X$ is canonically isomorphic to
$\operatorname{Spec}(\mathcal{S}_X^{\leq})$, and $E_X$ to
$\operatorname{Spec}(\mathcal{S}_X)$.
\end{corollary}

\begin{proof}
Since we know that the morphisms $p_X$ and $\pi_X$ are affine
(\hyperref[II.8.3.5]{8.3.5} and
\hyperref[II.8.3.6]{8.3.6}), it suffices (in view of
(\hyperref[II.1.3.1]{1.3.1})) to verify the corollary when
$Y=\operatorname{Spec}(A)$ is affine and $\mathcal{S}=\widetilde{S}$. The
first assertion follows from the existence of the canonical isomorphism
(\hyperref[II.8.2.7.2]{8.2.7.2})
$(S_f^{\geq})_{f/1}^{\leq}\xrightarrow{\sim}S_f^{\leq}$ and from the fact that these
isomorphisms are compatible with the passage from $f$ to $fg$ ($f$ and $g$
homogeneous in $S_+$). Similarly, the formula
(\hyperref[II.8.3.6.2]{8.3.6.2}), applied to $\pi_X$ instead of $\pi$,
shows that
$\pi_X^{-1}(u^{-1}(X_f))=
\operatorname{Spec}((S_f^{\geq})_{f/1})$ for $f$ homogeneous in $S_+$, and the
second assertion follows from the existence of the canonical isomorphism
(\hyperref[II.8.2.7.2]{8.2.7.2})
$(S_f^{\geq})_{f/1}\xrightarrow{\sim}S_f$.

We can thus say that $\widehat{C}_X$, considered as an $X$-scheme, is obtained
by \emph{gluing} the two $X$-schemes affine over $X$,
$C_X=\operatorname{Spec}(\mathcal{S}_X^{\geq})$ and
$\widehat{E}_X=\operatorname{Spec}(\mathcal{S}_X^{\leq})$, whose
intersection is the open subset $E_X=\operatorname{Spec}(\mathcal{S}_X)$.
\end{proof}

\begin{corollary}[8.6.4]
\label{II.8.6.4}
Suppose that $\mathcal{O}_X(1)$ is an invertible $\mathcal{O}_X$-module and
that $\mathcal{S}_X$ is isomorphic to
$\bigoplus_{n\in\mathbf{Z}}(\mathcal{O}_X(1))^{\otimes n}$ (which will in
particular be the case if $\mathcal{S}$ is generated by $\mathcal{S}_1$
(\hyperref[II.3.2.5]{3.2.5} and
\hyperref[II.3.2.7]{3.2.7})). Then the punctured projective cone
$\widehat{E}$ is identified with the vector bundle of rank one
$\mathbf{V}(\mathcal{O}_X(-1))$ over $X$, and the punctured affine cone $E$ with the
sub-prescheme of this vector bundle induced on the complement of the
zero section. With this identification, the canonical retraction
$\widehat{E}\to X$ is identified with the structural morphism of the $X$-scheme
$\mathbf{V}(\mathcal{O}_X(-1))$. Finally, there exists a canonical $Y$-morphism
$\mathbf{V}(\mathcal{O}_X(1))\to C$ whose restriction to the complement of
the zero section of $\mathbf{V}(\mathcal{O}_X(1))$ is an isomorphism of this
complement onto the punctured affine cone $E$.
\end{corollary}

\begin{proof}
Indeed, if we set $\mathcal{L}=\mathcal{O}_X(1)$,
$\mathcal{S}_X^{\geq}$ is then identical to
$\mathbf{S}_{\mathcal{O}_X}(\mathcal{L})$, so $\widehat{E}_X$ is canonically
identified with $\mathbf{V}(\mathcal{L}^{-1})$ by virtue of
(\hyperref[II.8.6.3]{8.6.3}), and $C_X$ with $\mathbf{V}(\mathcal{L})$. The
morphism $\mathbf{V}(\mathcal{L})\to C$ is the restriction of
$\operatorname{Proj}(\widehat{\alpha})$, and the assertions of the corollary are
therefore particular cases of (\hyperref[II.8.6.2]{8.6.2}).
\end{proof}

Note that the inverse image under the morphism
$\mathbf{V}(\mathcal{O}_X(1))\to C$ of the underlying space of the apex
prescheme of $C$ is the underlying space of the zero section of
$\mathbf{V}(\mathcal{O}_X(1))$ (\hyperref[II.8.5.5]{8.5.5}); but in
general the corresponding sub-preschemes of $C$ and of
$\mathbf{V}(\mathcal{O}_X(1))$ are not isomorphic. This question will be
studied below.

\subsection{Blow-up of projecting cones}
\label{subsection:II.8.7}

\begin{env}[8.7.1]
\label{II.8.7.1}
Under the conditions of (\hyperref[II.8.6.1]{8.6.1}), setting
$r=\operatorname{Proj}(\widehat{\alpha})$, we have a commutative diagram
\begin{equation}
\label{II.8.7.1.1}
  \xymatrix{
    X\ar[r]^{i_X\circ\varepsilon_X}\ar[d]_{q} &
    \widehat{C}_X\ar[d]^{r}\\
    Y\ar[r]_{i\circ\varepsilon} & \widehat{C}
  }
\tag{8.7.1.1}
\end{equation}

\oldpage[II]{174}
by virtue of (\hyperref[II.8.5.1.3]{8.5.1.3}) and
(\hyperref[II.8.5.3.2]{8.5.3.2}); moreover, the restriction of $r$ to the
complement $\widehat{C}_X-i_X(\varepsilon_X(X))$ of the zero section is an
\emph{isomorphism} onto the complement $\widehat{C}-i(\varepsilon(Y))$ of the
zero section, by virtue of (\hyperref[II.8.6.2]{8.6.2}). If we suppose for
simplicity that $Y$ is affine and that $\mathcal{S}$ is of finite type and
generated by $\mathcal{S}_1$, then $X$ is projective over $Y$ and
$\widehat{C}_X$ is projective over $X$ (\hyperref[II.5.5.1]{5.5.1}), hence
$\widehat{C}_X$ is projective over $Y$ (\hyperref[II.5.5.5]{5.5.5, (ii)}),
and \emph{a fortiori} over $\widehat{C}$ (\hyperref[II.5.5.5]{5.5.5, (v)}).
We thus have a projective $Y$-morphism $r:\widehat{C}_X\to\widehat{C}$ (whose
restriction to $C_X$ is a projective $Y$-morphism $C_X\to C$) which
\emph{contracts $X$ into $Y$} and induces an \emph{isomorphism} when it is
restricted to the \emph{complements of $X$ and of $Y$}. We therefore have a
relation between $C_X$ and $C$ analogous to the one that holds between a
blow-up prescheme and the prescheme it comes from
(\hyperref[II.8.1.3]{8.1.3}). We will indeed show that $C_X$ can be
identified with the homogeneous spectrum of a graded $\mathcal{O}_C$-algebra.
\end{env}

\begin{env}[8.7.2]
\label{II.8.7.2}
Keeping the notation of (\hyperref[II.8.6.1]{8.6.1}), consider, for each
$n\geqslant 0$, the quasi-coherent ideal
\begin{equation}
\label{II.8.7.2.1}
  \mathcal{S}_{[n]}=\bigoplus_{m\geqslant n}\mathcal{S}_m
\tag{8.7.2.1}
\end{equation}
of the graded $\mathcal{O}_Y$-algebra $\mathcal{S}$. It is clear that we have
\begin{equation}
\label{II.8.7.2.2}
  \mathcal{S}_{[0]}=\mathcal{S},\qquad
  \mathcal{S}_{[n]}\subset\mathcal{S}_{[m]}
  \qquad\text{for }m\leqslant n
\tag{8.7.2.2}
\end{equation}
\begin{equation}
\label{II.8.7.2.3}
  \mathcal{S}_{[n]}\mathcal{S}_{[m]}\subset\mathcal{S}_{[m+n]}.
\tag{8.7.2.3}
\end{equation}

Consider the $\mathcal{O}_C$-module associated to $\mathcal{S}_{[n]}$, which is
a quasi-coherent ideal of $\mathcal{O}_C=\widetilde{\mathcal{S}}$
(\hyperref[II.1.4.4]{1.4.4})
\begin{equation}
\label{II.8.7.2.4}
  \mathcal{I}_n=(\mathcal{S}_{[n]})^{\sim}.
\tag{8.7.2.4}
\end{equation}

We then deduce from (\hyperref[II.8.7.2.2]{8.7.2.2}) and
(\hyperref[II.8.7.2.3]{8.7.2.3}), using
(\hyperref[II.1.4.4]{1.4.4}) and (\hyperref[II.1.4.8.1]{1.4.8.1}),
the analogous formulas
\begin{equation}
\label{II.8.7.2.5}
  \mathcal{I}_0=\mathcal{O}_C,\qquad
  \mathcal{I}_n\subset\mathcal{I}_m
  \qquad\text{for }m\leqslant n
\tag{8.7.2.5}
\end{equation}
\begin{equation}
\label{II.8.7.2.6}
  \mathcal{I}_n\mathcal{I}_m\subset\mathcal{I}_{n+m}.
\tag{8.7.2.6}
\end{equation}

We are therefore under the conditions of (\hyperref[II.8.1.1]{8.1.1}),
which leads us to introduce the quasi-coherent graded $\mathcal{O}_C$-algebra
\begin{equation}
\label{II.8.7.2.7}
  \mathcal{S}^{\natural}
  =\bigoplus_{n\geqslant 0}\mathcal{I}_n
  =\left(\bigoplus_{n\geqslant 0}\mathcal{S}_{[n]}\right)^{\sim}.
\tag{8.7.2.7}
\end{equation}
\end{env}

\begin{proposition}[8.7.3]
\label{II.8.7.3}
There exists a canonical $C$-isomorphism
\begin{equation}
\label{II.8.7.3.1}
  h:C_X\xrightarrow{\sim}\operatorname{Proj}(\mathcal{S}^{\natural}).
\tag{8.7.3.1}
\end{equation}
\end{proposition}

\begin{proof}
First suppose $Y=\operatorname{Spec}(A)$ affine, whence
$\mathcal{S}=\widetilde{S}$, where $S$ is an $A$-algebra graded in positive
degrees, and $C=\operatorname{Spec}(S)$. The definition
(\hyperref[II.8.2.6.2]{8.2.6.2}) then shows that, with the notation of
(\hyperref[II.8.2.6]{8.2.6}), we have
$\mathcal{S}^{\natural}=(S^{\natural})^{\sim}$. To define
(\hyperref[II.8.7.3.1]{8.7.3.1}), consider a homogeneous element
$f\in S_d$ ($d>0$) and the corresponding element $f^{\natural}\in S^{\natural}$
(\hyperref[II.8.2.6]{8.2.6}); the $S$-isomorphism
(\hyperref[II.8.2.7.3]{8.2.7.3}) therefore defines a $C$-isomorphism
\begin{equation}
\label{II.8.7.3.2}
  \operatorname{Spec}(S_f^{\geqslant})
  \xrightarrow{\sim}
  \operatorname{Spec}(S_{(f^{\natural})}^{\natural}).
\tag{8.7.3.2}
\end{equation}

\oldpage[II]{175}
But with the notation of (\hyperref[II.8.6.2]{8.6.2}), if
$v:C_X\to X$ is the structural morphism, it follows from
(\hyperref[II.8.6.2.1]{8.6.2.1}) that
$v^{-1}(X_f)=\operatorname{Spec}(S_f^{\geqslant})$. On the other hand,
$\operatorname{Spec}(S_{(f^{\natural})}^{\natural})=D_+(f^{\natural})$, so
that (\hyperref[II.8.7.3.2]{8.7.3.2}) defines an isomorphism
$v^{-1}(X_f)\to D_+(f^{\natural})$. Moreover, if $g\in S_e$ ($e>0$), the
diagram
\[
  \xymatrix{
    v^{-1}(X_{fg})\ar[r]^{\sim}\ar[d] &
    D_+(f^{\natural}g^{\natural})\ar[d]\\
    v^{-1}(X_f)\ar[r]^{\sim} & D_+(f^{\natural})
  }
\]
is commutative, as follows immediately from the definition of the isomorphism
(\hyperref[II.8.2.7.3]{8.2.7.3}). Finally, by definition $S_+$ is generated
by the homogeneous $f$, so it follows from
(\hyperref[II.8.2.10]{8.2.10, (iv)}) and from
(\hyperref[II.2.3.14]{2.3.14}) that the $D_+(f^{\natural})$ form a
covering of $\operatorname{Proj}(S^{\natural})$ and the $v^{-1}(X_f)$ a
covering of $C_X$, since the $X_f$ form a covering of $X$; this therefore
completes the definition of the isomorphism
(\hyperref[II.8.7.3.1]{8.7.3.1}) in this case.

To prove (\hyperref[II.8.7.3]{8.7.3}) in the general case, it suffices
to see that if $U$, $U'$ are two affine open subsets of $Y$ such that
$U'\subset U$, with rings $A$ and $A'$, and if we set
$\mathcal{S}|U=\widetilde{S}$, $\mathcal{S}|U'=\widetilde{S'}$, then the
diagram
\begin{equation}
\label{II.8.7.3.3}
  \xymatrix{
    C_{U'}\ar[r]\ar[d] & \operatorname{Proj}(S'^{\natural})\ar[d]\\
    C_U\ar[r] & \operatorname{Proj}(S^{\natural})
  }
\tag{8.7.3.3}
\end{equation}
is commutative. But $S'$ is canonically identified with $S\otimes_A A'$, hence
$S'^{\natural}$ with
\[
  S^{\natural}\otimes_S S'=S^{\natural}\otimes_A A'~;
\]
we therefore have
$\operatorname{Proj}(S'^{\natural})=\operatorname{Proj}(S^{\natural})
\times_U U'$ (\hyperref[II.2.8.10]{2.8.10}); likewise, if
$X=\operatorname{Proj}(S)$, $X'=\operatorname{Proj}(S')$, we have
$X'=X\times_U U'$ and
$\mathcal{S}_{X'}=\mathcal{S}_X\otimes_{\mathcal{O}_U}\mathcal{O}_{U'}$
(\hyperref[II.3.5.4]{3.5.4}), in other words
$\mathcal{S}_{X'}=j^*(\mathcal{S}_X)$, where $j$ is the projection $X'\to X$.
Consequently (\hyperref[II.1.5.2]{1.5.2}) we have
$C_{U'}=C_U\times_X X'=C_U\times_U U'$, and the commutativity of
(\hyperref[II.8.7.3.3]{8.7.3.3}) is then immediate.
\end{proof}

\begin{remark}[8.7.4]
\label{II.8.7.4}
\begin{enumerate}
  \item[\rm{(i)}] The end of the argument of
  (\hyperref[II.8.7.3]{8.7.3}) generalizes immediately as
  follows. Let $g:Y'\to Y$ be a morphism,
  $\mathcal{S}'=g^*(\mathcal{S})$, $X'=\operatorname{Proj}(\mathcal{S}')$;
  we then have a commutative diagram
  \begin{equation}
  \label{II.8.7.4.1}
    \xymatrix{
      C_{X'}\ar[r]\ar[d] &
      \operatorname{Proj}(\mathcal{S}'^{\natural})\ar[d]\\
      C_X\ar[r] & \operatorname{Proj}(\mathcal{S}^{\natural})
    }
  \tag{8.7.4.1}
  \end{equation}

  On the other hand, let $\varphi:\mathcal{S}''\to\mathcal{S}$ be a
  homomorphism of graded $\mathcal{O}_Y$-algebras such that, if we set
  $X''=\operatorname{Proj}(\mathcal{S}'')$,
  $u=\operatorname{Proj}(\varphi):X\to X''$ is everywhere defined; we also
  have

  \oldpage[II]{176}
  a $Y$-morphism $v:C\to C''$ (with
  $C''=\operatorname{Spec}(\mathcal{S}'')$) such that
  $\mathcal{A}(v)=\varphi$, and since $\varphi$ is a homomorphism of graded
  algebras, we deduce from $\varphi$ a $v$-morphism of graded algebras
  $\psi:\mathcal{S}''^{\natural}\to\mathcal{S}^{\natural}$
  (\hyperref[II.1.4.1]{1.4.1}). Moreover, it follows easily from
  (\hyperref[II.8.2.10]{8.2.10, (iv)}) and from the hypothesis on $\varphi$
  that $\operatorname{Proj}(\psi)$ is everywhere defined. Finally, taking into
  account (\hyperref[II.3.5.6.1]{3.5.6.1}), we have a canonical
  $u$-morphism $\mathcal{S}_{X''}\to\mathcal{S}_X$, whence
  (\hyperref[II.1.5.6]{1.5.6}) a morphism $w:C_X\to C_{X''}$. This being
  so, the diagram
  \begin{equation}
  \label{II.8.7.4.2}
    \xymatrix{
      C_X\ar[r]^{\sim}\ar[d]_{w} &
      \operatorname{Proj}(\mathcal{S}^{\natural})
        \ar[d]^{\operatorname{Proj}(\psi)}\\
      C_{X''}\ar[r]^{\sim} & \operatorname{Proj}(\mathcal{S}''^{\natural})
    }
  \tag{8.7.4.2}
  \end{equation}
  is commutative, as is immediately verified by reducing to the case where
  $Y$ is affine.

  \item[\rm{(ii)}] Note that by virtue of
  (\hyperref[II.8.7.2.5]{8.7.2.5}) and
  (\hyperref[II.8.7.2.6]{8.7.2.6}), we have
  $\mathcal{I}_1^m\subset\mathcal{I}_m\subset\mathcal{I}_1$ for every $m>0$.
  Now, by definition, $\mathcal{I}_1=(\mathcal{S}_+)^{\sim}$, so
  $\mathcal{I}_1$ defines in $C$ the closed sub-prescheme
  $\varepsilon(Y)$ ((\hyperref[II.1.4.10]{1.4.10}) and
  (\hyperref[II.8.3.2]{8.3.2})); we conclude that for every $m>0$,
  \emph{the support of $\mathcal{O}_C/\mathcal{I}_m$ is contained in the
  underlying space of the apex prescheme $\varepsilon(Y)$}; in the inverse
  image of the affine punctured cone $E$, the structural morphism
  $\operatorname{Proj}(\mathcal{S}^{\natural})\to C$ therefore reduces to an
  \emph{isomorphism} (as also follows from
  (\hyperref[II.8.7.3]{8.7.3}) and
  (\hyperref[II.8.7.1]{8.7.1})). Furthermore, canonically identifying $C$
  with an open subset of $\widehat{C}$ (\hyperref[II.8.3.3]{8.3.3}), we can
  evidently extend the ideals $\mathcal{I}_m$ of $\mathcal{O}_C$ to
  ideals $\mathcal{J}_m$ of $\mathcal{O}_{\widehat{C}}$, by the condition that
  they coincide with $\mathcal{O}_{\widehat{C}}$ on the open subset
  $\widehat{E}$ of $\widehat{C}$. If we set
  $\mathcal{T}=\bigoplus_{n\geqslant 0}\mathcal{J}_n$, which is a
  quasi-coherent graded $\mathcal{O}_{\widehat{C}}$-algebra, we can
  extend the isomorphism (\hyperref[II.8.7.3.1]{8.7.3.1}) to a
  $\widehat{C}$-isomorphism
  \begin{equation}
  \label{II.8.7.4.3}
    \widehat{C}_X\xrightarrow{\sim}\operatorname{Proj}(\mathcal{T}).
  \tag{8.7.4.3}
  \end{equation}

  Indeed, over $\widehat{E}$, it follows from the above that
  $\operatorname{Proj}(\mathcal{T})$ is canonically identified with
  $\widehat{E}$, and we therefore define the isomorphism
  (\hyperref[II.8.7.4.3]{8.7.4.3}) over $\widehat{E}$ by requiring
  it to coincide with the canonical isomorphism
  $\widehat{E}_X\to\widehat{E}$ (\hyperref[II.8.6.2]{8.6.2}); it is clear
  that this isomorphism and (\hyperref[II.8.7.3.1]{8.7.3.1}) then coincide
  over $E$.
\end{enumerate}
\end{remark}

\begin{corollary}[8.7.5]
\label{II.8.7.5}
Suppose that there exists $n_0>0$ such that
\begin{equation}
\label{II.8.7.5.1}
  \mathcal{S}_{n+1}=\mathcal{S}_1\mathcal{S}_n
  \qquad\text{for }n\geqslant n_0.
\tag{8.7.5.1}
\end{equation}
Then the apex sub-prescheme of $C_X$ (isomorphic to $X$) is the inverse
image under the canonical morphism $r:C_X\to C$ of the apex sub-prescheme
of $C$ (isomorphic to $Y$). Conversely, if this property holds and if we
suppose moreover that $Y$ is Noetherian and $\mathcal{S}$ is of finite type,
there exists $n_0>0$ such that (\hyperref[II.8.7.5.1]{8.7.5.1}) holds.
\end{corollary}

\begin{proof}
The first assertion being local on $Y$, to establish it we may suppose
$Y=\operatorname{Spec}(A)$ affine and $\mathcal{S}=\widetilde{S}$, where $S$ is
an $A$-algebra graded in positive degrees. It then follows from
(\hyperref[II.8.2.12]{8.2.12}), since we have
$\operatorname{Proj}(S^{\natural}\otimes_S S_0)=
C_X\times_C\varepsilon(Y)$ (by virtue of the identification
(\hyperref[II.8.7.3.1]{8.7.3.1})), that is, this prescheme is the
inverse image of $\varepsilon(Y)$ in $C_X$
(\hyperref[I.4.4.1]{I, 4.4.1}). The converse also follows from
(\hyperref[II.8.2.12]{8.2.12}) when $Y$ is affine Noetherian and $S$ is of
finite type.

\oldpage[II]{177}
If $Y$ is Noetherian (not necessarily affine) and $\mathcal{S}$ is of finite
type, there is a finite covering of $Y$ by Noetherian affine open subsets
$U_i$, and we then deduce from the above that for every $i$ there is an
integer $n_i$ such that
$\mathcal{S}_{n+1}|U_i=(\mathcal{S}_1|U_i)(\mathcal{S}_n|U_i)$ as soon as
$n\geqslant n_i$; the largest of the $n_i$ therefore satisfies
(\hyperref[II.8.7.5.1]{8.7.5.1}).
\end{proof}

\begin{env}[8.7.6]
\label{II.8.7.6}
Now consider the $C$-prescheme $Z$ obtained by \emph{blowing up}, in the
affine cone $C$, the \emph{apex sub-prescheme} $\varepsilon(Y)$; by
definition (\hyperref[II.8.1.3]{8.1.3}) it is therefore the prescheme
$\operatorname{Proj}(\bigoplus_{n\geqslant 0}\mathcal{S}_+^n)$; the canonical
injection
\begin{equation}
\label{II.8.7.6.1}
  \iota:\bigoplus_{n\geqslant 0}\mathcal{S}_+^n
  \longrightarrow\mathcal{S}^{\natural}
\tag{8.7.6.1}
\end{equation}
defines (thanks to the identification (\hyperref[II.8.7.3]{8.7.3})) a
canonical dominant $C$-morphism
\begin{equation}
\label{II.8.7.6.2}
  G(\iota)\longrightarrow Z
\tag{8.7.6.2}
\end{equation}
where $G(\iota)$ is an open subset of $C_X$ (\hyperref[II.3.5.1]{3.5.1});
note that we can have $G(\iota)\ne C_X$, as is shown by the example where
$Y=\operatorname{Spec}(K)$, $K$ being a field, and $\mathcal{S}=\widetilde{S}$
with $S=K[\mathbf{y}]$, where $\mathbf{y}$ is an indeterminate \emph{of degree
2}; if $R_n$ denotes the set $(S_+)^n$, considered as a subset of
$S_{[n]}=S_n^{\natural}$, then $S_+^{\natural}$ is not the radical in
$S^{\natural}$ of the ideal generated by the union of the $R_n$ (cf.
(\hyperref[II.2.3.14]{2.3.14})).
\end{env}

\begin{corollary}[8.7.7]
\label{II.8.7.7}
Suppose that there exists $n_0>0$ such that
\begin{equation}
\label{II.8.7.7.1}
  \mathcal{S}_n=\mathcal{S}_1^n
  \qquad\text{for }n\geqslant n_0.
\tag{8.7.7.1}
\end{equation}
Then the canonical morphism (\hyperref[II.8.7.6.2]{8.7.6.2}) is everywhere
defined and is an isomorphism $C_X\xrightarrow{\sim}Z$. Conversely, if this
property holds and if we suppose moreover that $Y$ is Noetherian and
$\mathcal{S}$ is of finite type, there exists $n_0$ such that
(\hyperref[II.8.7.7.1]{8.7.7.1}) holds.
\end{corollary}

\begin{proof}
Indeed, the first assertion is local on $Y$, and therefore follows from
(\hyperref[II.8.2.14]{8.2.14}); the same is true of the converse,
reasoning as in (\hyperref[II.8.7.5]{8.7.5}).
\end{proof}

\begin{remark}[8.7.8]
\label{II.8.7.8}
Since condition (\hyperref[II.8.7.7.1]{8.7.7.1}) implies
(\hyperref[II.8.7.5.1]{8.7.5.1}), we see that when it is satisfied, not
only is $C_X$ identified with the prescheme obtained by blowing up the apex
(identified with $Y$) of the affine cone $C$, but also the apex (identified
with $X$) of $C_X$ is identified with the closed sub-prescheme that is the
inverse image of the apex $Y$ of $C$. Moreover, hypothesis
(\hyperref[II.8.7.7.1]{8.7.7.1}) implies that on
$X=\operatorname{Proj}(\mathcal{S})$ the $\mathcal{O}_X$-modules
$\mathcal{O}_X(n)$ are invertible (\hyperref[II.3.2.5]{3.2.5} and
\hyperref[II.3.2.9]{3.2.9}) and that
$\mathcal{O}_X(n)=\mathcal{L}^{\otimes n}$ with
$\mathcal{L}=\mathcal{O}_X(1)$ (\hyperref[II.3.2.7]{3.2.7} and
\hyperref[II.3.2.9]{3.2.9}); by definition
(\hyperref[II.8.6.1.1]{8.6.1.1}), $C_X$ is therefore the \emph{vector
bundle} $\mathbf{V}(\mathcal{L})$ over $X$, and its apex is the \emph{zero
section} of this vector bundle.
\end{remark}

\subsection{Ample sheaves and contractions}
\label{subsection:II.8.8}

\begin{env}[8.8.1]
\label{II.8.8.1}
Let $Y$ be a prescheme, $f:X\to Y$ a \emph{separated and quasi-compact}
morphism, and $\mathcal{L}$ an invertible $\mathcal{O}_X$-module \emph{ample
relative to $f$}. Consider the $\mathcal{O}_Y$-algebra graded in positive
degrees
\begin{equation}
\label{II.8.8.1.1}
  \mathcal{S}=\mathcal{O}_Y\oplus
  \bigoplus_{n\geqslant 1}f_*(\mathcal{L}^{\otimes n})
\tag{8.8.1.1}
\end{equation}

\oldpage[II]{178}
which is quasi-coherent (\hyperref[I.9.2.2]{I, 9.2.2, a)}). We have a
canonical homomorphism of graded $\mathcal{O}_X$-algebras
\begin{equation}
\label{II.8.8.1.2}
  \tau:f^*(\mathcal{S})\longrightarrow
  \bigoplus_{n\geqslant 0}\mathcal{L}^{\otimes n}
\tag{8.8.1.2}
\end{equation}
which, in degrees $\geqslant 1$, coincides with the canonical homomorphism
$\sigma:f^*(f_*(\mathcal{L}^{\otimes n}))\to\mathcal{L}^{\otimes n}$
(\hyperref[0.4.4.3]{0, 4.4.3}), and in degree $0$ is the identity.
The hypothesis that $\mathcal{L}$ is $f$-ample then implies
(\hyperref[II.4.6.3]{4.6.3} and \hyperref[II.3.6.1]{3.6.1}) that the
corresponding $Y$-morphism
\begin{equation}
\label{II.8.8.1.3}
  r=r_{\mathcal{L},\tau}:X\longrightarrow
  P=\operatorname{Proj}(\mathcal{S})
\tag{8.8.1.3}
\end{equation}
is everywhere defined and is a \emph{dominant open immersion}, and that we
have
\begin{equation}
\label{II.8.8.1.4}
  r^*(\mathcal{O}_P(n))=\mathcal{L}^{\otimes n}
  \qquad\text{for every }n\in\mathbf{Z}.
\tag{8.8.1.4}
\end{equation}
\end{env}

\begin{proposition}[8.8.2]
\label{II.8.8.2}
Let $C=\operatorname{Spec}(\mathcal{S})$ be \emph{the affine cone defined by
$\mathcal{S}$}; if $\mathcal{L}$ is $f$-ample, there exists a canonical
$Y$-morphism
\begin{equation}
\label{II.8.8.2.1}
  g:V=\mathbf{V}(\mathcal{L})\longrightarrow C
\tag{8.8.2.1}
\end{equation}
such that the diagram
\begin{equation}
\label{II.8.8.2.2}
  \xymatrix{
    X\ar[r]^{j}\ar[d]_{f} &
    \mathbf{V}(\mathcal{L})\ar[r]^{\pi}\ar[d]^{g} & X\ar[d]^{f}\\
    Y\ar[r]_{\varepsilon} & C\ar[r]_{\psi} & Y
  }
\tag{8.8.2.2}
\end{equation}
is commutative, $\psi$ and $\pi$ being the structural morphisms, and $j$ and
$\varepsilon$ the canonical immersions mapping $X$ and $Y$ respectively onto
the zero section of $\mathbf{V}(\mathcal{L})$ and the apex prescheme
of $C$. Moreover, the restriction of $g$ to
$\mathbf{V}(\mathcal{L})-j(X)$ is an open immersion
\begin{equation}
\label{II.8.8.2.3}
  \mathbf{V}(\mathcal{L})-j(X)\longrightarrow
  E=C-\varepsilon(Y)
\tag{8.8.2.3}
\end{equation}
into the affine punctured cone $E$ corresponding to $\mathcal{S}$.
\end{proposition}

\begin{proof}
With the notation of (\hyperref[II.8.8.1]{8.8.1}), let
$\mathcal{S}_P^{\geqslant}=\bigoplus_{n\geqslant 0}\mathcal{O}_P(n)$ and
$C_P=\operatorname{Spec}(\mathcal{S}_P^{\geqslant})$. We know
(\hyperref[II.8.6.2]{8.6.2}) that we have a canonical morphism
$h=\operatorname{Spec}(\alpha):C_P\to C$ such that the diagram
\begin{equation}
\label{II.8.8.2.4}
  \xymatrix{
    C_P\ar[r]\ar[d]_{h} & P\ar[d]^{p}\\
    C\ar[r]_{\psi} & Y
  }
\tag{8.8.2.4}
\end{equation}
is commutative; moreover, if $\varepsilon_P:P\to C_P$ is the canonical
immersion, the diagram
\begin{equation}
\label{II.8.8.2.5}
  \xymatrix{
    P\ar[r]^{\varepsilon_P}\ar[d]_{p} & C_P\ar[d]^{h}\\
    Y\ar[r]_{\varepsilon} & C
  }
\tag{8.8.2.5}
\end{equation}
is commutative (\hyperref[II.8.7.1.1]{8.7.1.1}), and finally, the
restriction of $h$ to the affine punctured cone $E_P$ is an
\emph{isomorphism} $E_P\xrightarrow{\sim}E$ (\hyperref[II.8.6.2]{8.6.2}).
On the other hand, it follows from (\hyperref[II.8.8.1.4]{8.8.1.4}) that
we have
\[
  r^*(\mathcal{S}_P^{\geqslant})=
  \mathbf{S}_{\mathcal{O}_X}(\mathcal{L})
\]

\oldpage[II]{179}
and consequently we have a canonical $P$-morphism
$q:\mathbf{V}(\mathcal{L})\to C_P$, the commutative diagram
\begin{equation}
\label{II.8.8.2.6}
  \xymatrix{
    \mathbf{V}(\mathcal{L})\ar[r]^{\pi}\ar[d]_{q} & X\ar[d]^{r}\\
    C_P\ar[r] & P
  }
\tag{8.8.2.6}
\end{equation}
identifying $\mathbf{V}(\mathcal{L})$ with the product $C_P\times_P X$
(\hyperref[II.1.5.2]{1.5.2}); since $r$ is an open immersion, so
therefore is $q$ (\hyperref[I.4.3.2]{I, 4.3.2}). Moreover, the
restriction of $q$ to $\mathbf{V}(\mathcal{L})-j(X)$ maps this prescheme into
$E_P$ by (\hyperref[II.8.5.2]{8.5.2}), and the diagram
\begin{equation}
\label{II.8.8.2.7}
  \xymatrix{
    X\ar[r]^{j}\ar[d]_{r} & \mathbf{V}(\mathcal{L})\ar[d]^{q}\\
    P\ar[r]_{\varepsilon_P} & C_P
  }
\tag{8.8.2.7}
\end{equation}
is commutative (a particular case of
(\hyperref[II.8.5.1.3]{8.5.1.3})). The assertions of
(\hyperref[II.8.8.2]{8.8.2}) follow immediately from these facts, by
taking for $g$ the composite morphism $h\circ q$.
\end{proof}

\begin{remark}[8.8.3]
\label{II.8.8.3}
Suppose moreover that $Y$ is a \emph{Noetherian} prescheme and $f$ a
\emph{proper} morphism. Since $r$ is then \emph{proper}
(\hyperref[II.5.4.4]{5.4.4}), hence closed, and is a dominant open
immersion, $r$ is necessarily an \emph{isomorphism}
$X\xrightarrow{\sim}P$. We will moreover see in Chapter III
(\hyperref[III.2.3.5.1]{III, 2.3.5.1}) that $\mathcal{S}$ is then
necessarily an $\mathcal{O}_Y$-algebra \emph{of finite type}. It then
follows that $\mathcal{S}^{\natural}$ is an
$\mathcal{S}_0^{\natural}$-algebra \emph{of finite type}
(\hyperref[II.8.2.10]{8.2.10, (i)} and
\hyperref[II.8.7.2.7]{8.7.2.7}); since $C_P$ is $C$-isomorphic to
$\operatorname{Proj}(\mathcal{S}^{\natural})$
(\hyperref[II.8.7.3]{8.7.3}), we see that the morphism $h:C_P\to C$ is
\emph{projective}; since the morphism $r$ is an isomorphism, so is
$q:\mathbf{V}(\mathcal{L})\to C_P$, and we conclude that the morphism
$g:\mathbf{V}(\mathcal{L})\to C$ is \emph{projective}. Moreover, since the
restriction of $h$ to $E_P$ is an isomorphism onto $E$ and $q$ is an
isomorphism, the restriction (\hyperref[II.8.8.2.3]{8.8.2.3}) of $g$ is an
\emph{isomorphism}
$\mathbf{V}(\mathcal{L})-j(X)\xrightarrow{\sim}E$.

If we suppose moreover that $\mathcal{L}$ is \emph{very ample} for $f$, we
will also see in Chapter III (\hyperref[III.2.3.5.1]{III, 2.3.5.1}) that
there exists an integer $n_0>0$ such that
$\mathcal{S}_n=\mathcal{S}_1^n$ for $n\geqslant n_0$. We will conclude from
this by (\hyperref[II.8.7.7]{8.7.7}) that $\mathbf{V}(\mathcal{L})$ is
identified with the prescheme $Z$ obtained by \emph{blowing up, in the affine
cone $C$, the apex sub-prescheme} (identified with $Y$), and that the
\emph{zero section} of $\mathbf{V}(\mathcal{L})$ (identified with $Y$) is
\emph{the inverse image} of the apex sub-prescheme $Y$ of $C$. Some of the
preceding results can in fact be established without a Noetherian
hypothesis:
\end{remark}

\begin{corollary}[8.8.4]
\label{II.8.8.4}
Let $Y$ be a prescheme (resp. a quasi-compact scheme), $f:X\to Y$ a proper
morphism, and $\mathcal{L}$ an invertible $\mathcal{O}_X$-module ample for
$f$. Then the morphism (\hyperref[II.8.8.2.1]{8.8.2.1}) is proper (resp.
projective) and its restriction (\hyperref[II.8.8.2.3]{8.8.2.3}) is an
isomorphism.
\end{corollary}

\begin{proof}
To prove that $g$ is proper, we may reduce to the case where $Y$ is affine,
and it therefore suffices to consider the case where $Y$ is a quasi-compact
scheme. The same argument as in (\hyperref[II.8.8.3]{8.8.3}) first shows
that $r$ is an \emph{isomorphism} $X\xrightarrow{\sim}P$; consequently $q$ is
also an isomorphism, and since the restriction of $h$ to $E_P$ is an
isomorphism $E_P\xrightarrow{\sim}E$, we already see that
(\hyperref[II.8.8.2.3]{8.8.2.3}) is an isomorphism. It remains to prove
that $g$ is \emph{projective}.

Since $f$ is of finite type by hypothesis, we can apply
(\hyperref[II.3.8.5]{3.8.5}) to the homomorphism

\oldpage[II]{180}
$\tau$ of (\hyperref[II.8.8.1.2]{8.8.1.2}): consequently there exist an
integer $d>0$ and a quasi-coherent sub-$\mathcal{O}_Y$-module of finite type
$\mathcal{E}$ of $\mathcal{S}_d$ such that, if $\mathcal{S}'$ is the
sub-$\mathcal{O}_Y$-algebra of $\mathcal{S}$ generated by $\mathcal{E}$ and
$\tau'=\tau\circ f^*(\varphi)$ ($\varphi$ being the canonical injection
$\mathcal{S}'\to\mathcal{S}$), then $r'=r_{\mathcal{L},\tau'}$ is an immersion
\[
  X\longrightarrow P'=\operatorname{Proj}(\mathcal{S}').
\]
Moreover, since $\varphi$ is injective, $r'$ is also a \emph{dominant
immersion} (\hyperref[II.3.7.6]{3.7.6}); the same argument as for $r$
therefore shows that $r'$ is a \emph{surjective closed immersion}; since $r'$
factors as
$X\xrightarrow{r}\operatorname{Proj}(\mathcal{S})
\xrightarrow{\Phi}\operatorname{Proj}(\mathcal{S}')$, where
$\Phi=\operatorname{Proj}(\varphi)$, we conclude that $\Phi$ is also a
\emph{surjective closed immersion}. But this implies that $\Phi$ is an
\emph{isomorphism}; indeed, we may reduce to the case where
$Y=\operatorname{Spec}(A)$ is affine, $\mathcal{S}=\widetilde{S}$,
$\mathcal{S}'=\widetilde{S'}$, where $S$ is a graded $A$-algebra and $S'$ a
graded subalgebra of $S$. For every homogeneous element $t\in S'$,
$S'_{(t)}$ is then a subring of $S_{(t)}$; going back to the
definition of $\operatorname{Proj}(\varphi)$
(\hyperref[II.2.8.1]{2.8.1}), we see that we are reduced to proving that if
$B'$ is a subring of a ring $B$, and if the morphism
$\operatorname{Spec}(B)\to\operatorname{Spec}(B')$ corresponding to the
canonical injection $B'\to B$ is a closed immersion, then this morphism is
necessarily an \emph{isomorphism}; but this follows from
(\hyperref[I.4.2.3]{I, 4.2.3}).

Moreover, we have
$\Phi^*(\mathcal{O}_{P'}(n))=\mathcal{O}_P(n)$
(\hyperref[II.3.5.2]{3.5.2, (ii)} and
\hyperref[II.3.5.4]{3.5.4}), hence
$r'^*(\mathcal{O}_{P'}(n))$ is isomorphic to $\mathcal{L}^{\otimes n}$
(\hyperref[II.4.6.3]{4.6.3}). Set
$\mathcal{S}''=\mathcal{S}'^{(d)}$, so that
(\hyperref[II.3.1.8]{3.1.8, (i)}) $X$ is canonically identified with
$P''=\operatorname{Proj}(\mathcal{S}'')$ and
$\mathcal{L}''=\mathcal{L}^{\otimes d}$ with $\mathcal{O}_{P''}(1)$
(\hyperref[II.3.2.9]{3.2.9, (ii)}).

This being so, if $C''=\operatorname{Spec}(\mathcal{S}'')$, then
$\mathcal{S}_{P''}^{\geqslant}=
\bigoplus_{n\geqslant 0}\mathcal{O}_{P''}(n)$ is identified with
$\bigoplus_{n\geqslant 0}\mathcal{L}''^{\otimes n}$, hence
$C_{P''}=\operatorname{Spec}(\mathcal{S}_{P''}^{\geqslant})$ with
$\mathbf{V}(\mathcal{L}'')$; on the other hand, we know
(\hyperref[II.8.7.3]{8.7.3}) that $C_{P''}$ is $C''$-isomorphic to
$\operatorname{Proj}(\mathcal{S}''^{\natural})$; by definition of
$\mathcal{S}''$, $\mathcal{S}''^{\natural}$ is generated by
$\mathcal{S}_1''^{\natural}$, and $\mathcal{S}_1''^{\natural}$ is of finite type
over $\mathcal{S}_0''^{\natural}=\mathcal{S}''$
(\hyperref[II.8.2.10]{8.2.10, (i) and (iii)}), hence
$\operatorname{Proj}(\mathcal{S}''^{\natural})$ is \emph{projective} over
$C''$ (\hyperref[II.5.5.1]{5.5.1}). Now consider the diagram
\begin{equation}
\label{II.8.8.4.1}
  \xymatrix{
    \mathbf{V}(\mathcal{L})\ar[r]^{g}\ar[d]_{u} &
    \operatorname{Spec}(\mathcal{S})=C\ar[d]^{v}\\
    \mathbf{V}(\mathcal{L}'')\ar[r]_{g''} &
    \operatorname{Spec}(\mathcal{S}'')=C''
  }
\tag{8.8.4.1}
\end{equation}
where $g$ and $g''$ correspond by (\hyperref[II.1.5.6]{1.5.6}) to the
canonical $f$-morphisms
\[
  \mathcal{S}\longrightarrow
  \bigoplus_{n\geqslant 0}\mathcal{L}^{\otimes n}
  \quad\text{and}\quad
  \mathcal{S}''\longrightarrow
  \bigoplus_{n\geqslant 0}\mathcal{L}''^{\otimes n}
\]
(\hyperref[II.3.3.2.3]{3.3.2.3}) (see
(\hyperref[II.8.8.5]{8.8.5}) below), $v$ to the inclusion morphism
$\mathcal{S}''\to\mathcal{S}$, and $u$ to the inclusion morphism
$\bigoplus_{n\geqslant 0}\mathcal{L}^{\otimes nd}\to
\bigoplus_{n\geqslant 0}\mathcal{L}^{\otimes n}$; it is immediate
(\hyperref[II.3.3.2]{3.3.2}) that this diagram is commutative. We have
just seen that $g''$ is a projective morphism; on the other hand, $u$ is a
\emph{finite} morphism. Indeed, the question being local on $X$, we may
suppose that $X$ is affine with ring $A$ and that $\mathcal{L}=\mathcal{O}_X$;
everything then comes down to remarking that the ring $A[T]$ is a module of
finite type over its subring $A[T^d]$ ($T$ an indeterminate). Since $Y$ is a
quasi-compact scheme and $C''$ is affine over $Y$, $C''$ is also a
quasi-compact scheme,

\oldpage[II]{181}
and consequently $g''\circ u$ is a projective morphism
(\hyperref[II.5.5.5]{5.5.5, (ii)}); by the commutativity of
(\hyperref[II.8.8.4.1]{8.8.4.1}), the same holds for $v\circ g$, and since
$v$ is affine, hence separated, we finally conclude that $g$ is projective
(\hyperref[II.5.5.5]{5.5.5, (v)}).
\end{proof}

\begin{env}[8.8.5]
\label{II.8.8.5}
Let us return to the situation of (\hyperref[II.8.8.1]{8.8.1}). We will see
that the morphism $g:\mathbf{V}(\mathcal{L})\to C$ can be given another
definition, valid for every invertible $\mathcal{O}_X$-module $\mathcal{L}$
(not necessarily ample). For this, consider the $f$-morphism
\begin{equation}
\label{II.8.8.5.1}
  \tau^\flat:\mathcal{S}\longrightarrow
  \bigoplus_{n\geqslant 0}\mathcal{L}^{\otimes n}
\tag{8.8.5.1}
\end{equation}
corresponding to the morphism $\tau$ of
(\hyperref[II.8.8.1.2]{8.8.1.2}). From it we deduce
(\hyperref[II.1.5.6]{1.5.6}) a morphism $g':V\to C$ such that, if
$\pi:V\to X$ and $\psi:C\to Y$ are the structural morphisms, the diagrams
\begin{equation}
\label{II.8.8.5.2}
  \xymatrix{
    X\ar[d]_{f} & V\ar[l]^{\pi}\ar[d]^{g'}\\
    Y & C\ar[l]_{\psi}
  }
  \qquad
  \xymatrix{
    X\ar[r]^{j}\ar[d]_{f} & V\ar[d]^{g'}\\
    Y\ar[r]_{\varepsilon} & C
  }
\tag{8.8.5.2}
\end{equation}
are commutative
(\hyperref[II.8.5.1.2]{8.5.1.2} and
\hyperref[II.8.5.1.3]{8.5.1.3}). We will show that (if we suppose
$\mathcal{L}$ ample for $f$) the morphisms $g$ and $g'$ are identical.

Indeed, the question is local on $Y$, so we may suppose that
$Y=\operatorname{Spec}(A)$ is affine and (because of
(\hyperref[II.8.8.1.3]{8.8.1.3})) identify $X$ with an open subset of
$P=\operatorname{Proj}(S)$, where
\[
  S=A\oplus\bigoplus_{n\geqslant 1}
  \Gamma(X,\mathcal{L}^{\otimes n})~;
\]
we then deduce from (\hyperref[II.8.8.1.4]{8.8.1.4}) that
$\Gamma(X,\mathcal{O}_P(n))=\Gamma(X,\mathcal{L}^{\otimes n})$ for every
$n\in\mathbf{Z}$. Taking into account the definition of
$h=\operatorname{Spec}(\alpha)$, where $\alpha$ is the canonical $p$-morphism
$\widetilde{S}\to\mathcal{S}_P^{\geqslant}$
(\hyperref[II.8.6.1.2]{8.6.1.2}), we must verify that the restriction to
$X$ of
$\alpha^\sharp:p^*(\widetilde{S})\to\mathcal{S}_P^{\geqslant}$ is
identical to $\tau$. Taking into account
(\hyperref[0.4.4.3]{0, 4.4.3}), we are reduced to seeing that if we compose
the canonical homomorphism
$\alpha_n:S_n\to\Gamma(P,\mathcal{O}_P(n))$ with the restriction
homomorphism
\[
  \Gamma(P,\mathcal{O}_P(n))\longrightarrow
  \Gamma(X,\mathcal{O}_P(n))=
  \Gamma(X,\mathcal{L}^{\otimes n}),
\]
we obtain the identity for every $n>0$; now this follows immediately from the
definition of the algebra $S$ and from that of $\alpha_n$
(\hyperref[II.2.6.2]{2.6.2}).
\end{env}

\begin{proposition}[8.8.6]
\label{II.8.8.6}
Suppose (with the notation of (\hyperref[II.8.8.5]{8.8.5})) that, if we
write $f=(f_0,\lambda)$, the homomorphism
$\lambda:\mathcal{O}_Y\to f_*(\mathcal{O}_X)$ is bijective; then:
\begin{enumerate}
  \item[\rm{(i)}] If we write $g=(g_0,\mu)$, then
  $\mu:\mathcal{O}_C\to g_*(\mathcal{O}_V)$ is an isomorphism.
  \item[\rm{(ii)}] If $X$ is integral (resp. locally integral and normal),
  then $C$ is integral (resp. normal).
\end{enumerate}
\end{proposition}

\begin{proof}
Indeed, the $f$-morphism $\tau^\flat$ is then an \emph{isomorphism}
\[
  \tau^\flat:\mathcal{S}=\psi_*(\mathcal{O}_C)
  \longrightarrow f_*(\pi_*(\mathcal{O}_V))
  =\psi_*(g_*(\mathcal{O}_V))
\]
and the $Y$-morphism $g$ can be regarded as the one for which the
homomorphism $\mathcal{A}(g)$
(\hyperref[II.1.1.2]{1.1.2}) is equal to $\tau^\flat$. To see that
$\mu$ is an isomorphism of $\mathcal{O}_C$-modules, it suffices
(\hyperref[II.1.4.2]{1.4.2}) to see that
\[
  \mathcal{A}(\mu):\psi_*(\mathcal{O}_C)
  \longrightarrow\psi_*(g_*(\mathcal{O}_V))
\]
is an isomorphism. But by definition
(\hyperref[II.1.1.2]{1.1.2}), we have
$\mathcal{A}(\mu)=\mathcal{A}(g)$, whence the conclusion of (i).

To prove (ii), we may reduce to the case where $Y$ is affine, so that
$\mathcal{S}=\widetilde{S}$, with

\oldpage[II]{182}
$S=\bigoplus_{n\geqslant 0}\Gamma(X,\mathcal{L}^{\otimes n})$;
the hypothesis that $X$ is integral implies that the ring $S$ is integral
(\hyperref[I.7.4.4]{I, 7.4.4}), hence so is $C$
(\hyperref[I.5.1.4]{I, 5.1.4}). To prove that $C$ is normal, we
will use the following lemma:

\begin{lemma}[8.8.6.1]
\label{II.8.8.6.1}
Let $Z$ be a normal integral prescheme. Then the ring
$\Gamma(Z,\mathcal{O}_Z)$ is integral and integrally closed.
\end{lemma}

\begin{proof}
Indeed, it follows from
(\hyperref[I.8.2.1.1]{I, 8.2.1.1}) that
$\Gamma(Z,\mathcal{O}_Z)$ is the intersection, in the field of rational
functions $R(Z)$, of the integrally closed rings $\mathcal{O}_z$ for
$z\in Z$.
\end{proof}

This being so, let us first show that $V$ is \emph{locally integral and
normal}; for this we may restrict to the case where $X=\operatorname{Spec}(A)$
is affine, with $A$ an integral and integrally closed ring
(\hyperref[II.6.3.8]{6.3.8}), and where
$\mathcal{L}=\mathcal{O}_X$. Since then
$V=\operatorname{Spec}(A[T])$ and $A[T]$ is integral and integrally closed
([24], p.~99), this establishes our assertion. On the other hand, for every
affine open subset $U$ of $C$, $g^{-1}(U)$ is quasi-compact since the morphism
$g$ is quasi-compact; since $V$ is locally integral, the connected components
of $g^{-1}(U)$ are integral preschemes open in $g^{-1}(U)$, hence finite in
number, and since $V$ is normal, these preschemes are normal
(\hyperref[II.6.3.8]{6.3.8}). Then $\Gamma(U,\mathcal{O}_C)$, which is equal
to $\Gamma(g^{-1}(U),\mathcal{O}_V)$ by virtue of (i), is the direct product
of finitely many integral and integrally closed rings
(\hyperref[II.8.8.6.1]{8.8.6.1}), which shows that $C$ is normal
(\hyperref[II.6.3.4]{6.3.4}).
\end{proof}

\subsection{Grauert's ampleness criterion: statement}
\label{subsection:II.8.9}

We propose to show that the properties proved in
(\hyperref[II.8.8.2]{8.8.2}) \emph{characterize} the
$\mathcal{O}_X$-modules ample for $f$, and, more precisely, to prove the
following criterion:

\begin{theorem}[8.9.1]
\label{II.8.9.1}
\emph{(Grauert's criterion).} Let $Y$ be a prescheme,
$p:X\to Y$ a separated and quasi-compact morphism, and $\mathcal{L}$ an
invertible $\mathcal{O}_X$-module. For $\mathcal{L}$ to be ample
relative to $p$, it is necessary and sufficient that there exist a
$Y$-prescheme $C$, a $Y$-section $\varepsilon:Y\to C$ of $C$, and a
$Y$-morphism $q:\mathbf{V}(\mathcal{L})\to C$, having the following
properties:
\begin{enumerate}
  \item[\rm{(i)}] The diagram
  \begin{equation}
  \label{II.8.9.1.1}
    \xymatrix{
      X\ar[r]^{j}\ar[d]_{p} & \mathbf{V}(\mathcal{L})\ar[d]^{q}\\
      Y\ar[r]_{\varepsilon} & C
    }
  \tag{8.9.1.1}
  \end{equation}
  where $j$ is the zero section of the vector bundle $\mathbf{V}(\mathcal{L})$,
  is commutative.
  \item[\rm{(ii)}] The restriction of $q$ to
  $\mathbf{V}(\mathcal{L})-j(X)$ is a quasi-compact open immersion
  \[
    \mathbf{V}(\mathcal{L})-j(X)\longrightarrow C
  \]
  whose image does not meet $\varepsilon(Y)$.
\end{enumerate}
\end{theorem}

Note that if $C$ is \emph{separated} over $Y$, we can, in condition
(ii), drop the hypothesis that the open immersion is quasi-compact; indeed,
to see that this latter property is a consequence of the other conditions,
we may restrict to the case where $Y$ is affine, and the assertion then
follows from (\hyperref[I.5.5.1]{I, 5.5.1, (i)} and
\hyperref[I.5.5.10]{5.5.10}). We can also drop

\oldpage[II]{183}
the same hypothesis when $X$ is supposed Noetherian, for then $V$ is also
Noetherian and the assertion this time follows from
(\hyperref[I.6.3.5]{I, 6.3.5}).

\begin{corollary}[8.9.2]
\label{II.8.9.2}
If the morphism $p:X\to Y$ is proper, then in the statement of
(\hyperref[II.8.9.1]{8.9.1}) we may suppose $q$ proper, and replace
``open immersion'' by ``isomorphism''.
\end{corollary}

Figuratively, one can say (when $p:X\to Y$ is proper) that
$\mathcal{L}$ is ample relative to $p$ if and only if the zero section of the
vector bundle $\mathbf{V}(\mathcal{L})$ can be ``contracted'' to the base
prescheme $Y$. An important particular case is the one where $Y$ is the
spectrum of a field, and where the operation of ``contraction'' therefore
consists in contracting the zero section of $\mathbf{V}(\mathcal{L})$ \emph{to
a single point}.

\begin{env}[8.9.3]
\label{II.8.9.3}
The necessity of the conditions of (\hyperref[II.8.9.1]{8.9.1}) and
Corollary (\hyperref[II.8.9.2]{8.9.2}) follow immediately from
(\hyperref[II.8.8.2]{8.8.2}) and
(\hyperref[II.8.8.4]{8.8.4}).

To prove the sufficiency in (\hyperref[II.8.9.1]{8.9.1}), we
will consider a slightly more general situation. For this, set (with the
notation of (\hyperref[II.8.8.2]{8.8.2}))
\[
  \mathcal{S}'=\bigoplus_{n\geqslant 0}\mathcal{L}^{\otimes n}
\]
and
\[
  V=\mathbf{V}(\mathcal{L})=\operatorname{Spec}(\mathcal{S}').
\]
The closed sub-prescheme $j(X)$, the zero section of $\mathbf{V}(\mathcal{L})$,
is defined by the quasi-coherent sheaf of ideals
$\mathcal{J}=(\mathcal{S}'_+)^{\sim}$ of $\mathcal{O}_V$
(\hyperref[II.1.4.10]{1.4.10}). This $\mathcal{O}_V$-module is
\emph{invertible}, since the question is local on $X$ and it comes down to
remarking that the ideal $TA[T]$ in a polynomial ring $A[T]$ is a monogenic
free $A[T]$-module. Moreover, it is immediate (again because the
question is local on $X$) that
\[
  \mathcal{L}=j^*(\mathcal{J})
\]
and
\[
  j_*(\mathcal{L})=\mathcal{J}/\mathcal{J}^2.
\]
On the other hand, if
\[
  \pi:\mathbf{V}(\mathcal{L})\longrightarrow X
\]
is the structural morphism, we have
$\pi_*(\mathcal{J})=\mathcal{S}'_+$ and
$\pi_*(\mathcal{J}/\mathcal{J}^2)=\mathcal{L}$; we therefore have two
canonical homomorphisms
$\mathcal{L}\to\pi_*(\mathcal{J})\to\mathcal{L}$, the first being
the canonical injection $\mathcal{L}\to\mathcal{S}'_+$, the second the
canonical projection of $\mathcal{S}'_+$ onto $\mathcal{S}'_1=\mathcal{L}$,
and their composite is the identity. Moreover, we can canonically embed
\[
  \pi_*(\mathcal{J})=\mathcal{S}'_+
  =\bigoplus_{n\geqslant 1}\mathcal{L}^{\otimes n}
\]
in the \emph{product}
\[
  \prod_{n\geqslant 1}\mathcal{L}^{\otimes n}
  =\varprojlim_n\pi_*(\mathcal{J}/\mathcal{J}^{n+1})
\]
(since
$\pi_*(\mathcal{J}/\mathcal{J}^{n+1})=
\mathcal{L}\oplus\mathcal{L}^{\otimes2}\oplus\cdots\oplus
\mathcal{L}^{\otimes n}$), and we therefore again have two canonical
homomorphisms
\begin{equation}
\label{II.8.9.3.1}
  \mathcal{L}\longrightarrow
  \varprojlim_n\pi_*(\mathcal{J}/\mathcal{J}^{n+1})
  \longrightarrow\mathcal{L}
\tag{8.9.3.1}
\end{equation}
whose composite is the identity.

This being so, the generalization of (\hyperref[II.8.9.1]{8.9.1}) that we
will prove is the following:
\end{env}

\begin{proposition}[8.9.4]
\label{II.8.9.4}
Let $Y$ be a prescheme, $V$ a $Y$-prescheme, and $X$ a closed sub-prescheme of
$V$, defined by an ideal $\mathcal{J}$ of $\mathcal{O}_V$ which is an
invertible $\mathcal{O}_V$-module; if $j:X\to V$ is

\oldpage[II]{184}
the canonical injection, set
$\mathcal{L}=j^*(\mathcal{J})=
\mathcal{J}\otimes_{\mathcal{O}_V}\mathcal{O}_X$, so that
$j_*(\mathcal{L})=\mathcal{J}/\mathcal{J}^2$. Suppose that the structural
morphism $p:X\to Y$ is separated and quasi-compact, and that the following
conditions are satisfied:
\begin{enumerate}
  \item[\rm{(i)}] There exists a $Y$-morphism of finite type $\pi:V\to X$
  such that $\pi\circ j=1_X$, and consequently
  $\pi_*(\mathcal{J}/\mathcal{J}^2)=\mathcal{L}$.
  \item[\rm{(ii)}] There exists a homomorphism of $\mathcal{O}_X$-modules
  \[
    \varphi:\mathcal{L}\longrightarrow
    \varprojlim_n\pi_*(\mathcal{J}/\mathcal{J}^{n+1})
  \]
  such that the composite
  \[
    \mathcal{L}\xrightarrow{\varphi}
    \varprojlim_n\pi_*(\mathcal{J}/\mathcal{J}^{n+1})
    \xrightarrow{\alpha}\pi_*(\mathcal{J}/\mathcal{J}^2)=\mathcal{L}
  \]
  (where $\alpha$ is the canonical homomorphism) is the identity.
  \item[\rm{(iii)}] There exist a $Y$-prescheme $C$, a $Y$-section
  $\varepsilon$ of $C$ and a $Y$-morphism $q:V\to C$ such that the diagram
  \begin{equation}
  \label{II.8.9.4.1}
    \xymatrix{
      X\ar[r]^{j}\ar[d]_{p} & V\ar[d]^{q}\\
      Y\ar[r]_{\varepsilon} & C
    }
  \tag{8.9.4.1}
  \end{equation}
  is commutative.
  \item[\rm{(iv)}] The restriction of $q$ to $W=V-j(X)$ is a quasi-compact
  open immersion into $C$, whose image does not meet
  $\varepsilon(Y)$.
\end{enumerate}
Under these conditions, $\mathcal{L}$ is ample relative to $p$.
\end{proposition}


\subsection{Grauert's ampleness criterion: proof}
\label{subsection:II.8.10}

\begin{lemma}[8.10.1]
\label{II.8.10.1}
Let $\pi:V\to X$ be a morphism, $j:X\to V$ an $X$-section of $V$ which is
a closed immersion, and $\mathcal{J}$ the quasi-coherent ideal of
$\mathcal{O}_V$ defining the closed sub-prescheme of $V$ associated to $j$.
\begin{enumerate}
  \item[\rm{(i)}] For every $n\geqslant 0$,
  $\pi_*(\mathcal{O}_V/\mathcal{J}^{n+1})$ and
  $\pi_*(\mathcal{J}/\mathcal{J}^{n+1})$ are quasi-coherent
  $\mathcal{O}_X$-modules, and we have
  $\pi_*(\mathcal{O}_V/\mathcal{J})=\mathcal{O}_X$,
  $\pi_*(\mathcal{J}/\mathcal{J}^2)=j^*(\mathcal{J})$.
  \item[\rm{(ii)}] If $X=\{\xi\}=\operatorname{Spec}(k)$, where $k$ is a
  field, then $\varprojlim_n\pi_*(\mathcal{O}_V/\mathcal{J}^{n+1})$ is
  isomorphic to the separated completion of the local ring
  $\mathcal{O}_{j(\xi)}$ for the $\mathfrak{m}_{j(\xi)}$-preadic topology.
  \item[\rm{(iii)}] Suppose that $\mathcal{J}$ is an invertible
  $\mathcal{O}_V$-module (which implies that
  \[
    \mathcal{L}=j^*(\mathcal{J})=
    \pi_*(\mathcal{J}/\mathcal{J}^2)
  \]
  is an invertible $\mathcal{O}_X$-module) and that there exists a
  homomorphism
  \[
    \varphi:\mathcal{L}\longrightarrow
    \varprojlim_n\pi_*(\mathcal{J}/\mathcal{J}^{n+1})
  \]
  such that the composite
  \[
    \mathcal{L}\xrightarrow{\varphi}
    \varprojlim_n\pi_*(\mathcal{J}/\mathcal{J}^{n+1})
    \xrightarrow{\alpha}\pi_*(\mathcal{J}/\mathcal{J}^2)
  \]
  ($\alpha$ the canonical homomorphism) is the identity. If we set
  $\mathcal{S}=\bigoplus_{n\geqslant0}\mathcal{L}^{\otimes n}$, we then
  canonically deduce from $\varphi$ an isomorphism of
  $\mathcal{O}_X$-algebras from the completion $\widehat{\mathcal{S}}$ of
  $\mathcal{S}$ with respect to its canonical filtration (a completion which
  is isomorphic to the product $\prod_{n\geqslant0}\mathcal{L}^{\otimes n}$)
  onto $\varprojlim_n\pi_*(\mathcal{O}_V/\mathcal{J}^{n+1})$.
\end{enumerate}
\end{lemma}

\begin{proof}
Note first that the support of the $\mathcal{O}_V$-module
$\mathcal{O}_V/\mathcal{J}^{n+1}$ is $j(X)$, and that of
$\mathcal{J}/\mathcal{J}^{n+1}$ is contained in $j(X)$. In case (ii),
$j(X)$ is a closed point $j(\xi)$ of $V$,

\oldpage[II]{185}
and by definition $\pi_*(\mathcal{O}_V/\mathcal{J}^{n+1})$ is the stalk of
$\mathcal{O}_V/\mathcal{J}^{n+1}$ at the point $j(\xi)$, that is, setting
$C=\mathcal{O}_{j(\xi)}$ and denoting by $\mathfrak{m}$ the maximal ideal
of $C$, the $C$-module $C/\mathfrak{m}^{n+1}$; assertion (ii) is then
evident.

To prove (i), note that the question is local on $X$; we may therefore
restrict to the case where $X$ is affine. Let $U$ be an affine open subset of
$V$; $j(X)\cap U$ is an affine open subset of $j(X)$, hence
$U_0=\pi(j(X)\cap U)$, which is isomorphic to it, is an affine open subset of
$X$; for every affine open subset $W_0\subset U_0$ of $X$,
$W=\pi^{-1}(W_0)\cap U$ is an affine open subset of $V$, since $X$ is a
scheme (\hyperref[I.5.5.10]{I, 5.5.10}); in particular
$U'=U\cap\pi^{-1}(U_0)$ is an affine open subset of $V$, and we evidently have
$\pi(U')=U_0$ and $j(U_0)=j(X)\cap U$. This being so, by definition
\[
  \Gamma(W_0,\pi_*(\mathcal{O}_V/\mathcal{J}^{n+1}))
  =\Gamma(\pi^{-1}(W_0),\mathcal{O}_V/\mathcal{J}^{n+1})~;
\]
but since every point of $\pi^{-1}(W_0)$ not belonging to $j(W_0)$ has an
open neighbourhood in $\pi^{-1}(W_0)$ which does not meet $j(X)$, and in
which $\mathcal{O}_V/\mathcal{J}^{n+1}$ is therefore zero, it is clear that
the sections of $\mathcal{O}_V/\mathcal{J}^{n+1}$ over $\pi^{-1}(W_0)$ and
over $W$ are in bijective correspondence. In other words, if $\pi'$ is
the restriction of $\pi$ to $U'$, the $(\mathcal{O}_X|U_0)$-modules
\[
  \pi_*(\mathcal{O}_V/\mathcal{J}^{n+1})|U_0
  \quad\hbox{and}\quad
  \pi'_*\bigl((\mathcal{O}_V/\mathcal{J}^{n+1})|U'\bigr)
\]
are identical. Since $U'$ and $U_0$ are affine and the $U_0$ cover
$X$, we conclude (\hyperref[I.1.6.3]{I, 1.6.3}) that
$\pi_*(\mathcal{O}_V/\mathcal{J}^{n+1})$ is quasi-coherent, and the
proof is identical for
$\pi_*(\mathcal{J}/\mathcal{J}^{n+1})$.

Finally, to prove (iii), note that $\mathcal{S}$ is none other than
$\mathbf{S}_{\mathcal{O}_X}(\mathcal{L})$; we therefore canonically deduce
from $\varphi$ a homomorphism of $\mathcal{O}_X$-algebras
\[
  \psi:\mathcal{S}\longrightarrow
  \varprojlim_n\pi_*(\mathcal{O}_V/\mathcal{J}^{n+1})
\]
(\hyperref[II.1.7.4]{1.7.4}); moreover, this homomorphism maps
$\mathcal{L}^{\otimes n}$ into
$\varprojlim_m\pi_*(\mathcal{J}^n/\mathcal{J}^{m+1})$, hence is continuous for
the topologies under consideration, and consequently does extend to a
homomorphism
\[
  \widehat{\psi}:\widehat{\mathcal{S}}\longrightarrow
  \varprojlim_n\pi_*(\mathcal{O}_V/\mathcal{J}^{n+1}).
\]
To see that it is an isomorphism, we may, as in the proof of (i), reduce to
the case where $X=\operatorname{Spec}(A)$ and
$V=\operatorname{Spec}(B)$ are affine, with
$\mathcal{J}=\widetilde{\mathfrak{J}}$, where $\mathfrak{J}$ is an ideal of
$B$; to $\pi$ there corresponds an injection $A\to B$, identifying $A$ with a
subring of $B$ complementary to $\mathfrak{J}$, and $\mathcal{L}$ (resp.
$\pi_*(\mathcal{O}_V/\mathcal{J}^{n+1})$) is the quasi-coherent
$\mathcal{O}_X$-module associated to the $A$-module
$L=\mathfrak{J}/\mathfrak{J}^2$ (resp. $B/\mathfrak{J}^{n+1}$). Since
$\mathcal{J}$ is an \emph{invertible} $\mathcal{O}_V$-module, we may
moreover suppose that $\mathfrak{J}=Bt$, where $t$ is not a zero divisor in
$B$. From the relation $B=A\oplus Bt$ we then deduce, for every $n>0$,
\[
  B=A\oplus At\oplus At^2\oplus\cdots\oplus At^n\oplus Bt^{n+1}
\]
so we have a canonical $A$-isomorphism of the formal power series ring
$A[[T]]$ onto $C=\varprojlim B/\mathfrak{J}^{n+1}$ mapping $T$ to $t$. On
the other hand, we have $L=A\bar t$, where $\bar t$ is the class of $t$
mod. $Bt^2$, and by hypothesis the homomorphism $\varphi$ maps $\bar t$ to an
element $t'\in C$ congruent to $t$ mod. $Ct^2$. We deduce by induction
on $n$ that
\[
  A\oplus At'\oplus\cdots\oplus At'{}^n\oplus Ct^{n+1}
  =A\oplus At\oplus\cdots\oplus At^n\oplus Ct^{n+1}
\]
which proves that the homomorphism $\widehat{\psi}$ indeed corresponds to an
isomorphism of $\prod_{n\geqslant0}L^{\otimes n}$ onto $C$.
\end{proof}

\begin{lemma}[8.10.2]
\label{II.8.10.2}
Under the hypotheses of (\hyperref[II.8.10.1]{8.10.1}), let
$g:X'\to X$ be a morphism,

\oldpage[II]{186}
let $V'=V\times_X X'$, and let $\pi':V'\to X'$,
$g':V'\to V$ be the canonical projections, so that we have the commutative
diagram
\[
  \xymatrix{
    V\ar[d]_{\pi} & V'\ar[l]^{g'}\ar[d]^{\pi'}\\
    X & X'\ar[l]_{g}
  }
\]
Then $j'=j\times1_{X'}$ is an $X'$-section of $V'$ which is a closed
immersion, and $\mathcal{J}'=(g')^*(\mathcal{J})\mathcal{O}_{V'}$ is the
quasi-coherent ideal of $\mathcal{O}_{V'}$ defining the closed sub-prescheme
of $V'$ associated to $j'$. Moreover, we have
\[
  \pi'_*(\mathcal{O}_{V'}/\mathcal{J}'^{n+1})
  =g^*(\pi_*(\mathcal{O}_V/\mathcal{J}^{n+1})).
\]
Finally, $\mathcal{J}'$ is an $\mathcal{O}_{V'}$-module canonically isomorphic
to $(g')^*(\mathcal{J})$, and in particular is invertible if $\mathcal{J}$ is
an invertible $\mathcal{O}_V$-module.
\end{lemma}

\begin{proof}
The fact that $j'$ is a closed immersion follows from
(\hyperref[I.4.3.1]{I, 4.3.1}), and it is an $X'$-section of $V'$ by
functoriality of extension of the base prescheme. Moreover, if $Z$ (resp.
$Z'$) is the closed sub-prescheme of $V$ (resp. $V'$) associated to $j$ (resp.
$j'$), we have $Z'=g'{}^{-1}(Z)$
(\hyperref[I.4.3.1]{I, 4.3.1}), and the second assertion then follows from
(\hyperref[I.4.4.5]{I, 4.4.5}). To prove the other assertions, we
see as in (\hyperref[II.8.10.1]{8.10.1}) that we may reduce to the case
where $X$, $V$ and $X'$ (hence also $V'$) are affine; let us keep the
notation of the proof of (\hyperref[II.8.10.1]{8.10.1}) and let
$X'=\operatorname{Spec}(A')$. We then have $V'=\operatorname{Spec}(B')$ where
$B'=B\otimes_A A'$, and $\mathcal{J}'=\widetilde{\mathfrak{J}'}$, with
$\mathfrak{J}'=\operatorname{Im}(\mathfrak{J}\otimes_A A')$. We therefore have
\[
  B'/\mathfrak{J}'^{n+1}=(B/\mathfrak{J}^{n+1})\otimes_A A'~;
\]
moreover, since $\mathfrak{J}$ is a direct factor (as an $A$-module) of
$B$, $\mathfrak{J}\otimes_A A'$ is a direct factor (as an $A'$-module)
of $B'$, and is therefore canonically identified with $\mathfrak{J}'$.
\end{proof}

\begin{corollary}[8.10.3]
\label{II.8.10.3}
Suppose that the hypotheses of (\hyperref[II.8.10.1]{8.10.1}) are
satisfied, and suppose moreover that $\pi$ is of finite type and that
$\mathcal{J}$ is an invertible $\mathcal{O}_V$-module. Then, for every
$x\in X$, the local ring at the point $j(x)$ of the fibre $\pi^{-1}(x)$ is a
regular (hence integral) ring of dimension 1, whose completion is isomorphic
to the formal power series ring $k(x)[[T]]$ ($T$ an indeterminate); moreover,
there is only one irreducible component of $\pi^{-1}(x)$ containing $j(x)$.
\end{corollary}

\begin{proof}
Since $\pi^{-1}(x)=V\times_X\operatorname{Spec}(k(x))$, we are reduced by
(\hyperref[II.8.10.2]{8.10.2}) to the case where $X$ is the spectrum of a
field $K$. Since $\pi$ is of finite type
(\hyperref[I.6.3.4]{I, 6.3.4, (iv)}), $\mathcal{O}_{j(x)}$ is then a
Noetherian local ring, hence separated for the
$\mathfrak{m}_{j(x)}$-preadic topology
(\hyperref[0.7.3.5]{0, 7.3.5}); it follows from
(\hyperref[II.8.10.1]{8.10.1, (ii) and (iii)}) that the completion of this
ring is isomorphic to $K[[T]]$, and consequently $\mathcal{O}_{j(x)}$ is
regular and of dimension 1 ([1], p.~17-01, th.~1); finally, since
$\mathcal{O}_{j(x)}$ is integral, $j(x)$ belongs to only one of the
(finitely many) irreducible components of $V$
(\hyperref[I.5.1.4]{I, 5.1.4}).
\end{proof}

\begin{corollary}[8.10.4]
\label{II.8.10.4}
Suppose that the hypotheses of (\hyperref[II.8.10.1]{8.10.1}) are
satisfied, and suppose moreover that $\mathcal{J}$ is an invertible
$\mathcal{O}_V$-module. Let $W=V-j(X)$; for every quasi-coherent ideal
$\mathcal{K}$ of $\mathcal{O}_X$, set
$\mathcal{K}_V=\pi^*(\mathcal{K})\mathcal{O}_V$ and
$\mathcal{K}_W=\mathcal{K}_V|W$. Then $\mathcal{K}_V$ is the largest
quasi-coherent ideal of $\mathcal{O}_V$ whose restriction to $W$ is
$\mathcal{K}_W$.
\end{corollary}

\begin{proof}
Indeed, we see as in (\hyperref[II.8.10.1]{8.10.1}) that the question
is local on $X$ and on $V$; we may therefore keep the notation of the
proof of (\hyperref[II.8.10.1]{8.10.1}), with
$\mathfrak{J}=Bt$, where $t$ is not a zero divisor in $B$. Moreover, we have
$W=\operatorname{Spec}(B_t)$ and $\mathcal{K}=\widetilde{\mathfrak{K}}$, where
$\mathfrak{K}$ is an ideal of $A$; whence
$\pi^*(\mathcal{K})\mathcal{O}_V=(\mathfrak{K}.B)^{\sim}$
(\hyperref[I.1.6.9]{I, 1.6.9}),
$\mathcal{K}_W=(\mathfrak{K}.B_t)^{\sim}$, and the largest ideal

\oldpage[II]{187}
of $B$ whose canonical image in $B_t$ is $\mathfrak{K}.B_t$ is the inverse
image of $\mathfrak{K}.B_t$, that is, the set of $s\in B$ such
that $t^ns\in\mathfrak{K}.B$ for some integer $n>0$. We must show that
this last relation implies $s\in\mathfrak{K}.B$, or equivalently that the
canonical image of $t$ is not a zero divisor in
$B/\mathfrak{K}B=(A/\mathfrak{K})\otimes_A B$, which follows from
(\hyperref[II.8.10.2]{8.10.2}) applied to
$X'=\operatorname{Spec}(A/\mathfrak{K})$.
\end{proof}

\begin{corollary}[8.10.5]
\label{II.8.10.5}
Suppose that the hypotheses of (\hyperref[II.8.10.3]{8.10.3}) are
satisfied; let $W=V-j(X)$, let $x$ be a point of $X$, $\mathcal{K}$ a
quasi-coherent ideal of $\mathcal{O}_X$, and $z$ the generic point of the
irreducible component of $\pi^{-1}(x)$ containing $j(x)$
(\hyperref[II.8.10.3]{8.10.3}).
\begin{enumerate}
  \item[\rm{(i)}] Let $g$ be a section of $\mathcal{O}_V$ over $V$
  such that $g|W$ is a section of $\mathcal{K}_W$ over $W$
  (notation of (\hyperref[II.8.10.4]{8.10.4})). Then $g$ is a section
  of $\mathcal{K}_V$; if moreover $g(z)\neq0$, and if, for every integer
  $m>0$, we denote by $g_m^x$ the germ at the point $x$ of the canonical
  image $g_m$ of $g$ in
  $\Gamma(X,\pi_*(\mathcal{O}_V/\mathcal{J}^{m+1}))$, then there exists an
  integer $m>0$ such that the image of $g_m^x$ in
  \[
    (\pi_*(\mathcal{O}_V/\mathcal{J}^{m+1}))_x
    \otimes_{\mathcal{O}_x}k(x)
  \]
  is $\neq0$.
  \item[\rm{(ii)}] Suppose moreover that the conditions of
  (\hyperref[II.8.10.1]{8.10.1, (iii)}) are satisfied. Then, if there exists
  a section $g$ of $\mathcal{K}_V$ over $V$ such that $g(z)\neq0$,
  there exist an integer $n\geqslant0$ and a section $f$ of
  $\mathcal{K}.\mathcal{L}^{\otimes n}=
  \mathcal{K}\otimes\mathcal{L}^{\otimes n}\subset\mathcal{L}^{\otimes n}$
  such that $f(x)\neq0$. If $g$ is a section of $\mathcal{J}$, we can
  take $n>0$.
\end{enumerate}
\end{corollary}

\begin{proof}
\begin{enumerate}
  \item[\rm{(i)}] Since the ideal of $\mathcal{O}_W$ generated by $g|W$ is
  contained in $\mathcal{K}_W$ by hypothesis, the ideal of $\mathcal{O}_V$
  generated by $g$ is contained in $\mathcal{K}_V$ by
  (\hyperref[II.8.10.4]{8.10.4}); in other words, $g$ is a section of
  $\mathcal{K}_V$. To prove the second assertion of (i), we may again
  suppose $X$ and $V$ affine and keep the notation of
  (\hyperref[II.8.10.1]{8.10.1}); the fibre $\pi^{-1}(x)$ is then affine
  with ring $B'=B\otimes_A k(x)$, and there exists in $B'$ an element $t'$
  which is not a zero divisor such that $B'=k(x)\oplus B't'$. Since $j(x)$ is a
  specialization of $z$ and $g(z)\neq0$, we necessarily have
  $g_{j(x)}\neq0$. But $\mathcal{O}_{j(x)}$ is a separated local ring
  (\hyperref[II.8.10.3]{8.10.3}), which therefore embeds in its completion,
  and consequently the image of $g$ in this completion is not zero. Now, this
  completion is isomorphic to $\varprojlim_n(B'/B't'^{n+1})$
  (\hyperref[II.8.10.3]{8.10.3}); if $g'=g\otimes1\in B'$, there therefore
  exists an integer $m$ such that $g'\notin B't'^{m+1}$, or in other words
  the image $g'_m$ of $g'$ in $B'/B't'^{m+1}$ is not zero. But since $g'_m$ is
  none other than the image of $g_m^x$, our assertion is proved.
  \item[\rm{(ii)}] By virtue of
  (\hyperref[II.8.10.1]{8.10.1, (iii)}),
  $\pi_*(\mathcal{O}_V/\mathcal{J}^{m+1})$ is isomorphic to the direct sum
  of the $\mathcal{L}^{\otimes k}$ for $0\leqslant k\leqslant m$; denote
  by $f_k$ the section of $\mathcal{L}^{\otimes k}$ over $X$ which is the
  component of the element of
  $\bigoplus_{k=0}^m\Gamma(X,\mathcal{L}^{\otimes k})$ corresponding to $g_m$
  under this isomorphism. Choosing $m$ as in (i), there is therefore an index
  $k$ such that $f_k(x)\neq0$, by (i). To see that $f_k$ is a section
  of $\mathcal{K}\mathcal{L}^{\otimes k}$, it suffices to consider, as
  above, the case where $X$ and $V$ are affine, and this then follows
  immediately from the fact that $g\in\mathfrak{K}.B$ (notation of
  (\hyperref[II.8.10.4]{8.10.4})). The last assertion follows from the fact
  that the hypothesis $g\in\Gamma(V,\mathcal{J})$ implies $f_0=0$.
\end{enumerate}
\end{proof}

\begin{env}[8.10.6]
\label{II.8.10.6}
\emph{Proof of (\hyperref[II.8.9.4]{8.9.4}).} The question is
local on $Y$ (\hyperref[II.4.6.4]{4.6.4}); since $\varepsilon$ is a
$Y$-section, we can therefore replace $C$ by an affine open neighbourhood $U$
of a point of $\varepsilon(Y)$ such that $\varepsilon(Y)\cap U$ is closed in
$U$. In other words, we may suppose $C$ affine and $Y$ a closed sub-prescheme
of $C$ (hence affine) defined by a quasi-

\oldpage[II]{188}
coherent sheaf $\mathcal{I}$ of ideals of $\mathcal{O}_C$. Since $p$ is
separated and quasi-compact, $X$ is then a quasi-compact scheme, and we are
reduced to proving that $\mathcal{L}$ is ample
(\hyperref[II.4.6.4]{4.6.4}). By virtue of the criterion
(\hyperref[II.4.5.2]{4.5.2, a)}), we must therefore prove the following: for
every quasi-coherent ideal $\mathcal{K}$ of $\mathcal{O}_X$ and every point
$x\in X$ not belonging to the support of $\mathcal{O}_X/\mathcal{K}$, there
exist an integer $n>0$ and a section $f$ of
$\mathcal{K}\otimes\mathcal{L}^{\otimes n}$ over $X$ such that
$f(x)\neq0$.

For this, set
\[
  \mathcal{K}_V=\pi^*(\mathcal{K})\mathcal{O}_V
\]
and
\[
  \mathcal{K}_W=\mathcal{K}_V|W,\qquad\hbox{where}\qquad W=V-j(X)~;
\]
since the restriction of $q$ to $W$ is a quasi-compact immersion into $C$,
it follows from (\hyperref[I.9.4.2]{I, 9.4.2}) that $\mathcal{K}_W$ is the
restriction to $W$ of a quasi-coherent ideal $\mathcal{K}'_V$ of
$\mathcal{O}_V$ of the form
\[
  \mathcal{K}'_V=q^*(\mathcal{K}_C)\mathcal{O}_V
\]
where $\mathcal{K}_C$ is a quasi-coherent ideal of $\mathcal{O}_C$. Moreover,
since by hypothesis $q^{-1}(Y)\subset j(X)$ and $Y$ is defined by the ideal
$\mathcal{I}$, the restriction to $W$ of $q^*(\mathcal{I})\mathcal{O}_V$ is
identical to that of $\mathcal{O}_V$, so $\mathcal{K}_W$ is also the
restriction to $W$ of $q^*(\mathcal{I}\mathcal{K}_C)\mathcal{O}_V$, and we
may therefore suppose that $\mathcal{K}_C\subset\mathcal{I}$, whence
\begin{equation}
\label{II.8.10.6.1}
  \mathcal{K}'_V\subset q^*(\mathcal{I})\mathcal{O}_V\subset\mathcal{J}
\tag{8.10.6.1}
\end{equation}
taking into account (\hyperref[I.4.4.6]{I, 4.4.6}) and the commutativity of
(\hyperref[II.8.9.4.1]{8.9.4.1}). Moreover, we deduce from
(\hyperref[II.8.10.4]{8.10.4}) that
\begin{equation}
\label{II.8.10.6.2}
  \mathcal{K}'_V\subset\mathcal{K}_V.
\tag{8.10.6.2}
\end{equation}

This being so, it follows from (\hyperref[II.8.10.3]{8.10.3}) that $j(x)$
belongs to only one irreducible component of $\pi^{-1}(x)$; let $z$ be the
generic point of this component, and set $z'=q(z)$. By virtue of
(\hyperref[II.8.10.5]{8.10.5}), the proof will be complete (taking into
account (\hyperref[II.8.10.6.1]{8.10.6.1}) and
(\hyperref[II.8.10.6.2]{8.10.6.2})) if we prove the existence of a
section $g$ of $\mathcal{K}'_V$ over $V$ such that $g(z)\neq0$. Now,
by hypothesis, $\mathcal{K}$ has a restriction equal to that of
$\mathcal{O}_X$ in an open neighbourhood of $x$; on the other hand, it follows
from (\hyperref[II.8.10.3]{8.10.3}) that $z\neq j(x)$, hence $z\in W$, and
consequently $(\mathcal{K}_W)_z=\mathcal{O}_{V,z}$, whence by definition
$(\mathcal{K}_C)_{z'}=\mathcal{O}_{C,z'}$. Since $C$ is affine, there is
therefore a section $g'$ of $\mathcal{K}_C$ over $C$ such that
$g'(z')\neq0$, and by taking for $g$ the section of $\mathcal{K}'_V$
corresponding canonically to $g'$, we indeed have $g(z)\neq0$, which completes
the proof.
\end{env}

\begin{remark}[8.10.7]
\label{II.8.10.7}
We do not know whether, in the statement of (\hyperref[II.8.9.4]{8.9.4}),
condition (ii) is superfluous or not. In any case, the conclusion no longer
holds if we do not suppose the existence of a $Y$-morphism $\pi:V\to X$ such
that $\pi\circ j=1_X$; let us indicate briefly how one can indeed construct
a counterexample, whose details can only be developed later.
Take $Y=\operatorname{Spec}(k)$ where $k$ is a field,
$C=\operatorname{Spec}(A)$ where $A=k[T_1,T_2]$, and the $Y$-section
$\varepsilon$ corresponding to the augmentation homomorphism $A\to k$. Denote
by $C'$ the scheme deduced from $C$ by blowing up the closed point
$a=\varepsilon(Y)$ of $C$; if $D$ is the inverse image of $a$ in $C'$,
consider a closed point $b$ in $D$,

\oldpage[II]{189}
and denote by $V$ the scheme deduced from $C'$ by blowing up $b$; $X$ is
the closed sub-prescheme of $V$ that is the inverse image of $a$ under the
structural morphism $q:V\to C$. One shows that $X$ is the union of two
irreducible components $X_1$, $X_2$, where $X_1$ is the inverse image of $b$
in $V$. It is immediate that the ideal $\mathcal{J}$ of $\mathcal{O}_V$ which
defines $X$ is again invertible, but one proves that
$j^*(\mathcal{J})=\mathcal{L}$ ($j$ the canonical injection $X\to V$) is not
ample, by considering the ``degree'' of the inverse image of
$\mathcal{L}$ on $X_1$, which should be $>0$ if $\mathcal{L}$ were ample,
and which one shows (by an elementary intersection computation) to be
equal to 0.
\end{remark}

\subsection{Uniqueness of contractions}
\label{subsection:II.8.11}

\begin{lemma}[8.11.1]
\label{II.8.11.1}
Let $U$, $V$ be two preschemes and
$h=(h_0,\lambda):U\to V$ a surjective morphism. Suppose that:
\begin{enumerate}
  \item[$1^\circ$]
  $\lambda:\mathcal{O}_V\to h_*(\mathcal{O}_U)
  =(h_0)_*(\mathcal{O}_U)$ is an isomorphism.
  \item[$2^\circ$] The underlying space of $V$ is identified with the
  quotient space of the underlying space of $U$ by the relation
  $R:h_0(x)=h_0(y)$ (a condition always satisfied when the morphism $h$ is
  open, or closed, or \emph{a fortiori} when $h$ is proper).
\end{enumerate}
Then, for every prescheme $W$, the map
\begin{equation}
\label{II.8.11.1.1}
  \operatorname{Hom}(V,W)\longrightarrow\operatorname{Hom}(U,W)
\tag{8.11.1.1}
\end{equation}
which sends every morphism $v=(v_0,\nu)$ from $V$ to $W$ to the
morphism $u=v\circ h=(u_0,\mu)$ is a bijection of
$\operatorname{Hom}(V,W)$ onto the set of $u$ such that $u_0$ is
constant on every fibre $h_0^{-1}(x)$.
\end{lemma}

\begin{proof}
It is clear that if $u=v\circ h$, so that $u_0=v_0\circ h_0$, then $u_0$ is
constant on every set $h_0^{-1}(x)$. Conversely, if $u$ has this
property, let us show that there exists one and only one
$v\in\operatorname{Hom}(V,W)$ such that $u=v\circ h$. The existence and
uniqueness of the continuous map $v_0:V\to W$ such that $u_0=v_0\circ h_0$
follow from the hypothesis, since $h_0$ is identified with the canonical
map of $U$ onto $U/R$. On the other hand, replacing $V$ by an isomorphic
prescheme if necessary, we may suppose that $\lambda$ is the identity; by
hypothesis, $\mu$ is then a homomorphism
\[
  \mu:\mathcal{O}_W\longrightarrow(u_0)_*(\mathcal{O}_U)
  =(v_0)_*((h_0)_*(\mathcal{O}_U))
\]
such that the corresponding homomorphism
$\mu^\sharp:u_0^*(\mathcal{O}_W)\to\mathcal{O}_U$ is local on each
stalk. Since
$(v_0)_*((h_0)_*(\mathcal{O}_U))=(v_0)_*(\mathcal{O}_V)$, we must
necessarily have $\nu=\mu$, and everything comes down to seeing that the
corresponding homomorphism $\nu^\sharp:v_0^*(\mathcal{O}_W)\to\mathcal{O}_V$
is local on each stalk. Now, every $y\in V$ is of the form $h_0(x)$ for some
$x\in U$; set $z=v_0(y)=u_0(x)$. Then
(\hyperref[0.3.5.5]{0, 3.5.5}) the homomorphism $\mu_x^\sharp$ factors
as
\[
  \mu_x^\sharp:\mathcal{O}_z
  \xrightarrow{\nu_y^\sharp}\mathcal{O}_y
  \xrightarrow{\lambda_x^\sharp}\mathcal{O}_x.
\]
By hypothesis $\lambda_x^\sharp$ and $\mu_x^\sharp$ are local
homomorphisms; hence $\lambda_x^\sharp$ maps every invertible element of
$\mathcal{O}_y$ to an invertible element of $\mathcal{O}_x$; if
$\nu_y^\sharp$ mapped a non-invertible element of $\mathcal{O}_z$ to
an invertible element of $\mathcal{O}_y$, $\mu_x^\sharp$ would map this
element of $\mathcal{O}_z$ to an invertible element of $\mathcal{O}_x$,
contrary to the hypothesis, whence the lemma.
\end{proof}

\begin{corollary}[8.11.2]
\label{II.8.11.2}
Let $U$ be an integral prescheme and $V$ a normal prescheme; then every
morphism $h:U\to V$ which is universally closed, birational and radicial is an
isomorphism.
\end{corollary}

\begin{proof}
If $h=(h_0,\lambda)$, it follows from the hypotheses that $h_0$ is injective
and closed, and that $h_0(U)$

\oldpage[II]{190}
is dense in $V$, so $h_0$ is a \emph{homeomorphism} of $U$ onto $V$.
To prove the corollary, it will suffice to see that
$\lambda:\mathcal{O}_V\to(h_0)_*(\mathcal{O}_U)$ is an isomorphism: we
can then apply (\hyperref[II.8.11.1]{8.11.1}), which will prove that
the map (\hyperref[II.8.11.1.1]{8.11.1.1}) is bijective (the fibres
$h_0^{-1}(x)$ each being reduced to a single point); hence $h$ will be an
isomorphism. The question being evidently local on $V$, we may suppose
that $V=\operatorname{Spec}(A)$ is affine with an integral and integrally
closed ring (\hyperref[II.8.8.6.1]{8.8.6.1}); $h$ then corresponds
(\hyperref[I.2.2.4]{I, 2.2.4}) to a homomorphism
$\varphi:A\to\Gamma(U,\mathcal{O}_U)$, and everything comes down to seeing
that $\varphi$ is an isomorphism. Now, if $K$ is the field of fractions of
$A$, then by hypothesis $\Gamma(U,\mathcal{O}_U)$ has $K$ as its field of
fractions, and $A$ is a subring of $\Gamma(U,\mathcal{O}_U)$, $\varphi$ being
the canonical injection (\hyperref[I.8.2.7]{I, 8.2.7}). Since the morphism
$h$ satisfies the hypotheses of (\hyperref[II.7.3.11]{7.3.11}),
$\Gamma(U,\mathcal{O}_U)$ is a subring of the integral closure of $A$
in $K$, and is therefore identical to $A$ by hypothesis.
\end{proof}

\begin{remark}[8.11.3]
\label{II.8.11.3}
We will see in Chapter III (\hyperref[III.4.4.11]{III, 4.4.11}) that when
$V$ is a \emph{locally Noetherian} prescheme, every morphism $h:U\to V$
which is proper and quasi-finite (in particular every morphism satisfying the
hypotheses of (\hyperref[II.8.11.2]{8.11.2})) is necessarily
\emph{finite}. The conclusion of (\hyperref[II.8.11.2]{8.11.2}) therefore
follows in this case from (\hyperref[II.6.1.15]{6.1.15}).
\end{remark}

\begin{env}[8.11.4]
\label{II.8.11.4}
We will now see that, in Grauert's criterion
(\hyperref[II.8.9.1]{8.9.1}), one can often assert that the prescheme
$C$ and the ``contraction'' $q$ are determined in an essentially
unique way.
\end{env}

\begin{lemma}[8.11.5]
\label{II.8.11.5}
Let $Y$ be a prescheme, $p:X\to Y$ a proper morphism, $\mathcal{L}$ a
$p$-ample invertible $\mathcal{O}_X$-module, $C$ a $Y$-prescheme,
$\varepsilon:Y\to C$ a $Y$-section, and
$q:V=\mathbf{V}(\mathcal{L})\to C$ a $Y$-morphism, such that the diagram
(\hyperref[II.8.9.1.1]{8.9.1.1}) is commutative. Suppose moreover that,
if $p=(p_0,\theta)$, then
$\theta:\mathcal{O}_Y\to p_*(\mathcal{O}_X)$ is an isomorphism. Now let
$\mathcal{S}'=\bigoplus_{n\geqslant0}p_*(\mathcal{L}^{\otimes n})$,
$C'=\operatorname{Spec}(\mathcal{S}')$, and let
$q':\mathbf{V}(\mathcal{L})\to C'$ be the canonical $Y$-morphism
(\hyperref[II.8.8.5]{8.8.5}). There exists one and only one $Y$-morphism
$u:C'\to C$ such that $q=u\circ q'$.
\end{lemma}

\begin{proof}
The hypothesis on $\theta$ implies in particular that $p$ is surjective;
since by virtue of (\hyperref[II.8.8.4]{8.8.4}) the restriction of $q'$ to
$\mathbf{V}(\mathcal{L})-j(X)$ is an \emph{isomorphism} onto
$C'-\varepsilon'(Y)$ ($\varepsilon'$ being the apex section of $C'$), it
follows from (\hyperref[II.8.8.4]{8.8.4}) that $q'$ is \emph{proper and
surjective}; moreover, by virtue of (\hyperref[II.8.8.6]{8.8.6}), if we write
$q'=(q'_0,\tau)$, then
$\tau:\mathcal{O}_{C'}\to q'_*(\mathcal{O}_V)$ is an isomorphism. We are
therefore in a position to apply Lemma
(\hyperref[II.8.11.1]{8.11.1}), and the lemma will be proved if we prove
that $q$ is constant on every fibre $q'{}^{-1}(z')$, where $z'\in C'$. Now, the
condition is trivially satisfied for $z'\notin\varepsilon'(Y)$. On the other
hand, if $z'\in\varepsilon'(Y)$, there exists one and only one $y\in Y$ such
that $z'=\varepsilon'(y)$; and by virtue of the commutativity of
(\hyperref[II.8.8.5.2]{8.8.5.2}) and the fact that $q'$ maps
$\mathbf{V}(\mathcal{L})-j(X)$ into $C'-\varepsilon'(Y)$, we have
$q'{}^{-1}(z')=j(p^{-1}(y))$; the commutativity of the diagram
(\hyperref[II.8.9.1.1]{8.9.1.1}) therefore establishes our assertion.
\end{proof}

\begin{corollary}[8.11.6]
\label{II.8.11.6}
Under the hypotheses of (\hyperref[II.8.11.5]{8.11.5}), suppose moreover
that $q$ is proper and that the restriction of $q$ to
$\mathbf{V}(\mathcal{L})-j(X)$ is an isomorphism onto
$C-\varepsilon(Y)$. Then the morphism $u$ is universally closed,
surjective and radicial, and its restriction to $C'-\varepsilon'(Y)$ is an
isomorphism onto $C-\varepsilon(Y)$.
\end{corollary}

\begin{proof}
Since $q'$ is an isomorphism of $\mathbf{V}(\mathcal{L})-j(X)$ onto
$C'-\varepsilon'(Y)$ (\hyperref[II.8.8.4]{8.8.4}), the last assertion
follows immediately from $q=u\circ q'$. Moreover, the commutativity of the
diagrams

\oldpage[II]{191}
(\hyperref[II.8.8.5.2]{8.8.5.2}) and
(\hyperref[II.8.9.1.1]{8.9.1.1}) shows that the restriction of $u$ to the
closed sub-prescheme $\varepsilon'(Y)$ of $C'$ is an isomorphism onto the
closed sub-prescheme $\varepsilon(Y)$ of $C$, from which it follows
immediately that for every $z'\in\varepsilon'(Y)$, if $z=u(z')$, then $u$
defines an isomorphism of $k(z)$ onto $k(z')$. These remarks prove that $u$
is bijective and radicial; moreover, if $\psi$ and $\psi'$ are the structural
morphisms $C\to Y$, $C'\to Y$, we have $\psi'=\psi\circ u$, and since $\psi'$
is separated (\hyperref[II.1.2.4]{1.2.4}), so is $u$
(\hyperref[I.5.5.1]{I, 5.5.1, (v)}). On the other hand, we have seen in the
proof of Lemma (\hyperref[II.8.11.5]{8.11.5}) that $q'$ is
surjective; since $q=u\circ q'$ is proper, we finally conclude from
(\hyperref[II.5.4.3]{5.4.3}) and (\hyperref[II.5.4.9]{5.4.9}) that $u$
is universally closed.
\end{proof}

\begin{proposition}[8.11.7]
\label{II.8.11.7}
Let $Y$ be a prescheme, $X$ an integral prescheme, $p:X\to Y$ a proper
morphism, $\mathcal{L}$ a $p$-ample invertible $\mathcal{O}_X$-module, $C$ a
normal $Y$-prescheme, $\varepsilon:Y\to C$ a $Y$-section, and
$q:V=\mathbf{V}(\mathcal{L})\to C$ a $Y$-morphism, such that the diagram
(\hyperref[II.8.9.1.1]{8.9.1.1}) is commutative. Suppose that $q$ is
proper and that the restriction of $q$ to $V-j(X)$ is an isomorphism onto
$C-\varepsilon(Y)$. Suppose moreover that,
if $p=(p_0,\theta)$, then $\theta:\mathcal{O}_Y\to p_*(\mathcal{O}_X)$ is an
isomorphism. Now let
$\mathcal{S}'=\bigoplus_{n\geqslant0}p_*(\mathcal{L}^{\otimes n})$,
$C'=\operatorname{Spec}(\mathcal{S}')$, and let
$q':\mathbf{V}(\mathcal{L})\to C'$ be the canonical $Y$-morphism
(\hyperref[II.8.8.5]{8.8.5}). Then the unique $Y$-morphism $u:C'\to C$ such
that $q=u\circ q'$ is an isomorphism.
\end{proposition}

\begin{proof}
It follows from (\hyperref[II.8.8.6]{8.8.6}) that $C'$ is integral; since
$u$ is a homeomorphism of the underlying spaces $C'\to C$ ($u$ being
bijective and closed by (\hyperref[II.8.11.6]{8.11.6})), $C$ is
irreducible, hence integral, and since the restriction of $u$ to a non-empty
open subset of $C'$ is an isomorphism onto an open subset of $C$, $u$ is
birational. Since $C$ is assumed to be normal, it suffices to apply
(\hyperref[II.8.11.2]{8.11.2}) to obtain the conclusion.
\end{proof}

\begin{remark}[8.11.8]
\label{II.8.11.8}
\begin{enumerate}
  \item[\rm{(i)}] Note that the hypothesis that $C$ is normal implies that
  $X$ is normal as well. Indeed, $C'=\operatorname{Spec}(\mathcal{S}')$ is
  then normal, being isomorphic to $C$, and integral by virtue of
  (\hyperref[II.8.8.6]{8.8.6}); we conclude that
  $\operatorname{Proj}(\mathcal{S}')$ is \emph{normal}. Indeed, the question
  is local on $Y$; if $Y$ is affine and $\mathcal{S}'=\widetilde{S'}$,
  the ring $S'=\Gamma(C',\mathcal{O}_{C'})$ is integral and integrally closed
  (\hyperref[II.8.8.6.1]{8.8.6.1}), so, for every homogeneous element
  $f\in S'_+$, the graded ring $S'_f$ is integral and integrally closed
  ([13], t.~I, p.~257 and 261), and consequently so is the ring $S'_{(f)}$
  of its terms of degree 0, since the intersection of $S'_f$ and the field
  of fractions of $S'_{(f)}$ is equal to $S'_{(f)}$; this proves our
  assertion (\hyperref[II.6.3.4]{6.3.4}). Finally, since $X$ is isomorphic
  to an open sub-prescheme of $\operatorname{Proj}(\mathcal{S}')$
  (\hyperref[II.8.8.1]{8.8.1}), $X$ is indeed normal. We can therefore
  express Proposition (\hyperref[II.8.11.7]{8.11.7}) in the following form:
  \emph{If $X$ is integral and normal, and
  $p=(p_0,\theta):X\to Y$ is a proper morphism such that
  $\theta:\mathcal{O}_Y\to p_*(\mathcal{O}_X)$ is an isomorphism, then,
  for every $p$-ample $\mathcal{O}_X$-module $\mathcal{L}$, there is one
  and only one way of contracting the zero section of
  $V=\mathbf{V}(\mathcal{L})$ so as to obtain a normal $Y$-scheme $C$ and
  a proper $Y$-morphism $q:V\to C$.}
  \item[\rm{(ii)}] When $p$ is proper, the hypothesis
  $p_*(\mathcal{O}_X)=\mathcal{O}_Y$ can be regarded as an auxiliary
  hypothesis, not constituting a genuine restriction of the generality of the
  result. Indeed, if it is not satisfied, it suffices to replace $Y$
  by the $Y$-scheme
  $Y'=\operatorname{Spec}(p_*(\mathcal{O}_X))$ and to regard $X$ as a
  $Y'$-scheme. We will return to this general method in Chapter III,
  \S~4.
\end{enumerate}
\end{remark}


\subsection{Quasi-coherent sheaves on projecting cones}
\label{subsection:II.8.12}

\begin{env}[8.12.1]
\label{II.8.12.1}
Let us return to the hypotheses and notation of
(\hyperref[II.8.3.1]{8.3.1}). Let $\mathcal{M}$ be a quasi-coherent
graded $\mathcal{S}$-module; to avoid any confusion, we will denote by
$\widetilde{\mathcal{M}}$ the quasi-coherent $\mathcal{O}_C$-module

\oldpage[II]{192}
associated to $\mathcal{M}$
(\hyperref[II.1.4.3]{1.4.3}) when $\mathcal{M}$ is regarded as an
\emph{ungraded} $\mathcal{S}$-module, and by
$\operatorname{Proj}_0(\mathcal{M})$ the quasi-coherent
$\mathcal{O}_X$-module associated to $\mathcal{M}$, $\mathcal{M}$ this time
being regarded as a graded $\mathcal{S}$-module (in other words, the
$\mathcal{O}_X$-module denoted $\widetilde{\mathcal{M}}$ in
(\hyperref[II.3.2.2]{3.2.2})). We will moreover set
\begin{equation}
\label{II.8.12.1.1}
  \mathcal{M}_X=\operatorname{Proj}(\mathcal{M})
  =\bigoplus_{n\in\mathbf{Z}}\operatorname{Proj}_0(\mathcal{M}(n))~;
\tag{8.12.1.1}
\end{equation}
the quasi-coherent graded $\mathcal{O}_X$-algebra $\mathcal{S}_X$ being
defined by (\hyperref[II.8.6.1.1]{8.6.1.1}),
$\operatorname{Proj}(\mathcal{M})$ is endowed with a structure of
(quasi-coherent) graded $\mathcal{S}_X$-module by means of the canonical
homomorphisms (\hyperref[II.3.2.6.1]{3.2.6.1})
\begin{equation}
\label{II.8.12.1.2}
  \mathcal{O}_X(m)\otimes_{\mathcal{O}_X}
  \operatorname{Proj}_0(\mathcal{M}(n))
  \to \operatorname{Proj}_0(\mathcal{S}(m)\otimes_{\mathcal{S}}\mathcal{M}(n))
  \to \operatorname{Proj}_0(\mathcal{M}(m+n))
\tag{8.12.1.2}
\end{equation}
the verification of the module axioms being carried out with the help of the
commutative diagram (\hyperref[II.2.5.11.4]{2.5.11.4}).

If $Y=\operatorname{Spec}(A)$ is affine, $\mathcal{S}=\widetilde{S}$ and
$\mathcal{M}=\widetilde{M}$, where $S$ is a graded $A$-algebra and $M$ a
graded $S$-module, then, for every homogeneous element $f\in S_+$, we have
\begin{equation}
\label{II.8.12.1.3}
  \Gamma(X_f,\operatorname{Proj}(\widetilde{M}))=M_f
\tag{8.12.1.3}
\end{equation}
by virtue of the definitions and of (\hyperref[II.8.2.9.1]{8.2.9.1}).

Now consider the quasi-coherent graded
$\widehat{\mathcal{S}}$-module
\begin{equation}
\label{II.8.12.1.4}
  \widehat{\mathcal{M}}
  =\mathcal{M}\otimes_{\mathcal{S}}\widehat{\mathcal{S}}
\tag{8.12.1.4}
\end{equation}
($\widehat{\mathcal{S}}$ being defined by
(\hyperref[II.8.3.1.1]{8.3.1.1})); from it we deduce a quasi-coherent
graded $\mathcal{O}_{\widehat{C}}$-module
$\operatorname{Proj}_0(\widehat{\mathcal{M}})$, which we will also denote by
\begin{equation}
\label{II.8.12.1.5}
  \mathcal{M}^{\square}=\operatorname{Proj}_0(\widehat{\mathcal{M}}).
\tag{8.12.1.5}
\end{equation}
It is clear (\hyperref[II.3.2.4]{3.2.4}) that
$\mathcal{M}^{\square}$ is an additive exact functor in $\mathcal{M}$,
commuting with direct sums and inductive limits.
\end{env}

\begin{proposition}[8.12.2]
\label{II.8.12.2}
With the notation of (\hyperref[II.8.3.2]{8.3.2}), we have functorial
canonical isomorphisms
\begin{equation}
\label{II.8.12.2.1}
  i^*(\mathcal{M}^{\square})\xrightarrow{\sim}\widetilde{\mathcal{M}},
  \qquad
  j^*(\mathcal{M}^{\square})\xrightarrow{\sim}
  \operatorname{Proj}_0(\mathcal{M}).
\tag{8.12.2.1}
\end{equation}
\end{proposition}

\begin{proof}
Indeed, $i^*(\mathcal{M}^{\square})$ is canonically identified with
$\left(\widehat{\mathcal{M}}/(z-1)\widehat{\mathcal{M}}\right)^{\sim}$
on
$\operatorname{Spec}(\widehat{\mathcal{S}}/(z-1)\widehat{\mathcal{S}})$
by virtue of (\hyperref[II.3.2.3]{3.2.3}); the first of the canonical
isomorphisms (\hyperref[II.8.12.2.1]{8.12.2.1}) is then deduced immediately
(\hyperref[II.1.4.1]{1.4.1}) from the canonical isomorphism
$\widehat{\mathcal{M}}/(z-1)\widehat{\mathcal{M}}
\xrightarrow{\sim}\mathcal{M}$. On the other hand, the canonical immersion
$j:X\to\widehat{C}$ corresponds to the canonical homomorphism
$\widehat{\mathcal{S}}\to\mathcal{S}$ with kernel
$z\widehat{\mathcal{S}}$ (\hyperref[II.8.3.2]{8.3.2}); the second
homomorphism (\hyperref[II.8.12.2.1]{8.12.2.1}) is the particular case
of the canonical homomorphism (\hyperref[II.3.5.2]{3.5.2, (ii)}), since
here we have
$\widehat{\mathcal{M}}\otimes_{\widehat{\mathcal{S}}}\mathcal{S}
=\mathcal{M}$; to verify that it is an isomorphism, we may restrict to the
case where $Y=\operatorname{Spec}(A)$ is affine,
$\mathcal{S}=\widetilde{S}$, $\mathcal{M}=\widetilde{M}$; referring to
(\hyperref[II.2.8.8]{2.8.8}), the verification that for every homogeneous
$f$ in $S_+$ the preceding homomorphism, restricted to $X_f$, reduces to an
isomorphism is then immediate.
\end{proof}

\oldpage[II]{193}
By abuse of language, we will again say, because of the existence of the first
isomorphism (\hyperref[II.8.12.2.1]{8.12.2.1}), that
$\mathcal{M}^{\square}$ is the \emph{projective closure} of the
$\mathcal{O}_C$-module $\widetilde{\mathcal{M}}$ (it being understood that the
datum of the $\mathcal{O}_C$-module $\widetilde{\mathcal{M}}$ includes the
grading of the $\mathcal{S}$-module $\mathcal{M}$).

\begin{env}[8.12.3]
\label{II.8.12.3}
With the notation of (\hyperref[II.8.3.5]{8.3.5}), we have a functorial
canonical homomorphism
\begin{equation}
\label{II.8.12.3.1}
  p^*(\operatorname{Proj}_0(\mathcal{M}))
  \to\mathcal{M}^{\square}|\widehat{E}.
\tag{8.12.3.1}
\end{equation}
Indeed, it is a particular case of the homomorphism $u^{\sharp}$ defined in
general in (\hyperref[II.3.5.6]{3.5.6}). If
$Y=\operatorname{Spec}(A)$ is affine, $\mathcal{S}=\widetilde{S}$,
$\mathcal{M}=\widetilde{M}$, we see, referring to
(\hyperref[II.2.8.8]{2.8.8}), that the restriction of
(\hyperref[II.8.12.3.1]{8.12.3.1}) to
$p^{-1}(X_f)=\widehat{C}_f$ (for a homogeneous $f$ in $S_+$) corresponds to
the canonical homomorphism
\begin{equation}
\label{II.8.12.3.2}
  M_{(f)}\otimes_{S_{(f)}}S_f^{\leq}\to M_f^{\leq}
\tag{8.12.3.2}
\end{equation}
taking into account (\hyperref[II.8.2.3.2]{8.2.3.2}) and
(\hyperref[II.8.2.5.2]{8.2.5.2}).
\end{env}

\begin{env}[8.12.4]
\label{II.8.12.4}
Let us place ourselves under the hypotheses of (\hyperref[II.8.5.1]{8.5.1}),
whose notation we keep. It follows from
(\hyperref[II.1.5.6]{1.5.6}) that for every quasi-coherent graded
$\mathcal{S}$-module $\mathcal{M}$ we have, on the one hand, a canonical
isomorphism
\begin{equation}
\label{II.8.12.4.1}
  \Phi^*(\widetilde{\mathcal{M}})
  \xrightarrow{\sim}
  \left(q^*(\mathcal{M})\otimes_{q^*(\mathcal{S})}\mathcal{S}'\right)^{\sim}
\tag{8.12.4.1}
\end{equation}
of $\mathcal{O}_{C'}$-modules; and, on the other hand,
(\hyperref[II.3.5.6]{3.5.6}) implies the existence of a canonical
$\operatorname{Proj}(\varphi)$-morphism
\begin{equation}
\label{II.8.12.4.2}
  \operatorname{Proj}_0(\mathcal{M})\to
  \left(\operatorname{Proj}_0
  (q^*(\mathcal{M})\otimes_{q^*(\mathcal{S})}\mathcal{S}')\right)
  |G(\varphi)
\tag{8.12.4.2}
\end{equation}
and also of a canonical $\widehat{\Phi}$-morphism
\begin{equation}
\label{II.8.12.4.3}
  \operatorname{Proj}_0(\widehat{\mathcal{M}})\to
  \left(\operatorname{Proj}_0
  (q^*(\widehat{\mathcal{M}})
  \otimes_{q^*(\widehat{\mathcal{S}})}\widehat{\mathcal{S}'})\right)
  |G(\widehat{\varphi}).
\tag{8.12.4.3}
\end{equation}
\end{env}

\begin{env}[8.12.5]
\label{II.8.12.5}
Now consider the situation of (\hyperref[II.8.6.1]{8.6.1}),
with the same notation; in the preceding we therefore take $Y'=X$, the
morphism $q:X\to Y$ being the structural morphism, and $\varphi$ is the
canonical $q$-morphism (\hyperref[II.8.6.1.2]{8.6.1.2}). We then have a
canonical isomorphism
\begin{equation}
\label{II.8.12.5.1}
  q^*(\mathcal{M})\otimes_{q^*(\mathcal{S})}\mathcal{S}_X^{\geq}
  \xrightarrow{\sim}\mathcal{M}_X^{\geq}
\tag{8.12.5.1}
\end{equation}
setting
$\mathcal{M}_X^{\geq}=\bigoplus_{n\geq0}\operatorname{Proj}_0(\mathcal{M}(n))$.
Indeed, we may restrict to the case where $Y=\operatorname{Spec}(A)$ is
affine, $\mathcal{S}=\widetilde{S}$ and $\mathcal{M}=\widetilde{M}$, and
define the isomorphism (\hyperref[II.8.12.5.1]{8.12.5.1}) on each of the
affine open subsets $X_f$ ($f$ homogeneous in $S_+$), verifying the
compatibility with passage to a homogeneous multiple of $f$. Now, the
restriction to $X_f$ of the left-hand side of
(\hyperref[II.8.12.5.1]{8.12.5.1}) is
$\widetilde{M}'=
\left((M\otimes_A S_{(f)})
\otimes_{S\otimes_A S_{(f)}}S_f^{\geq}\right)^{\sim}$ by
(\hyperref[II.8.6.2.1]{8.6.2.1}); since we have a canonical isomorphism of
$M\otimes_A S_{(f)}$ onto $M\otimes_S(S\otimes_A S_{(f)})$, we deduce from
it an isomorphism of $\widetilde{M}'$ onto
$\left(M\otimes_S S_f^{\geq}\right)^{\sim}$, and the latter is
canonically isomorphic by (\hyperref[II.8.2.9.1]{8.2.9.1}) to the
restriction to $X_f$ of the right-hand side of
(\hyperref[II.8.12.5.1]{8.12.5.1}), and satisfies the required
compatibility conditions.

\oldpage[II]{194}
Replacing $\mathcal{M}$ by $\widehat{\mathcal{M}}$, $\mathcal{S}$ by
$\widehat{\mathcal{S}}$ and $\mathcal{S}_X$ by
$(\mathcal{S}_X^{\geq})^{\wedge}$ in the preceding argument, we likewise
have a canonical isomorphism
\begin{equation}
\label{II.8.12.5.2}
  q^*(\widehat{\mathcal{M}})
  \otimes_{q^*(\widehat{\mathcal{S}})}(\mathcal{S}_X^{\geq})^{\wedge}
  \xrightarrow{\sim}(\mathcal{M}_X^{\geq})^{\wedge}.
\tag{8.12.5.2}
\end{equation}
If we recall (\hyperref[II.8.6.2]{8.6.2}) that the structural morphism
$u:\operatorname{Proj}(\mathcal{S}_X^{\geq})\to X$ is an
isomorphism, we first deduce from the above that we have a canonical
$u$-isomorphism
\begin{equation}
\label{II.8.12.5.3}
  \operatorname{Proj}_0\mathcal{M}
  \xrightarrow{\sim}\operatorname{Proj}_0(\mathcal{M}_X^{\geq})
\tag{8.12.5.3}
\end{equation}
as a particular case of (\hyperref[II.8.12.4.2]{8.12.4.2}). Indeed, we
observe, with the notation of the proof of
(\hyperref[II.8.6.2]{8.6.2}), that this amounts to seeing that
the canonical homomorphism
$M_{(f)}\otimes_{S_{(f)}}(S_f^{\geq})^{(d)}\to(M_f^{\geq})^{(d)}$ is an
isomorphism when $f\in S_d$, which is immediate.

Secondly, the isomorphism (\hyperref[II.8.12.5.2]{8.12.5.2}) makes it
possible, by applying this time (\hyperref[II.8.12.4.3]{8.12.4.3}) to the
canonical morphism $r=\operatorname{Proj}(\widehat{\alpha}):\widehat{C}_X\to
\widehat{C}$, to obtain a canonical $r$-morphism
\begin{equation}
\label{II.8.12.5.4}
  \mathcal{M}^{\square}\to(\mathcal{M}_X^{\geq})^{\square}.
\tag{8.12.5.4}
\end{equation}
Let us now recall (\hyperref[II.8.6.2]{8.6.2}) that the restrictions of
$r$ to the punctured cones $\widehat{E}_X$ and $E_X$ are \emph{isomorphisms}
onto $\widehat{E}$ and $E$ respectively. Moreover:
\end{env}

\begin{proposition}[8.12.6]
\label{II.8.12.6}
The restrictions to $\widehat{E}_X$ and to $E_X$ of the canonical $r$-morphism
(\hyperref[II.8.12.5.4]{8.12.5.4}) are isomorphisms
\begin{equation}
\label{II.8.12.6.1}
  \mathcal{M}^{\square}|\widehat{E}
  \xrightarrow{\sim}(\mathcal{M}_X^{\geq})^{\square}|\widehat{E}_X
\tag{8.12.6.1}
\end{equation}
\begin{equation}
\label{II.8.12.6.2}
  \widetilde{\mathcal{M}}|E
  \xrightarrow{\sim}(\mathcal{M}_X^{\geq})^{\sim}|E_X.
\tag{8.12.6.2}
\end{equation}
\end{proposition}

\begin{proof}
We reduce to the case where $Y$ is affine as in the proof of
(\hyperref[II.8.6.2]{8.6.2}); with the notation of that proof, and
going back to the definitions (\hyperref[II.2.8.8]{2.8.8}), we must
show that the canonical homomorphism
\[
  \widehat{M}_{(f)}\otimes_{\widehat{S}_{(f)}}
  (S_f^{\geq})_{(f/1)}^{\wedge}
  \to(M\otimes_S S_f^{\geq})_{(f/1)}^{\wedge}
\]
is an isomorphism; but by virtue of
(\hyperref[II.8.2.3.2]{8.2.3.2}) and
(\hyperref[II.8.2.5.2]{8.2.5.2}), the left-hand side is canonically
identified with
$M_f^{\leq}\otimes_{S_f^{\leq}}(S_f^{\geq})_{f/1}^{\leq}$, hence with
$M_f^{\leq}$ by virtue of (\hyperref[II.8.2.7.2]{8.2.7.2}), and the
right-hand side with $(M_f^{\geq})_{f/1}^{\leq}$, hence also with $M_f^{\leq}$
by virtue of (\hyperref[II.8.2.9.2]{8.2.9.2}), whence the conclusion as
regards (\hyperref[II.8.12.6.1]{8.12.6.1}); formula
(\hyperref[II.8.12.6.2]{8.12.6.2}) then follows from
(\hyperref[II.8.12.6.1]{8.12.6.1}) and from
(\hyperref[II.8.12.2.1]{8.12.2.1}).
\end{proof}

\begin{corollary}[8.12.7]
\label{II.8.12.7}
With the identifications of (\hyperref[II.8.6.3]{8.6.3}), the restriction
of $(\mathcal{M}_X^{\geq})^{\square}$ to $\widehat{E}_X$ is identified with
$(\mathcal{M}_X^{\leq})^{\sim}$, and the restriction of
$(\mathcal{M}_X^{\geq})^{\square}$ to $E_X$ is identified with
$\widetilde{\mathcal{M}}_X$.
\end{corollary}

\begin{proof}
Indeed, we are reduced to the affine case, and this follows from the
identification of $(M_f^{\geq})_{f/1}^{\leq}$ with $M_f^{\leq}$ and of
$(M_f^{\geq})_{f/1}$ with $M_f$ (\hyperref[II.8.2.9.2]{8.2.9.2}).
\end{proof}

\begin{proposition}[8.12.8]
\label{II.8.12.8}
Under the hypotheses of (\hyperref[II.8.6.4]{8.6.4}), the canonical
homomorphism (\hyperref[II.8.12.3.1]{8.12.3.1}) is an isomorphism.
\end{proposition}

\begin{proof}
Taking into account the fact that $\operatorname{Proj}(\mathcal{S}_X^{\geq})\to X$
is an isomorphism (\hyperref[II.8.6.2]{8.6.2}) and the

\oldpage[II]{195}
isomorphisms (\hyperref[II.8.12.5.4]{8.12.5.4}) and
(\hyperref[II.8.12.6.1]{8.12.6.1}), we are reduced to proving the
corresponding proposition for the canonical homomorphism
$p_X^*(\operatorname{Proj}_0(\mathcal{M}_X^{\geq}))
\to(\mathcal{M}_X^{\geq})^{\square}|\widehat{E}_X$; in other words, we are
reduced to the case where $\mathcal{S}_1$ is an invertible
$\mathcal{O}_Y$-module and $\mathcal{S}$ is generated by $\mathcal{S}_1$.
With the notation of (\hyperref[II.8.12.3]{8.12.3}), we then have, for an
$f\in S_1$, $S_f^{\leq}=S_{(f)}[1/f]$, and the canonical homomorphism
$M_{(f)}\otimes_{S_{(f)}}S_f^{\leq}\to M_f^{\leq}$ is an isomorphism by
definition of $M_f^{\leq}$.
\end{proof}

\begin{env}[8.12.9]
\label{II.8.12.9}
Now consider the quasi-coherent $\mathcal{S}$-modules
\[
  \mathcal{M}_{[n]}=\bigoplus_{m\geqslant n}\mathcal{M}_m
\]
and (with the notation of (\hyperref[II.8.7.2]{8.7.2})) the
quasi-coherent graded $\mathcal{S}^{\natural}$-module
\begin{equation}
\label{II.8.12.9.1}
  \mathcal{M}^{\natural}
  =\left(\bigoplus_{n\geqslant0}\mathcal{M}_{[n]}\right)^{\sim}.
\tag{8.12.9.1}
\end{equation}
We have seen (\hyperref[II.8.7.3]{8.7.3}) that there exists a canonical
$C$-isomorphism
$h:C_X\xrightarrow{\sim}\operatorname{Proj}(\mathcal{S}^{\natural})$.
Moreover:
\end{env}

\begin{proposition}[8.12.10]
\label{II.8.12.10}
There exists a canonical $h$-isomorphism
\begin{equation}
\label{II.8.12.10.1}
  \operatorname{Proj}_0(\mathcal{M}^{\natural})
  \xrightarrow{\sim}\widetilde{\mathcal{M}}_X.
\tag{8.12.10.1}
\end{equation}
\end{proposition}

\begin{proof}
We argue as in (\hyperref[II.8.7.3]{8.7.3}), this time using the
existence of the di-isomorphism (\hyperref[II.8.2.9.3]{8.2.9.3}) instead
of (\hyperref[II.8.2.7.3]{8.2.7.3}). We leave the details to the
reader.
\end{proof}

\subsection{Projective closures of sub-sheaves and of closed
subschemes}
\label{subsection:II.8.13}

\begin{env}[8.13.1]
\label{II.8.13.1}
The hypotheses and notation being those of
(\hyperref[II.8.12.1]{8.12.1}), consider a quasi-coherent
sub-$\mathcal{S}$-module $\mathcal{N}$ of $\mathcal{M}$, \emph{not
necessarily graded}. We can therefore consider the quasi-coherent
$\mathcal{O}_C$-module $\widetilde{\mathcal{N}}$ associated to $\mathcal{N}$,
which is a sub-$\mathcal{O}_C$-module of $\widetilde{\mathcal{M}}$. On the
other hand, we have seen (\hyperref[II.8.12.2.1]{8.12.2.1}) that
$\widetilde{\mathcal{M}}$ is identified with the restriction of
$\mathcal{M}^{\square}$ to $C$. Since the canonical injection
$i:C\to\widehat{C}$ is an affine morphism
(\hyperref[II.8.3.2]{8.3.2}), and \emph{a fortiori} quasi-compact, the
\emph{canonical extension} $(\widetilde{\mathcal{N}})^{-}$, the largest
sub-$\mathcal{O}_{\widehat{C}}$-module contained in
$\mathcal{M}^{\square}$ and inducing $\widetilde{\mathcal{N}}$ on $C$, is
a quasi-coherent $\mathcal{O}_{\widehat{C}}$-module
(\hyperref[I.9.4.2]{I, 9.4.2}). We will give an explicit description of
it by means of a graded $\widehat{\mathcal{S}}$-module.
\end{env}

\begin{env}[8.13.2]
\label{II.8.13.2}
For this, consider, for every integer $n\geqslant0$, the homomorphism
$\bigoplus_{i\leqslant n}\mathcal{M}_i\to\mathcal{M}$ which, for every open
subset $U$ of $Y$, maps the family
\[
  (s_i)\in\bigoplus_{i\leqslant n}\Gamma(U,\mathcal{M}_i)
\]
to the section $\sum_i s_i\in\Gamma(U,\mathcal{M})$. Denote by
$\mathcal{N}'_n$ the inverse image of $\mathcal{N}$ under this homomorphism,
which is a quasi-coherent sub-$\mathcal{O}_Y$-module of
$\bigoplus_{i\leqslant n}\mathcal{M}_i$. Now consider the
homomorphism
$\bigoplus_{i\leqslant n}\mathcal{M}_i\to
\widehat{\mathcal{M}}=\mathcal{M}[z]$ which maps $(s_i)$ to the
section
$\sum_{i\leqslant n}s_i z^{n-i}\in\Gamma(U,\widehat{\mathcal{M}}_n)$, and
let $\mathcal{N}_n$ be the image of $\mathcal{N}'_n$ under this homomorphism;
one immediately verifies that
$\overline{\mathcal{N}}=\bigoplus_{n\geqslant0}\mathcal{N}_n$ is a
(quasi-coherent) graded sub-$\widehat{\mathcal{S}}$-module of
$\widehat{\mathcal{M}}$; we say that $\overline{\mathcal{N}}$ is deduced from
$\mathcal{N}$ by \emph{homogenization}, by means of the ``homogenizing
variable'' $z$. Note

\oldpage[II]{196}
that if $\mathcal{N}$ is already a \emph{graded} sub-$\mathcal{S}$-module of
$\mathcal{M}$, then $\overline{\mathcal{N}}$ is identified with the direct
sum of the components $\widehat{\mathcal{N}}_n$ of degree $n\geqslant0$ in
$\widehat{\mathcal{N}}=\mathcal{N}[z]$.
\end{env}

\begin{proposition}[8.13.3]
\label{II.8.13.3}
The $\mathcal{O}_{\widehat{C}}$-module
$\operatorname{Proj}_0(\overline{\mathcal{N}})$ is the canonical extension
$(\widetilde{\mathcal{N}})^{-}$ of $\widetilde{\mathcal{N}}$ to
$\widehat{C}$.
\end{proposition}

\begin{proof}
The question is local on $Y$ and $\widehat{C}$ by virtue of the definition of
the canonical extension (\hyperref[I.9.4.1]{I, 9.4.1}). We may therefore
already suppose that $Y=\operatorname{Spec}(A)$ is affine, with
$\mathcal{S}=\widetilde{S}$, $\mathcal{M}=\widetilde{M}$ and
$\mathcal{N}=\widetilde{N}$, where $N$ is a not necessarily graded
sub-$S$-module of $M$. Moreover
(\hyperref[II.8.3.2.6]{8.3.2.6}), $\widehat{C}$ is the union of the affine
open subsets $\widehat{C}_z=C$ and
$\widehat{C}_f=\operatorname{Spec}(S_f^{\leq})$ ($f$ homogeneous in $S_+$).
It will therefore suffice to show that: $1^\circ$ the restriction of
$\operatorname{Proj}_0(\overline{\mathcal{N}})$ to $C$ is
$\widetilde{\mathcal{N}}$; $2^\circ$ the restriction of
$\operatorname{Proj}_0(\overline{\mathcal{N}})$ to each $\widehat{C}_f$ is
the canonical extension of the restriction of $\widetilde{\mathcal{N}}$ to
$C\cap\widehat{C}_f=\operatorname{Spec}(S_f)$
(\hyperref[II.8.3.2.6]{8.3.2.6}). For the first point, note that
$\operatorname{Proj}_0(\overline{\mathcal{N}})|C$ is identified with
$(\overline{N}_{(z)})^{\sim}$
(\hyperref[II.8.3.2.4]{8.3.2.4}); now $\overline{N}_{(z)}$ is canonically
identified (\hyperref[II.2.2.5]{2.2.5}) with the image of
$\overline{N}$ in $\widehat{M}/(z-1)\widehat{M}$, and under the canonical
isomorphism of the latter onto $M$ (\hyperref[II.8.2.5]{8.2.5}), this
image is identified with $N$, by virtue of the definition of $\overline{N}$
given in (\hyperref[II.8.13.2]{8.13.2}).

To prove the second point, note that the injection
$i:C\cap\widehat{C}_f\to\widehat{C}_f$ corresponds to the canonical injection
$S_f^{\leq}\to S_f$ (\hyperref[II.8.3.2.6]{8.3.2.6}); on the other hand, we
have $\Gamma(\widehat{C}_f,\mathcal{M}^{\square})=M_f^{\leq}$,
$\Gamma(\widehat{C}_f,i_*(\widetilde{\mathcal{N}}))=N_f$ and (by
(\hyperref[II.8.12.2.1]{8.12.2.1}))
$\Gamma(\widehat{C}_f,i_*(i^*(\mathcal{M}^{\square})))=M_f$. Taking into
account (\hyperref[I.9.4.2]{I, 9.4.2}), we are therefore reduced to showing
that $\overline{N}_{(f)}\subset\widehat{M}_{(f)}=M_f^{\leq}$ is canonically
identified with the inverse image of $N_f$ under the canonical injection
$M_f^{\leq}\to M_f$. Indeed, let $d=\deg(f)>0$, and suppose that an
element $(\sum_{k\leqslant md}x_k)/f^m$ of $M_f$ (with $x_k\in M_k$) is of
the form $y/f^m$ with $y\in N$. By multiplying $y$ and the $x_k$ by one and
the same suitable $f^h$, we may already suppose that
$\sum_{k\leqslant md}x_k=y$. But in the identification
(\hyperref[II.8.2.5.2]{8.2.5.2}), $(\sum_{k\leqslant md}x_k)/f^m$
corresponds to $\sum_{k\leqslant md}x_kz^{md-k}/f^m$, and this element indeed
belongs to $\overline{N}_{(f)}$, since $\sum_{k\leqslant md}x_k\in N$; the
converse is evident.
\end{proof}

\begin{remark}[8.13.4]
\label{II.8.13.4}
\begin{enumerate}
  \item[\rm{(i)}] The most important case of application of
  (\hyperref[II.8.13.3]{8.13.3}) is the one where
  $\mathcal{M}=\mathcal{S}$, $\widetilde{\mathcal{N}}$ therefore being an
  \emph{arbitrary} quasi-coherent sheaf of ideals $\mathcal{J}$ of
  $\mathcal{O}_C$ (\hyperref[II.1.4.3]{1.4.3}), corresponding
  bijectively to a \emph{closed sub-prescheme} $Z$ of $C$. Then the
  canonical extension $\overline{\mathcal{J}}$ of $\mathcal{J}$ is the
  quasi-coherent sheaf of ideals of $\mathcal{O}_{\widehat{C}}$ which defines
  the \emph{closure} $\overline{Z}$ of $Z$ in $\widehat{C}$
  (\hyperref[I.9.5.10]{I, 9.5.10}); Proposition
  (\hyperref[II.8.13.3]{8.13.3}) gives a canonical way of defining
  $\overline{Z}$ by means of a graded ideal in
  $\widehat{\mathcal{S}}=\mathcal{S}[z]$.
  \item[\rm{(ii)}] Suppose for simplicity that $Y$ is affine, and keep
  the notation of the proof of
  (\hyperref[II.8.13.3]{8.13.3}). For every nonzero $x\in N$, let $d(x)$
  be the largest of the degrees of the homogeneous components $x_i$ of $x$ in
  $M$; $\overline{N}$ is by definition the submodule of $\widehat{M}$
  consisting of $0$ and the elements of the form
  $h(x,k)=z^k\sum_{i\leqslant d(x)}x_i z^{d(x)-i}$ ($k$ an integer
  $\geqslant0$); it is therefore generated, as a module over
  $\widehat{S}=S[z]$, by the elements of the form
  \[
    h(x,0)=\sum_{i\leqslant d(x)}x_i z^{d(x)-i}.
  \]

  \oldpage[II]{197}
  We say that $h(x,0)$ is deduced from $x$ by \emph{homogenization} by
  means of the ``homogenizing variable'' $z$. But since $h(x,0)$ does not
  depend additively on $x$ (nor, \emph{a fortiori}, $S$-linearly), \emph{one
  must beware of believing} (even when $M=S$) that the $h(x,0)$ run through a
  \emph{system of generators} of the graded $\widehat{S}$-module
  $\overline{N}$ when $x$ runs through a \emph{system of
  generators} of the $S$-module $N$. This is nevertheless indeed the case in
  the situation (the only one considered in elementary algebraic geometry)
  where $N$ is a \emph{monogenic free} $S$-module, since if $t$ is a basis of
  $N$, $h(t,0)$ generates the $\widehat{S}$-module $\overline{N}$.
\end{enumerate}
\end{remark}


\subsection{Supplement on the sheaves associated to
graded $\mathcal{S}$-modules}
\label{subsection:II.8.14}

\begin{env}[8.14.1]
\label{II.8.14.1}
Let $Y$ be a prescheme, $\mathcal{S}$ a quasi-coherent $\mathcal{O}_Y$-algebra
graded in \emph{positive degrees},
$X=\operatorname{Proj}(\mathcal{S})$, and $q:X\to Y$ the structural morphism
(separated by virtue of (\hyperref[II.3.1.3]{3.1.3})). We keep the
notation of (\hyperref[II.8.12.1]{8.12.1}), so that we have
defined a functor
$\mathcal{M}_X=\operatorname{Proj}(\mathcal{M})$ in $\mathcal{M}$, from the
category of quasi-coherent graded $\mathcal{S}$-modules to that of
quasi-coherent graded $\mathcal{S}_X$-modules; it is moreover clear
(\hyperref[II.3.2.4]{3.2.4}) that it is an \emph{exact additive} functor,
commuting with inductive limits.

Note moreover that it follows immediately from the definition
(\hyperref[II.8.12.1.1]{8.12.1.1}) that we have
\begin{equation}
\label{II.8.14.1.1}
  \operatorname{Proj}(\mathcal{M}(n))
  =(\operatorname{Proj}(\mathcal{M}))(n)
  \qquad\text{for every }n\in\mathbf{Z}.
\tag{8.14.1.1}
\end{equation}
\end{env}

\begin{env}[8.14.2]
\label{II.8.14.2}
We first extend to the $\mathcal{S}_X$-modules of the form
$\operatorname{Proj}(\mathcal{M})$ the canonical homomorphisms $\lambda$
and $\mu$ defined in (\hyperref[II.3.2.6]{3.2.6}). To this end, note
that, for all $m\in\mathbf{Z}$ and $n\in\mathbf{Z}$, we have, by virtue
of (\hyperref[II.2.1.2.1]{2.1.2.1}), a canonical homomorphism of
$\mathcal{O}_X$-modules
\begin{equation}
\label{II.8.14.2.1}
  \lambda_{mn}:\operatorname{Proj}_0(\mathcal{M}(m))
  \otimes_{\mathcal{O}_X}\operatorname{Proj}_0(\mathcal{N}(n))
  \to\operatorname{Proj}_0
  ((\mathcal{M}\otimes_{\mathcal{S}}\mathcal{N})(m+n))
\tag{8.14.2.1}
\end{equation}
and likewise, by virtue of (\hyperref[II.2.1.2.2]{2.1.2.2}), a
canonical homomorphism of $\mathcal{O}_X$-modules
\begin{equation}
\label{II.8.14.2.2}
  \mu_{mn}:\operatorname{Proj}_0
  ((\mathcal{H}om_{\mathcal{S}}(\mathcal{M},\mathcal{N}))(n-m))
  \to\mathcal{H}om_{\mathcal{O}_X}
  (\operatorname{Proj}_0(\mathcal{M}(m)),
   \operatorname{Proj}_0(\mathcal{N}(n)))
\tag{8.14.2.2}
\end{equation}
for all quasi-coherent graded $\mathcal{S}$-modules
$\mathcal{M}$, $\mathcal{N}$. From this we deduce a homomorphism
\[
  \mu_k:\operatorname{Proj}_0
  ((\mathcal{H}om_{\mathcal{S}}(\mathcal{M},\mathcal{N}))(k))
  \to
  (\mathcal{H}om_{\mathcal{S}_X}
  (\operatorname{Proj}(\mathcal{M}),\operatorname{Proj}(\mathcal{N})))_k
\]
obtained by assigning to every
$u\in\Gamma(U,\operatorname{Proj}_0
((\mathcal{H}om_{\mathcal{S}}(\mathcal{M},\mathcal{N}))(k)))$
the homomorphism $\mu_k(u)$, of degree $k$, of graded $\mathbf{Z}$-modules
$\Gamma(U,\operatorname{Proj}(\mathcal{M}))\to
\Gamma(U,\operatorname{Proj}(\mathcal{N}))$ ($U$ an open subset of $X$) which
on each $\Gamma(U,\operatorname{Proj}_0(\mathcal{M}(m)))$ coincides with
$\mu_{m,m+k}(u)$; moreover, going back to the definition of the $\mu_{mn}$
(\hyperref[II.2.5.12.1]{2.5.12.1}), we see immediately that $\mu_k(u)$ is
in fact a homomorphism of degree $k$ of graded
$\Gamma(U,\mathcal{S}_X)$-modules, and moreover that the $\mu_k$
define a homomorphism of graded $\mathcal{S}_X$-modules
\begin{equation}
\label{II.8.14.2.3}
  \operatorname{Proj}
  (\mathcal{H}om_{\mathcal{S}}(\mathcal{M},\mathcal{N}))
  \to\mathcal{H}om_{\mathcal{S}_X}
  (\operatorname{Proj}(\mathcal{M}),\operatorname{Proj}(\mathcal{N})).
\tag{8.14.2.3}
\end{equation}

Likewise, taking into account the associativity diagram
(\hyperref[II.2.5.11.4]{2.5.11.4}), the homomorphisms
(\hyperref[II.8.14.2.1]{8.14.2.1}) give a homomorphism of
graded $\mathcal{S}_X$-modules
\begin{equation}
\label{II.8.14.2.4}
  \lambda:\operatorname{Proj}(\mathcal{M})
  \otimes_{\mathcal{S}_X}\operatorname{Proj}(\mathcal{N})
  \to\operatorname{Proj}(\mathcal{M}\otimes_{\mathcal{S}}\mathcal{N}).
\tag{8.14.2.4}
\end{equation}
\end{env}

\oldpage[II]{198}
\begin{proposition}[8.14.3]
\label{II.8.14.3}
The homomorphism (\hyperref[II.8.14.2.4]{8.14.2.4}) is bijective; the
same is true of (\hyperref[II.8.14.2.3]{8.14.2.3}) when the
graded $\mathcal{S}$-module $\mathcal{M}$ admits a finite presentation
(\hyperref[II.3.1.1]{3.1.1}).
\end{proposition}

\begin{proof}
The question is evidently local on $X$ and on $Y$; we may therefore suppose
$Y=\operatorname{Spec}(A)$ affine, $\mathcal{S}=\widetilde{S}$,
$\mathcal{M}=\widetilde{M}$, $\mathcal{N}=\widetilde{N}$, where $S$ is an
$A$-algebra graded in positive degrees and $M$ and $N$ are two graded
$S$-modules. If $f$ is a homogeneous element of $S_+$, the homomorphisms
(\hyperref[II.8.14.2.1]{8.14.2.1}) and
(\hyperref[II.8.14.2.2]{8.14.2.2}), restricted to the affine open subset
$D_+(f)$, correspond to the canonical homomorphisms
(\hyperref[II.2.5.11.1]{2.5.11.1}) and
(\hyperref[II.2.5.12.1]{2.5.12.1}):
\begin{align*}
  M(m)_{(f)}\otimes_{S_{(f)}}N(n)_{(f)}
  &\to(M\otimes_S N)(m+n)_{(f)},\\
  (\operatorname{Hom}_S(M,N))(n-m)_{(f)}
  &\to\operatorname{Hom}_{S_{(f)}}(M(m)_{(f)},N(n)_{(f)}).
\end{align*}
Referring to the definitions of these homomorphisms, we therefore see (taking
into account (\hyperref[II.8.2.9.1]{8.2.9.1})) that the restriction of
(\hyperref[II.8.14.2.4]{8.14.2.4}) to $D_+(f)$ corresponds to
the canonical homomorphism
\[
  M_f\otimes_{S_f}N_f\to(M\otimes_S N)_f
\]
defined in (\hyperref[0.1.3.4]{0, 1.3.4}), and we know that the latter is
an isomorphism. Likewise, the restriction of
(\hyperref[II.8.14.2.3]{8.14.2.3}) to $D_+(f)$ corresponds to
the canonical homomorphism (\hyperref[0.1.3.5]{0, 1.3.5})
\[
  (\operatorname{Hom}_S(M,N))_f\to
  \operatorname{Hom}_{S_f}(M_f,N_f)
\]
taking into account that, $M$ being of finite type, the module
$\operatorname{Hom}_S(M,N)$, the direct sum of the subgroups consisting of the
\emph{homogeneous} homomorphisms of $S$-modules
(\hyperref[II.2.1.2]{2.1.2}), coincides with the set of \emph{all}
homomorphisms $M\to N$ of $S$-modules. The hypothesis that $M$ admits a
finite presentation then implies (\hyperref[0.1.3.5]{0, 1.3.5}) that
the canonical homomorphism under consideration is indeed an isomorphism.
\end{proof}

\begin{proposition}[8.14.4]
\label{II.8.14.4}
If $U$ is a quasi-compact open subset of $X$, there exists an integer $d$
such that for every integer $n$ that is a multiple of $d$,
$\mathcal{O}_X(n)|U$ is invertible, its inverse being
$\mathcal{O}_X(-n)|U$.
\end{proposition}

\begin{proof}
Since $q(U)$ is quasi-compact, it is covered by finitely many affine open
subsets $V_i$, so every $x\in U$ is contained in an affine open subset of the
form $D_+(f)$, where $f$ is a homogeneous element of degree $>0$ of one of the
rings $\Gamma(V_i,\mathcal{S})$. Since $U$ is quasi-compact, it can be
covered by finitely many such open subsets $D_+(f_j)$; let $d$ be a
common multiple of the degrees of the $f_j$. The number $d$ answers the
question by virtue of (\hyperref[II.2.5.17]{2.5.17}).
\end{proof}

\begin{env}[8.14.5]
\label{II.8.14.5}
With the hypotheses and notation of (\hyperref[II.8.14.1]{8.14.1}), we
have defined in (\hyperref[II.3.3.2]{3.3.2}) canonical homomorphisms of
$\mathcal{O}_Y$-modules
\begin{equation}
\label{II.8.14.5.1}
  \alpha_n:\mathcal{M}_n\to
  q_*(\operatorname{Proj}_0(\mathcal{M}(n)))
  \qquad(n\in\mathbf{Z})
\tag{8.14.5.1}
\end{equation}
Generalizing the notation of (\hyperref[II.3.3.1]{3.3.1}), we
set, for every graded $\mathcal{S}_X$-module $\mathcal{F}$,
\begin{equation}
\label{II.8.14.5.2}
  \Gamma_*(\mathcal{F})=
  \bigoplus_{n\in\mathbf{Z}}q_*(\mathcal{F}_n).
\tag{8.14.5.2}
\end{equation}
In particular,
$\Gamma_*(\mathcal{S}_X)=\bigoplus_{n\in\mathbf{Z}}q_*(\mathcal{O}_X(n))$
is the graded $\mathcal{O}_Y$-algebra denoted
$\Gamma_*(\mathcal{O}_X)$ in
(\hyperref[II.3.3.1.2]{3.3.1.2}); it is clear that
$\Gamma_*(\mathcal{F})$ is a graded $\Gamma_*(\mathcal{S}_X)$-module
(\hyperref[0.4.2.2]{0, 4.2.2}). When

\oldpage[II]{199}
in the homomorphisms (\hyperref[II.8.14.5.1]{8.14.5.1}) we take
$\mathcal{M}=\mathcal{S}$, we obtain the homomorphism of
graded $\mathcal{O}_Y$-algebras
\begin{equation}
\label{II.8.14.5.3}
  \alpha:\mathcal{S}\to\Gamma_*(\mathcal{S}_X)
\tag{8.14.5.3}
\end{equation}
already defined in (\hyperref[II.3.3.2]{3.3.2}), which therefore makes
$\Gamma_*(\mathcal{F})$ a graded $\mathcal{S}$-module; the homomorphisms
(\hyperref[II.8.14.5.1]{8.14.5.1}) then define a homomorphism (of
degree $0$) of graded $\mathcal{S}$-modules
\begin{equation}
\label{II.8.14.5.4}
  \alpha:\mathcal{M}\to\Gamma_*(\operatorname{Proj}(\mathcal{M})).
\tag{8.14.5.4}
\end{equation}
\end{env}

\begin{env}[8.14.6]
\label{II.8.14.6}
In general, for a quasi-coherent graded $\mathcal{S}_X$-module
$\mathcal{F}$, it is not certain that the graded $\mathcal{S}$-module
$\Gamma_*(\mathcal{F})$ is necessarily quasi-coherent. Consider an
open subset $X'$ of $X$ such that the restriction $q':X'\to Y$ of $q$ to $X'$
is a \emph{quasi-compact} morphism. Since $q'$ is moreover separated,
$q'_*(\mathcal{F}')$ is then a quasi-coherent $\mathcal{O}_Y$-module for
every quasi-coherent $\mathcal{O}_{X'}$-module $\mathcal{F}'$
(\hyperref[I.9.2.2]{I, 9.2.2, b)}). We set
\begin{equation}
\label{II.8.14.6.1}
  \mathcal{S}_{X'}=\mathcal{S}_X|X'
  =\bigoplus_{n\in\mathbf{Z}}\mathcal{O}_X(n)|X'
\tag{8.14.6.1}
\end{equation}
and, for every graded $\mathcal{S}_{X'}$-module $\mathcal{F}'$,
\begin{equation}
\label{II.8.14.6.2}
  \Gamma'_*(\mathcal{F}')=
  \bigoplus_{n\in\mathbf{Z}}q'_*(\mathcal{F}'_n).
\tag{8.14.6.2}
\end{equation}
The preceding remark then shows that if $\mathcal{F}'$ is a
quasi-coherent $\mathcal{S}_{X'}$-module, $\Gamma'_*(\mathcal{F}')$ is a
quasi-coherent graded $\mathcal{S}$-module
(\hyperref[I.9.6.1]{I, 9.6.1}).

Note moreover that the canonical injection $j:X'\to X$ is
\emph{quasi-compact}, since $q'=q\circ j$ is quasi-compact and $q$ is separated
(\hyperref[I.6.6.4]{I, 6.6.4, (v)}). Consequently
$\mathcal{F}=j_*(\mathcal{F}')$ is a quasi-coherent graded
$\mathcal{S}_X$-module for every quasi-coherent graded
$\mathcal{S}_{X'}$-module $\mathcal{F}'$, and it follows from the preceding
definitions that we have
\begin{equation}
\label{II.8.14.6.3}
  \Gamma'_*(\mathcal{F}')=\Gamma_*(\mathcal{F}).
\tag{8.14.6.3}
\end{equation}
With the same hypotheses on $X'$, for every quasi-coherent graded
$\mathcal{S}$-module $\mathcal{M}$, we set
\begin{equation}
\label{II.8.14.6.4}
  \operatorname{Proj}'(\mathcal{M})
  =\operatorname{Proj}(\mathcal{M})|X'
\tag{8.14.6.4}
\end{equation}
which is a quasi-coherent graded $\mathcal{S}_{X'}$-module. The canonical
homomorphism
\[
  \operatorname{Proj}(\mathcal{M})\to
  j_*(\operatorname{Proj}'(\mathcal{M}))
\]
(\hyperref[0.4.4.3]{0, 4.4.3}) therefore gives a canonical homomorphism
$\Gamma_*(\operatorname{Proj}(\mathcal{M}))\to
\Gamma'_*(\operatorname{Proj}'(\mathcal{M}))$ of graded
$\mathcal{S}$-modules, and consequently, by composing it with
(\hyperref[II.8.14.5.4]{8.14.5.4}), we obtain a canonical functorial
homomorphism (of degree $0$) of quasi-coherent graded $\mathcal{S}$-modules
\begin{equation}
\label{II.8.14.6.5}
  \alpha':\mathcal{M}\to
  \Gamma'_*(\operatorname{Proj}'(\mathcal{M})).
\tag{8.14.6.5}
\end{equation}
\end{env}

\begin{env}[8.14.7]
\label{II.8.14.7}
Keep on $X'$ the hypotheses of
(\hyperref[II.8.14.6]{8.14.6}), and let $\mathcal{F}'$ be a
quasi-coherent graded $\mathcal{S}_{X'}$-module, so that
$\operatorname{Proj}'(\Gamma'_*(\mathcal{F}'))$ is also a
quasi-coherent graded $\mathcal{S}_{X'}$-module. We are going to define a
canonical functorial homomorphism (of degree $0$) of
graded $\mathcal{S}_{X'}$-modules
\begin{equation}
\label{II.8.14.7.1}
  \beta':\operatorname{Proj}'(\Gamma'_*(\mathcal{F}'))\to\mathcal{F}'.
\tag{8.14.7.1}
\end{equation}

\oldpage[II]{200}
Suppose first that $Y=\operatorname{Spec}(A)$ is affine and
$\mathcal{S}=\widetilde{S}$, where $S$ is an $A$-algebra graded in positive
degrees; then $\Gamma'_*(\mathcal{F}')=\widetilde{M}$, where
$M=\bigoplus_{n\in\mathbf{Z}}\Gamma(X',\mathcal{F}'_n)$ is a graded
$S$-module. Let $f\in S_d$ be such that $D_+(f)\subset X'$; by definition
(\hyperref[II.2.6.2]{2.6.2}), $\alpha_d(f)$ restricted to $D_+(f)$ is the
section of $\mathcal{O}_X(d)$ over $D_+(f)$ corresponding to
the element $f/1$ of $(S(d))_{(f)}$, and is consequently invertible; the
same is therefore true of $\alpha_{nd}(f^n)$ for every $n>0$. We conclude
immediately that we define an $S_f$-homomorphism (of degree $0$) of graded
modules $\beta_f:M_f\to\Gamma(D_+(f),\mathcal{F}')$ by assigning to every
element $z/f^n\in M_f$ (where $z\in M$) the section
$(z|D_+(f))(\alpha_{nd}(f^n)|D_+(f))^{-1}$ of $\mathcal{F}'$ over
$D_+(f)$. Moreover, we have a commutative diagram corresponding to
(\hyperref[II.2.6.4.1]{2.6.4.1}), whence the definition of $\beta'$ in this
case. To pass to the general case, one must consider an $A$-algebra $A'$,
the graded $A'$-algebra $S'=S\otimes_A A'$, and use a commutative diagram
analogous to (\hyperref[II.2.8.13.2]{2.8.13.2}); we leave
the details to the reader.
\end{env}

\begin{proposition}[8.14.8]
\label{II.8.14.8}
If $X'$ is an open subset of $X=\operatorname{Proj}(\mathcal{S})$ such that
$q':X'\to Y$ is quasi-compact, the homomorphism $\beta'$ defined in
(\hyperref[II.8.14.7]{8.14.7}) is bijective.
\end{proposition}

\begin{proof}
We may evidently restrict to the case where $Y$ is affine, and everything
comes down to proving (with the notation of (\hyperref[II.8.14.7]{8.14.7}))
that the homomorphism $\beta_f:M_f\to\Gamma(D_+(f),\mathcal{F}')$ is an
isomorphism. Now, when we replace $f$ by one of its powers, we change neither
$D_+(f)$ nor $\beta_f$; since $X'$ is \emph{quasi-compact} by
hypothesis, we may therefore always suppose, by virtue of
(\hyperref[II.8.14.4]{8.14.4}), that the sheaf $\mathcal{O}_X(d)$ is
\emph{invertible}. Since $X'$ is a scheme (because $q'$ is separated), the
proposition is then nothing other than (\hyperref[I.9.3.1]{I, 9.3.1}).
\end{proof}

\begin{corollary}[8.14.9]
\label{II.8.14.9}
Under the hypotheses of (\hyperref[II.8.14.8]{8.14.8}), every
quasi-coherent graded $\mathcal{S}_{X'}$-module $\mathcal{F}'$ is isomorphic
to a graded $\mathcal{S}_{X'}$-module of the form
$\operatorname{Proj}'(\mathcal{M})$, where $\mathcal{M}$ is a
quasi-coherent graded $\mathcal{S}$-module. If moreover $\mathcal{F}'$ is of
finite type, and if we suppose that $Y$ is a quasi-compact scheme or a
prescheme whose underlying space is Noetherian, we may suppose
$\mathcal{M}$ of finite type.
\end{corollary}

\begin{proof}
The proof from (\hyperref[II.8.14.8]{8.14.8}) follows
exactly the course of that of (\hyperref[II.3.4.5]{3.4.5}) from
(\hyperref[II.3.4.4]{3.4.4}), and we leave the details to the reader.
\end{proof}

\begin{proposition}[8.14.10]
\label{II.8.14.10}
Under the hypotheses of (\hyperref[II.8.14.7]{8.14.7}), let
$\mathcal{M}$ be a quasi-coherent graded $\mathcal{S}$-module and
$\mathcal{F}'$ a quasi-coherent graded $\mathcal{S}_{X'}$-module; the
composite homomorphisms
\begin{equation}
\label{II.8.14.10.1}
  \operatorname{Proj}'(\mathcal{M})
  \xrightarrow{\operatorname{Proj}'(\alpha')}
  \operatorname{Proj}'(\Gamma'_*(\operatorname{Proj}'(\mathcal{M})))
  \xrightarrow{\beta'}\operatorname{Proj}'(\mathcal{M})
\tag{8.14.10.1}
\end{equation}
\begin{equation}
\label{II.8.14.10.2}
  \Gamma'_*(\mathcal{F}')
  \xrightarrow{\alpha'}
  \Gamma'_*(\operatorname{Proj}'(\Gamma'_*(\mathcal{F}')))
  \xrightarrow{\Gamma'_*(\beta')}\Gamma'_*(\mathcal{F}')
\tag{8.14.10.2}
\end{equation}
are the identity isomorphisms.
\end{proposition}

\begin{proof}
The question is local on $Y$, and the verification is done as in
(\hyperref[II.2.6.5]{2.6.5}); we leave the details to the reader.
\end{proof}

\begin{remark}[8.14.11]
\label{II.8.14.11}
In Chapter III (\hyperref[III.2.3.1]{III, 2.3.1}), we will see that
when $Y$ is \emph{locally Noetherian} and $\mathcal{S}$ is a
quasi-coherent graded $\mathcal{O}_Y$-algebra \emph{of finite type} (in which
case

\oldpage[II]{201}
we can take $X'=X$), the homomorphism $\alpha$
(\hyperref[II.8.14.5.4]{8.14.5.4}) is \textbf{(TN)}-\emph{bijective} for
every quasi-coherent graded $\mathcal{S}$-module $\mathcal{M}$ satisfying
condition \textbf{(TF)}.
\end{remark}

\begin{remark}[8.14.12]
\label{II.8.14.12}
The situation described in (\hyperref[II.8.14.4]{8.14.4}) is a particular
case of the following one. Let $X$ be a ringed space and $\mathcal{S}$ a
graded $\mathcal{O}_X$-algebra (in positive and negative degrees); suppose
that there exists an integer $d>0$ such that $\mathcal{S}_d$ and
$\mathcal{S}_{-d}$ are \emph{invertible}, the canonical homomorphism
\begin{equation}
\label{II.8.14.12.1}
  \mathcal{S}_d\otimes_{\mathcal{O}_X}\mathcal{S}_{-d}
  \to\mathcal{O}_X
\tag{8.14.12.1}
\end{equation}
being an \emph{isomorphism} (so that $\mathcal{S}_{-d}$ is identified with
$\mathcal{S}_d^{-1}$). We then say that the graded $\mathcal{O}_X$-algebra
$\mathcal{S}$ is \emph{periodic of period $d$}. This name
comes from the following property: \emph{under the preceding hypotheses,
for every graded $\mathcal{S}$-module $\mathcal{F}$, the canonical
homomorphism}
\begin{equation}
\label{II.8.14.12.2}
  \mathcal{S}_d\otimes\mathcal{F}_n\to\mathcal{F}_{n+d}
\tag{8.14.12.2}
\end{equation}
\emph{is an isomorphism for every $n\in\mathbf{Z}$.} Indeed, the question
is local on $X$, and we may therefore suppose that $\mathcal{S}_d$ has an
\emph{invertible} section $s$ over $X$, its inverse $s'$ being a
section of $\mathcal{S}_{-d}$. The homomorphism
$\mathcal{F}_{n+d}\to\mathcal{S}_d\otimes\mathcal{F}_n$ which assigns to every
section $z\in\Gamma(U,\mathcal{F}_{n+d})$ the section
$(s|U)\otimes(s'|U)z$ of $\mathcal{S}_d\otimes\mathcal{F}_n$ over $U$ is
then the inverse of (\hyperref[II.8.14.12.2]{8.14.12.2}), whence our
assertion. From this we deduce, for every $k\in\mathbf{Z}$, a canonical
isomorphism
\[
  (\mathcal{S}_d)^{\otimes k}\otimes\mathcal{F}_n
  \xrightarrow{\sim}\mathcal{F}_{n+kd}.
\]
Consequently, \emph{a graded $\mathcal{S}$-module $\mathcal{F}$ is known
when one knows the $\mathcal{S}_0$-modules $\mathcal{F}_i$
$(0\leqslant i\leqslant d-1)$ and the canonical homomorphisms}
\[
  \mathcal{S}_i\otimes\mathcal{F}_j\to\mathcal{F}_{i+j}
  \qquad\text{for }0\leqslant i,j\leqslant d-1
\]
(setting
$\mathcal{F}_{i+j}=\mathcal{S}_d\otimes_{\mathcal{S}_0}
\mathcal{F}_{i+j-d}$ when $i+j\geqslant d$). Of course, for these
homomorphisms to define an $\mathcal{S}$-module structure on the
direct sum of the $(\mathcal{S}_d)^{\otimes k}\otimes\mathcal{F}_i$
$(k\in\mathbf{Z},\ 0\leqslant i\leqslant d-1)$, they must satisfy
associativity conditions which we will not make explicit.

In the case $d=1$ (which is the one considered in
(\hyperref[subsection:II.3.3]{3.3})) we can therefore say that the category
of graded $\mathcal{S}$-modules (resp. of quasi-coherent ones, if $X$ is a
prescheme and $\mathcal{S}$ is quasi-coherent) is \emph{equivalent} to the
category of arbitrary (resp. quasi-coherent) $\mathcal{S}_0$-modules; it is in
this sense that the developments of this number can be regarded as
generalizing those of \S~\hyperref[section:II.3]{3}. Moreover, we see that,
under suitable finiteness conditions, the results of this number
(together with (\hyperref[II.8.14.11]{8.14.11})) reduce to a certain
extent the study of quasi-coherent graded algebras over a prescheme, and of
graded modules ``mod. \textbf{(TN)}'' over such algebras, to the study of the
particular case where the algebras under consideration are \emph{periodic}
(and where condition \textbf{(TN)} for $\mathcal{M}$
(\hyperref[II.3.4.2]{3.4.2}) therefore implies $\mathcal{M}=0$).
\end{remark}

\begin{remark}[8.14.13]
\label{II.8.14.13}
Under the hypotheses of (\hyperref[II.8.14.1]{8.14.1}), let $d$ be an
integer $>0$; we have defined a canonical $Y$-isomorphism $h$ of $X$ onto
$X^{(d)}=\operatorname{Proj}(\mathcal{S}^{(d)})$
(\hyperref[II.3.1.8]{3.1.8}). For every

\oldpage[II]{202}
quasi-coherent graded $\mathcal{S}$-module $\mathcal{M}$ and every integer $k$
such that $0\leqslant k\leqslant d-1$, we also have a canonical
$h$-isomorphism (with the notation of (\hyperref[II.3.1.1]{3.1.1}))
\begin{equation}
\label{II.8.14.13.1}
  (\operatorname{Proj}(\mathcal{M}))^{(d,k)}
  \xleftarrow{\sim}\operatorname{Proj}(\mathcal{M}^{(d,k)}).
\tag{8.14.13.1}
\end{equation}

Suppose first that $Y=\operatorname{Spec}(A)$ is affine,
$\mathcal{S}=\widetilde{S}$ and $\mathcal{M}=\widetilde{M}$, where $S$ is an
$A$-algebra graded in positive degrees and $M$ a graded $S$-module. We know
that for every $f\in S_e$ ($e>0$), $h$ maps $D_+(f)$ onto $D_+(f^d)$ and
corresponds to the canonical isomorphism
$S_{(f^d)}\xrightarrow{\sim}S_{(f)}$
(\hyperref[II.2.2.2]{2.2.2}). The restriction of
(\hyperref[II.8.14.13.1]{8.14.13.1}) to $D_+(f^d)$ then corresponds to the
canonical di-isomorphism $M_{f^d}\xrightarrow{\sim}M_f$, restricted to the
elements of $M_{f^d}$ of degree congruent to $k$ (mod. $d$). We leave it to
the reader to verify that these isomorphisms are compatible with passing
from $f$ to one of its homogeneous multiples $fg$, and then that there is an
analogous compatibility with passing from $S$ to a graded $A'$-algebra
$S'=S\otimes_A A'$, when $A'$ is an $A$-algebra.

In particular, we thus have an $h$-isomorphism
\begin{equation}
\label{II.8.14.13.2}
  (\mathcal{S}^{(d)})_{X^{(d)}}
  \xrightarrow{\sim}(\mathcal{S}_X)^{(d)}
\tag{8.14.13.2}
\end{equation}
which respects the multiplicative structures of the two sides, and by means
of which (\hyperref[II.8.14.13.1]{8.14.13.1}) becomes an
$h$-di-isomorphism of a graded $(\mathcal{S}^{(d)})_{X^{(d)}}$-module onto a
graded $(\mathcal{S}_X)^{(d)}$-module. Likewise, we have an $h$-isomorphism
\begin{equation}
\label{II.8.14.13.3}
  \operatorname{Proj}_0(\mathcal{S}^{(d,k)}(n))
  \xrightarrow{\sim}\mathcal{O}_X(nd+k)
\tag{8.14.13.3}
\end{equation}
which completes the result of (\hyperref[II.3.2.9]{3.2.9, (ii)}).

From the isomorphism (\hyperref[II.8.14.13.1]{8.14.13.1}) we immediately
deduce an isomorphism of graded $\mathcal{S}^{(d)}$-modules
\begin{equation}
\label{II.8.14.13.4}
  \Gamma_*^{(d)}(\operatorname{Proj}(\mathcal{M}^{(d,k)}))
  \xrightarrow{\sim}
  \Gamma_*((\operatorname{Proj}(\mathcal{M}))^{(d,k)})
\tag{8.14.13.4}
\end{equation}
where $\Gamma_*^{(d)}$ corresponds to the structural morphism
$q^{(d)}:X^{(d)}\to Y$; it is immediate to verify that the canonical
homomorphism $\alpha$ (\hyperref[II.8.14.5.4]{8.14.5.4}) and the
analogous homomorphism $\alpha^{(d)}$ for $X^{(d)}$ make the diagram
\begin{equation}
\label{II.8.14.13.5}
  \xymatrix{
    & \mathcal{M}^{(d,k)}\ar[dl]_{\alpha^{(d)}}\ar[dr]^{\alpha} & \\
    \Gamma_*^{(d)}(\operatorname{Proj}(\mathcal{M}^{(d,k)}))
      \ar[rr]^{\sim}
      && \Gamma_*((\operatorname{Proj}(\mathcal{M}))^{(d,k)})
  }
\tag{8.14.13.5}
\end{equation}
commutative, the verification being done by supposing $Y$ affine and
computing the restrictions of the images under $\alpha^{(d)}$ and $\alpha$ of
the same element of $M^{(d,k)}$ to the open subsets $D_+(f^d)$ and $D_+(f)$
(with the same notation as above).
\end{remark}

\begin{proposition}[8.14.14]
\label{II.8.14.14}
Let $Y$ be a quasi-compact prescheme, $\mathcal{S}$ a quasi-coherent graded
$\mathcal{O}_Y$-algebra of finite type, and $\mathcal{M}$ a quasi-coherent
graded $\mathcal{S}$-module satisfying condition
\textbf{(TF)}; let $X=\operatorname{Proj}(\mathcal{S})$. Then
$\mathcal{S}_X$ is a periodic graded $\mathcal{O}_X$-algebra
(\hyperref[II.8.14.12]{8.14.12}), and there exists a

\oldpage[II]{203}
period $d$ of $\mathcal{S}_X$ such that the
$(\operatorname{Proj}(\mathcal{M}))^{(d,k)}$
$(0\leqslant k\leqslant d-1)$ are
$(\mathcal{S}_X)^{(d)}$-modules of finite type.
\end{proposition}

\begin{proof}
Indeed, (\hyperref[II.3.1.10]{3.1.10}) proves that there exists $d$ such
that $\mathcal{S}^{(d)}$ is generated by
$\mathcal{S}_d=(\mathcal{S}^{(d)})_1$, the latter being an
$\mathcal{S}_0$-module of finite type. To prove the first assertion, we
may therefore, by virtue of (\hyperref[II.8.14.13.2]{8.14.13.2}), restrict
to the case $d=1$, and then the proposition follows from
(\hyperref[II.3.2.7]{3.2.7}). Moreover, taking into account
(\hyperref[II.8.14.13.1]{8.14.13.1}), the second assertion is a consequence
of (\hyperref[II.2.1.6]{2.1.6, (iii)}) and of
(\hyperref[II.3.4.3]{3.4.3}).
\end{proof}
